% concatenated from Hopf_Subgroup_MCG_ArXiv.tex
\documentclass[dvipdfmx]{amsart}




\makeatletter
\@addtoreset{equation}{section}
\def\theequation{\thesection.\arabic{equation}}
\makeatother


%\usepackage{refcheck}
\usepackage{amsfonts, amsmath, amssymb, amsthm}
%\usepackage{pdfsync}
\usepackage{url}
\usepackage{cleveref}
\usepackage[dvipdfmx]{graphicx}
%\usepackage[dvipdfmx]{color}
\usepackage[dvipdfmx]{color}
\usepackage{amscd}
\usepackage[right]{lineno}
%\modulolinenumbers[6]
%\usepackage{showkeys}
%\usepackage{eucal}
%\usepackage{xy}
\usepackage{mathrsfs}
\usepackage{MnSymbol}
\setcounter{tocdepth}{1}
\usepackage{enumerate}
\usepackage[all]{xy}



\newtheorem{theorem}{Theorem}[section]
\newtheorem{lemma}{Lemma}[section]
\newtheorem{proposition}{Proposition}[section]
\newtheorem{corollary}{Corollary}[section]
\newtheorem{definition}{Definition}[section]
\theoremstyle{remark}
\newtheorem{example}{Example}[section]
\newtheorem{remark}{Remark}[section]
\newtheorem{assumption}{Assumption}
\newtheorem{conjecture}{Conjecture}
\newtheorem*{convention}{Convention}
\newtheorem{claim}{Claim}
\newtheorem{problem}{Problem}

%%%%%%%%%%%%%%%%%%%%%%%%%%%%%%%%%%%%%%%%%%%%%%%%%
\newcommand{\ext}{{\rm Ext}}

\newcommand{\teich}{\mathcal{T}}
\newcommand{\torelli}{\mathscr{T}}

\newcommand{\mcg}{\operatorname{MCG}}
\newcommand{\tmod}{\operatorname{Mod}}
\newcommand{\torelliG}{\mathcal{I}}

\newcommand{\ThursM}{\mu_{Th}}
\newcommand{\PThursM}{\hat{\mu}_{Th}}
% \newcommand{\pushThursM}{{\boldsymbol \mu}}

\newcommand{\PSiegelThursM}{\mathscr{M}_{Th}}
\newcommand{\HSmeas}{{\bf m}}

\newcommand{\convgenus}{{\boldsymbol \xi}}

\newcommand{\ray}{{\boldsymbol r}}


\newcommand{\ml}{\mathcal{ML}}
\newcommand{\pml}{\mathcal{PML}}
\newcommand{\syp}{\mathop{Sp}}
\newcommand{\psyp}{\mathop{PSp}}
\newcommand{\gl}{\mathop{GL}}
\newcommand{\spl}{\mathop{Sp}}
% \newcommand{\gra}{\mathop{\mathbf{Gr}}}
\newcommand{\symm}{\mathop{\mathbf{Sym}}}
\renewcommand{\hom}{\mathop{\mathrm{Hom}}}


\newcommand{\rank}{\mathop{rank}}


\newcommand{\lllangle}{\left\llangle}
\newcommand{\rrrangle}{\right\rrangle}

\newcommand{\mat}{\mathrm{Mat}}

\newcommand{\symrep}{\mathop{\boldsymbol{\varrho}}}
\newcommand{\Cayley}{\Phi_{C\!a\!y}}


\newcommand{\abeljacobi}{\mathcal{A}}
\newcommand{\tree}{\boldsymbol{T}}
\newcommand{\protree}{\boldsymbol{p}}
\newcommand{\latticehomo}{\boldsymbol{L}}
\newcommand{\hypell}{\operatorname{\mathbf{h\ell}}}
\newcommand{\prG}{\operatorname{\varpi}}
\newcommand{\torpr}{\operatorname{\varpi}_{tor}}

\newcommand{\busemann}{\boldsymbol{b}}


\newcommand{\euclidvolume}{Evol}
\newcommand{\QThurs}{{\mathbf m}_{Th}}
\newcommand{\proje}{\boldsymbol{\psi}}

\newcommand{\distanceml}{\boldsymbol{\delta}}

\newcommand{\pushThursM}{{\boldsymbol \mu}}

\newcommand{\HM}{v}
\newcommand{\HMh}{h}

\newcommand{\horo}{\mathfrak{horo}}

\newcommand{\Dirichlet}{\mathscr{D}}
\newcommand{\horos}{\mathscr{H}}

\newcommand{\conservativepart}{\mathfrak{Cons}}
\newcommand{\dissipativepart}{\mathfrak{Diss}}

\begin{document}
%\frontmatter




\title[Hopf decomposition]{Hopf decomposition of the actions of subgroups of the mapping class group}
\author{Hideki Miyachi}


\date{\today}
\address{School of Mathematics and Physics,
College of Science and Engineering,
Kanazawa University,
Kakuma-machi, Kanazawa,
Ishikawa, 920-1192, Japan
}
\email{miyachi@se.kanazawa-u.ac.jp}
\thanks{This work is partially supported by JSPS KAKENHI Grant Numbers
25K00909,
26K21806,
23K22396 }
\subjclass[2020]{31C10, 32G05, 32G15, 32U35, 57M50}
\keywords{Teichm\"uller space, Thurston boundary,
Thurston measure, Hopf decomposition,
horospherical limit sets, Dirichlet polyhedra,
Torelli group, period map}


\begin{abstract}
We study the Hopf decomposition of subgroup actions of the
Teichm\"uller modular group on the Thurston boundary with respect to the Thurston measure class. Kaimanovich's general
Radon--Nikodym criteria allow us to describe the conservative and dissipative parts by the divergence and convergence of a series expressed in terms of extremal length. We identify the conservative part with the big horospherical limit set modulo null sets. For any basepoint with trivial stabilizer in the subgroup, we identify the dissipative part, modulo null sets, with the set of Dirichlet points and with the union of the subgroup translates of the ideal boundary of the associated
Dirichlet polyhedron. The description via Dirichlet polyhedra uses the fact that level sets of extremal-length ratios at distinct points of Teichm\"uller space have measure zero.

For the Torelli group of a closed surface of genus at least two, we use radial limits of the period map to prove that its conical limit set has measure zero. Combining our geometric characterization with the conservativity of its boundary action established by Choi, Gekhtman, Yang, and Zheng, we obtain that its big horospherical limit set has full measure and that the ideal boundary of each Dirichlet polyhedron has measure zero.
\end{abstract}




\maketitle
\tableofcontents


%%%%%%%%%%%%%%%%%%%%%%%%%%%%%%%%%%%%%%%%%%%%%
%%%%%%%%%%%%%%%%%%%%%%%%%%%%%%%%%%%%%%%%%%%%%
%%%%%%%%%%%%%%%%%%%%%%%%%%%%%%%%%%%%%%%%%%%%%
%%%%%%%%%%%%%%%%%%%%%%%%%%%%%%%%%%%%%%%%%%%%%


\section{Introduction}

The action of a subgroup of the mapping class group on the Thurston boundary can be studied from both geometric and measure-theoretic perspectives. In this paper, we relate these perspectives by describing the conservative and dissipative parts of the action in terms of extremal length and the geometry of subgroup orbits in Teichm\"uller space.
\Cref{table:dictionary} summarizes classical analogies between the geometry of the hyperbolic plane and that of Teichm\"uller space.
These analogies have served as a guiding principle in the author's
previous work and continue to inform the present study.
For example, this perspective underlies the author's alternative approach
to Krushkal's formula for the pluricomplex Green function
\cite{Miyachi2019}, as well as Poisson integral representations
\cite{Miyachi2023,miyachi2025bddpluriharmonicfunctionII-ToAppear}
and results on the radial limits of bounded pluriharmonic functions
\cite{Miyachi2024Bounded} on Teichm\"uller space.
\begin{table}
\centering
\begin{tabular}{|l|l|l|}
\hline
  & Hyperbolic plane $\mathbb{H}$ & $\teich_{g,m}$ \\
\hline\hline
Distance & Hyperbolic distance & Teichm\"uller--Kobayashi distance \\ \hline
Metric & Hyperbolic metric & Teichm\"uller--Kobayashi metric \\ \hline
Discrete groups & Fuchsian groups
& Subgroups of $\tmod(g,m)$ \\ \hline
Compactification & $\overline{\mathbb{H}}=\mathbb{H}\cup \hat{\mathbb{R}}$ ($\hat{\mathbb{R}}=\mathbb{R}\cup\{\infty\}$) & Thurston compactification
\\
\hline
Boundary & $\hat{\mathbb{R}}$ &
$\pml$
\\
\hline
Boundary measure & harmonic (visual) measure & normalized Thurston measure
\\
\hline
\end{tabular}
\caption{Analogies motivating the present research}
\label{table:dictionary}
\end{table}
% In the present paper, we ask how the Hopf decomposition of the boundary action is reflected in the geometry of subgroup orbits.

McCarthy and Papadopoulos studied limit sets and domains of discontinuity
for subgroup actions on the Thurston boundary
\cite{McCarthyPapadopoulos1989}. Their construction of fundamental
domains in Teichm\"uller space \cite{McCarthyPapadopoulos1996} provides
an analogue of Dirichlet polyhedra. Kent and Leininger
\cite{KentLeininger2008} studied conical limit points in connection with
convex cocompactness. Horospherical limit points in the Thurston boundary
were studied in \cite{miyachi2025functiontheorydynamicsergodic}
and in joint work with Ohshika \cite{miyachi2026limitsetsmappingclass}.
In that joint work, we give a complete classification of limit points in terms of the geometry of measured laminations.

Sullivan's and Tukia's descriptions of boundary actions of
Kleinian groups relate the Hopf decomposition to the geometry of horospherical limit sets \cite{Sullivan1981,Tukia1997}. These results motivate the question of whether an analogous geometric description holds for subgroup actions on the Thurston boundary.
Kaimanovich \cite{Kaimanovich2010} gave general criteria for the Hopf decomposition of nonsingular actions in terms of Radon--Nikodym derivatives.
Choi, Gekhtman, Yang, and Zheng \cite{choi2024confinedsubgroupsgroupscontracting} developed related results for groups with contracting elements; their Lemma~2.48 identifies the infinitely conservative part with a big horospherical limit set in their boundary framework.

A related measure-classification problem was studied by Choi and Kim \cite{ChoiKim2025InvariantMeasures}, who classified invariant Radon measures concentrated on the recurrence locus in the space of measured laminations for non-elementary subgroups of the mapping class group. Their results concern invariant measures on the non-projectivized space, while the present paper studies the Hopf decomposition of the action on the projectivized space with respect to the Thurston measure class.

\subsection{Main results}
We give two descriptions of the Hopf decomposition: one in terms of a series of Radon--Nikodym derivatives, and the other in terms of horospherical limit sets and Dirichlet points.

Let $\teich_{g,m}$ denote the Teichm\"uller space of Riemann surfaces of analytically finite type $(g,m)$, where $2g-2+m>0$ and $\convgenus=3g-3+m\ge 1$, and let $\pml=\pml_{g,m}$ denote the space of projective measured laminations on a fixed complete hyperbolic surface of genus $g$ with $m$ cusps. We identify $\pml$ with the Thurston boundary of $\teich_{g,m}$.
Let $\tmod(g,m)$ denote the Teichm\"uller modular group, namely the image of the mapping class group in the group of automorphisms of $\teich_{g,m}$.

\subsubsection{Hopf decomposition}
For $x\in\teich_{g,m}$, let $\ext_x(\lambda)$ denote the extremal length of a measured lamination $\lambda$ on $x$ (cf. \S\ref{subsec:HMdiffExtremalLength}).
The extremal-length unit ball in the space of measured laminations determines a probability measure $\PThursM^x$ on $\pml$ by projectivizing and normalizing the Thurston measure (see \eqref{eq:def_Th_measure}).
The measures $\PThursM^x$ are mutually absolutely continuous and define the Thurston measure class on $\pml$.
The action of $\tmod(g,m)$ is nonsingular with respect to this measure
class, and its Radon--Nikodym derivatives are given by
\[
\frac{d([\omega]_*\PThursM^x)}{d\PThursM^x}([\lambda])
=\left(
\frac{\ext_x(\lambda)}{\ext_{[\omega](x)}(\lambda)}
\right)^{\convgenus}
\qquad ([\omega]\in\tmod(g,m));
\]
see \eqref{eq:Radon-Nikodym2}.

This formula leads us to consider the sum of these derivatives over
a subgroup. For a nontrivial subgroup $G\le\tmod(g,m)$, define
\begin{align*}
\conservativepart(G)
&=\conservativepart_x(G)
=\left\{[\lambda]\in\pml\mathrel{}\middle|\mathrel{}
\sum_{[\omega]\in G}
\left(
\frac{\ext_x(\lambda)}
{\ext_{[\omega](x)}(\lambda)}
\right)^{\convgenus}
=\infty
\right\},\\
\dissipativepart(G)
&=\dissipativepart_x(G)
=\left\{[\lambda]\in\pml\mathrel{}\middle|\mathrel{}
\sum_{[\omega]\in G}
\left(
\frac{\ext_x(\lambda)}
{\ext_{[\omega](x)}(\lambda)}
\right)^{\convgenus}
<\infty
\right\}.
\end{align*}
These sets are independent of the choice of $x$
(see \Cref{lem:Cons-Diss-ElementaryProperty}).
Our first main result applies the Radon--Nikodym criterion to identify
these sets with the conservative and dissipative parts of the action.

\begin{theorem}[Hopf decomposition]
\label{thm:Hopf_decomposition_main}
Let $G\le\tmod(g,m)$ be a nontrivial subgroup.
For every $x\in\teich_{g,m}$, the decomposition
\[
\pml=\conservativepart(G)\sqcup\dissipativepart(G)
\]
agrees, modulo $\PThursM^x$-null sets, with the Hopf decomposition
of the nonsingular action of $G$ on $(\pml,\PThursM^x)$.
More precisely,
\begin{enumerate}
\item the action of $G$ on $\conservativepart(G)$ is conservative;
\item the action of $G$ on $\dissipativepart(G)$ is dissipative.
\end{enumerate}
\end{theorem}

Our second main result gives a geometric characterization of these
two parts, in analogy with classical results for Kleinian groups.
Sullivan \cite{Sullivan1981} identified the conservative part of the boundary action, with respect to spherical Lebesgue measure, with the horospherical limit set modulo null sets.
Tukia \cite[Theorem~1]{Tukia1997} extended this description to non-atomic conformal measures by using the big horospherical limit set.
% The following theorem gives the corresponding description of the conservative part on the Thurston boundary and identifies the dissipative part with the set of Dirichlet points, modulo Thurston null sets.
In the setting of the Thurston boundary, we obtain a parallel
description of the Hopf decomposition. The conservative part
agrees with the big horospherical limit set. For a basepoint
with trivial stabilizer, the dissipative part agrees with
the set of Dirichlet points and with the union of the
subgroup translates of the ideal boundary of the associated
Dirichlet polyhedron. All these identifications are
understood modulo Thurston null sets.

The big horospherical limit set $\Lambda_H(G)$ consists
of the boundary points for which some horoball centered at that point
contains infinitely many points of a $G$-orbit.
For $x_0\in\teich_{g,m}$, let $\operatorname{Dir}_G(x_0)$ denote
the set of Dirichlet points: these are the boundary points whose
associated Busemann function attains its infimum on $G(x_0)$.
For uniquely ergodic $[\lambda]$, this is equivalent to the
attainment of the minimum of $\ext_{[\omega](x_0)}(\lambda)$ over
$[\omega]\in G$; see \eqref{eq:horo_definition4}.

\begin{theorem}[Geometric characterization of the Hopf decomposition]
\label{thm:main_classfication_limit_sets}
Let $G\le\tmod(g,m)$ be a nontrivial subgroup.
Then $\conservativepart(G)$ agrees with $\Lambda_H(G)$ modulo
Thurston null sets.
If $x_0\in\teich_{g,m}$ satisfies
$\operatorname{Stab}_G(x_0)=\{\mathrm{id}\}$, then
$\dissipativepart(G)$ agrees with $\operatorname{Dir}_G(x_0)$
modulo Thurston null sets.
\end{theorem}

Here and below, null sets on $\pml$ are understood with respect to the Thurston measure class, or equivalently with respect to any of the measures $\PThursM^x$.
Under the same assumption on $x_0$, the dissipative part also agrees, modulo null sets, with the union of the $G$-translates of the ideal boundary of the Dirichlet polyhedron (in the sense of McCarthy--Papadopoulos \cite{McCarthyPapadopoulos1996}) centered at $x_0$ (see \Cref{thm:Dissipative_part}).
Both main theorems are proved in \S\ref{sec:Hopf_decomposition}.

The identification of the conservative part uses Kaimanovich's
Radon--Nikodym criterion \cite[Theorem~29]{Kaimanovich2010}, together
with the essential freeness of the boundary action and the expression of the relevant Busemann functions in terms of extremal length on the uniquely ergodic locus.
For the description of the dissipative part, we use the nullity of level loci of extremal-length ratios at distinct points of Teichm\"uller space (\Cref{thm:level_set}). 
In particular, for almost every $[\lambda]\in\pml$
and any representative $\lambda\in\ml\setminus\{0\}$,
the function $y\mapsto\ext_y(\lambda)$ takes distinct
values at distinct points of a fixed subgroup orbit.
Consequently, whenever its minimum over the orbit is attained,
the minimizing orbit point is unique.
This conclusion is independent of the choice of representative.
% In particular, outside a Thurston null set, no two distinct points of a fixed subgroup orbit have equal extremal length for a representative of the same projective measured lamination. Thus, whenever the minimum over the orbit is attained, the minimizing orbit point is unique.
This observation is used to show that the ideal boundary of a Dirichlet polyhedron is wandering and to identify its translates with the dissipative part modulo null sets. 
The nullity theorem follows from local real analyticity
in common train-track coordinates and the fact that
extremal length functions at distinct points cannot
be proportional on a nonempty open set.

In \cite[Proposition~8.2]{miyachi2025functiontheorydynamicsergodic},
the author showed that the big horospherical limit set contains
no wandering set of positive measure.
The present results identify both parts of the Hopf decomposition
and describe the dissipative part using Dirichlet polyhedra.
In particular, \Cref{coro:conservativity_Dirichletset} establishes the equivalence between conservativity, full measure of
the big horospherical limit set, and nullity of the ideal boundary
of a Dirichlet polyhedron. This establishes the equivalence asked for in \cite[Problem~9]{miyachi2025functiontheorydynamicsergodic},
with the boundary fundamental domain interpreted as the ideal boundary of a Dirichlet polyhedron centered at a basepoint with trivial stabilizer.

\subsubsection{Dynamics of the Torelli group}
As an application, in \S\ref{sec:dynamics_of_Torelli_group} we consider the Torelli group $\torelliG_g$ of a closed surface of genus $g\ge 2$.

\begin{corollary}[Dynamics of the Torelli group]
\label{coro:torelli_main}
For $g\ge 2$, the action of $\torelliG_g$ on $\pml_{g,0}$ is
conservative but not ergodic with respect to the Thurston measure class, and its big horospherical limit set has full measure.
Moreover, we have
\begin{itemize}
\item[(a)] the conical limit set of the Torelli group has measure zero;
\item[(b)] the ideal boundary of a Dirichlet polyhedron of the Torelli group has measure zero for every choice of basepoint.
\end{itemize}
\end{corollary}

Since the Torelli group is torsion-free, the stabilizer condition in the geometric characterization is satisfied at every basepoint.
\Cref{table:comparison} gives a dynamical comparison among $\operatorname{PSL}_2(\mathbb{Z})$, the Teichm\"uller modular group and the Torelli group.

\begin{table}
\centering
\begin{tabular}{|c|c|c|c|}
\hline
& $\operatorname{PSL}_2(\mathbb{Z})$
& $\tmod(g)$
& $\mathcal{I}_g$ \\ \hline\hline
Limit set
& $\partial \mathbb{H}$
& \multicolumn{2}{c|}{$\pml$} \\
\hline
Ergodicity
& \multicolumn{2}{c|}{Yes}
& No \\ \hline
Conservativity
& \multicolumn{3}{c|}{Yes}
 \\ \hline
Measure of $\Lambda_C$
& \multicolumn{2}{c|}{$1$}
& $0$ \\ \hline
Measure of $\Lambda_H$
& \multicolumn{3}{c|}{$1$}\\ \hline
\end{tabular}
\caption{A dynamical comparison for $g\ge2$. For $\operatorname{PSL}_2(\mathbb Z)$,
measure refers to normalized visual measure on $\partial\mathbb H$;
for the other two groups, it refers to $\PThursM^{x_0}$ on $\pml_{g,0}$.
Ergodicity and conservativity are understood with respect to these
measure classes.}
\label{table:comparison}
\end{table}
The non-ergodicity of this action was established in
\cite{Miyachi2024Bounded}.
Since the Torelli group is an infinite normal subgroup of the mapping
class group, it is confined. Its conservativity therefore follows from
\cite[Theorem~1.6(1)]{choi2024confinedsubgroupsgroupscontracting}.
The full-measure assertion for the big horospherical limit set and
the nullity of the ideal boundary of a Dirichlet polyhedron then
follow from \Cref{coro:conservativity_Dirichletset}.
The measure-zero statement for the conical limit set is proved
using radial limits of the period map
(\Cref{cor:conical_limit_sets}).
The proof relates conical limit points to the limiting
behavior of the period map along Teichm\"uller geodesic rays.
Thus, almost every boundary point is a big horospherical limit point
of the Torelli group, whereas its conical limit set has measure zero.
This distinguishes the horospherical behavior that describes
conservativity of the boundary action from conical recurrence.

\subsection*{Acknowledgements}
The author is deeply grateful to Professor Athanase Papadopoulos and Professor Ken'ichi Ohshika for their continuous encouragement and generous support.

%%%%%%%%%%%%%%%%%%%%%%%%%%%%%%%%%%%%%%%%%%%%%
%%%%%%%%%%%%%%%%%%%%%%%%%%%%%%%%%%%%%%%%%%%%%
%%%%%%%%%%%%%%%%%%%%%%%%%%%%%%%%%%%%%%%%%%%%%
%%%%%%%%%%%%%%%%%%%%%%%%%%%%%%%%%%%%%%%%%%%%%
% The action of a subgroup of the mapping class group on the Thurston boundary can be studied from both geometric and measure-theoretic perspectives. In this paper, we relate these perspectives by describing the conservative and dissipative parts of the action in terms of extremal length and the geometry of subgroup orbits in Teichm\"uller space.

% Let $\teich_{g,m}$ be the Teichm\"uller space of Riemann surfaces of analytically finite type $(g,m)$ with $\convgenus=3g-3+m\ge 1$,
% and $\pml=\pml_{g,m}$ be the space of projective measured laminations on a (fixed) complete hyperbolic surface of genus $g$ with $m$-cusps.
% We think of $\pml$ as the Thurston boundary of $\teich_{g,m}$.
% Let $\tmod(g,m)$ be the Teichm\"uller modular group acting on the Thurston compactification.

% For $x\in \teich_{g,m}$, we denote by $\ext_x(\lambda)$ be the extremal length of a measured lamination $\lambda$ on $x$.
% For a nontrivial subgroup $G$ of $\tmod(g,m)$ and $x\in \teich_{g,m}$, we define
% \begin{align*}
% \conservativepart(G)
% &=\conservativepart_x(G)
% =\left\{[\lambda]\in\pml\mathrel{}\middle|\mathrel{}
% \sum_{[\omega]\in G}
% \left(
% \frac{\ext_x(\lambda)}
% {\ext_{[\omega](x)}(\lambda)}
% \right)^{\convgenus}
% =\infty
% \right\},\\
% \dissipativepart(G)
% &=\dissipativepart_x(G)
% =\left\{[\lambda]\in\pml\mathrel{}\middle|\mathrel{}
% \sum_{[\omega]\in G}
% \left(
% \frac{\ext_x(\lambda)}
% {\ext_{[\omega](x)}(\lambda)}
% \right)^{\convgenus}
% <\infty
% \right\}.
% \end{align*}
% The main results of this paper is to establish the following Hopf decomposition for the actions of subgroups of the Teichm\"uller modular group. This is an analogue of Sullivan's dichotomy of the mesurable action of the Kleinian groups \cite{Sullivan1981}.

% \begin{theorem}[Hopf decomposition]
% \label{thm:Hopf_decomposition_main}
% Let $G\le\tmod(g,m)$ be a nontrivial subgroup. Then, modulo
% null sets of the Thurston measure, the decomposition
% \[
% \pml
% =
% \conservativepart(G)\sqcup\dissipativepart(G)
% \]
% is the Hopf decomposition of the action of $G$ on
% $(\pml,\PThursM^x)$. More precisely,
% \begin{enumerate}
% \item the action of $G$ on $\conservativepart(G)$ is conservative;
% \item the action of $G$ on $\dissipativepart(G)$ is dissipative.
% \end{enumerate}
% \end{theorem}

% \begin{theorem}[Classification of limit sets]
% \label{thm:main_classfication_limit_sets}
% Modulo null sets of the Thurston measure, the conservative part and dissipative part agree with the big horospherical limit sets and the Dirichlet points, respectively.
% % Moreover, the action on $\conservativepart(G)$ is infinitely
% % conservative in the sense of Kaimanovich \cite{Kaimanovich2010}.
% \end{theorem}

% \Cref{thm:Hopf_decomposition_main} is proved in \S\ref{sec:Hopf_decomposition}.

% \begin{corollary}[Dynamics of Torelli group]
% \label{coro:torelli_main}
% For $g\ge 2$,
% the action of the Torelli group $\torelliG_g$ for closed surfaces of genus $g$ is conservative but not ergodic. Furthemore, the big horospherical limit set is full, but the conical limit set of the Torelli group is null in terms of the Thurston measure.
% \end{corollary}

% The non-ergodicity of the action of $\torelliG_g$ is obtained in \cite[Corollary 1.4]{Miyachi2023}.
% The conservativity of the action follows from a recent result by Choi, Gekhtman, Yang, and Zheng \cite{choi2024confinedsubgroupsgroupscontracting}.
% The dynamics of the Torelli group is discussed in \S\ref{sec:dynamics_of_Torelli_group}.








\section{Notation and preliminaries}
Henceforth, we fix non-negative integers $g$ and $m$ such that $2g-2+m>0$ and $\convgenus=3g-3+m>0$.
Let $\Sigma_{g,m}$ be an oriented surface of genus $g$ with $m$ punctures.
When $m=0$, we write $\Sigma_g$ for $\Sigma_{g,0}$.


\subsection{Teichm\"uller space}
A \emph{marked Riemann surface} $(M,f)$ of type $(g,m)$ is a pair consisting of an analytically finite Riemann surface $M$ of type $(g,m)$ and an orientation-preserving homeomorphism $f\colon \Sigma_{g,m}\to M$.
Two marked Riemann surfaces $(M_1,f_1)$ and $(M_2,f_2)$ are said to be \emph{Teichm\"uller equivalent} if there exists a biholomorphism $h\colon M_1\to M_2$ such that $h\circ f_1$ is homotopic to $f_2$ rel punctures.
By the uniformization theorem, any $x=(M,f)$ can be regarded as a marked complete hyperbolic surface of finite area by equipping $M$ with the complete hyperbolic metric in its conformal class.

\subsubsection{Topology of $\teich_{g,m}$}
The \emph{Teichm\"uller space} $\teich_{g,m}$ of Riemann surfaces of analytically finite type $(g,m)$ is the set of Teichm\"uller equivalence classes of marked Riemann surfaces of type $(g,m)$.
The Teichm\"uller space admits a canonical distance, called the \emph{Teichm\"uller distance},
defined by
\[
d_T(x_1,x_2)=\frac{1}{2}\log \inf_h K(h)
\]
for $x_i=(M_i,f_i)\in \teich_{g,m}$ ($i=1,2$), where $h$ runs over all quasiconformal mappings $h\colon M_1\to M_2$ such that $h\circ f_1$ is homotopic to $f_2$ rel punctures and $K(h)$ is the maximal dilatation of $h$.
The metric space $(\teich_{g,m},d_T)$ is known to be complete and homeomorphic to Euclidean space of dimension $2\convgenus$.

\subsection{Measured laminations}
\label{subsec:measured_laminations}
Let $\mathcal{S}$ be the set of homotopy classes of 
essential simple closed curves on $\Sigma_{g,m}$.
Let $i(\alpha,\beta)$ denote the \emph{(geometric) intersection number} of simple closed curves $\alpha,\beta\in \mathcal{S}$.
For $t,s\ge0$ and $\alpha,\beta\in\mathcal S$, set
\begin{equation}
\label{eq:intersection-number-WS}
i(t\alpha,s\beta)=ts\,i(\alpha,\beta).
\end{equation}

\subsubsection{Measured laminations}
Henceforth, we fix a complete hyperbolic structure of finite area on $\Sigma_{g,m}$.
A \emph{geodesic lamination} $L$ on $\Sigma_{g,m}$ is a nonempty closed subset that is a disjoint union of complete simple geodesics, where a geodesic is said to be \emph{complete} if it is either closed or has infinite length in both directions.
The geodesics in $L$ are called the \emph{leaves} of $L$.
A \emph{transverse measure} on a geodesic lamination $L$ is an assignment of a Borel measure to each arc transverse to $L$, subject to the following two conditions: if a transverse arc $k'$ is contained in a transverse arc $k$, the measure assigned to $k'$ is the restriction of the measure assigned to $k$; and if two transverse arcs $k$ and $k'$ are homotopic through transverse arcs, the homotopy carries the measure assigned to $k$ to that assigned to $k'$.
A transverse measure on a geodesic lamination $L$ is said to have \emph{full support} if the support of the measure assigned to each transverse arc $k$ is exactly $k\cap L$.

A \emph{measured lamination} $\lambda$ is a pair consisting of a compact geodesic lamination $|\lambda|$, called the \emph{support} of $\lambda$, and a transverse measure of full support on $|\lambda|$.
For simplicity, we use the same symbol $\lambda$ for the measured lamination and its transverse measure.

Let $\ml$ be the set of measured laminations on $\Sigma_{g,m}$, together with the zero measured lamination $0$.
A weighted simple closed curve $t\alpha$ is identified with
a measured lamination supported on the simple closed geodesic homotopic to $\alpha$, with transverse measure equal to $t$ times the counting measure on its intersections with transverse arcs. A measured lamination $\lambda$ is said to be \emph{essentially complete} if its support is maximal among geodesic laminations. Equivalently, $\lambda$ is essentially complete if and only if every component of $\Sigma_{g,m}\setminus |\lambda|$ is either an ideal triangle or a punctured monogon (cf. \cite[\S9.5]{Thurston-LectureNote}).

\subsubsection{Topology on $\ml$}
Birman and Series \cite{BirmanSeries1985} showed that the union of all complete simple geodesics on $\Sigma_{g,m}$ has Hausdorff dimension $1$.
Consequently, one may choose geodesic arcs transverse to every complete simple geodesic; we call such arcs \emph{generic}.
The topology on $\ml$ is defined as follows:
a sequence $\{\lambda_n\}_n\subset \ml$ converges to $\lambda\in \ml$ if, for every generic transverse arc $k$ and every continuous function $\varphi$ on $k$,
\[
\int_k \varphi\,d\lambda_n\to \int_k\varphi\,d\lambda
\]
as $n\to \infty$ (cf. \cite[Part II]{Bonahon2001}).
The intersection number \eqref{eq:intersection-number-WS}
on weighted simple closed curves extends continuously to $\ml\times \ml$.
Indeed, for $\lambda\in \ml$ and $\alpha\in \mathcal{S}$,
\[
i(\lambda,\alpha)=\inf_{\alpha'\in \alpha}\int_{\alpha'}d\lambda
\]
(e.g. \cite{Bonahon2001} and \cite{Thurston-LectureNote}).

The multiplicative group $\mathbb{R}_{>0}$ acts on $\ml$ by multiplication. The quotient space $\pml=(\ml-\{0\})/\mathbb{R}_{>0}$ is called the \emph{space of projective measured laminations}.
Thurston showed that $\ml$ and $\pml$ are homeomorphic to $\mathbb{R}^{2\convgenus}$ and $\mathbb{S}^{2\convgenus-1}$, respectively.

The space $\ml$ of measured laminations is embedded in the space of non-negative functions on $\mathcal{S}$ by
\begin{equation}
\label{eq:ml-emb}
\ml\ni \lambda\mapsto (\mathcal{S}\ni \alpha\mapsto i(\lambda,\alpha))\in \mathbb{R}_{\ge 0}^{\mathcal{S}}.
\end{equation}

\subsubsection{Train track coordinates}
We follow Penner and Harer \cite{PennerHarer1992}
for terminology and conventions concerning train tracks.
We recall the notions used below.

A train track on $\Sigma_{g,m}$ is an embedded $1$-complex $\tau$ satisfying the following properties. Each edge (called a \emph{branch}) is a smooth path with well-defined tangent vectors at the endpoints, and at any vertex (called a \emph{switch}) the incident edges are mutually tangent. The tangent vector at the switch pointing toward the interior of an edge can have two possible directions, and this divides the ends of edges at the switch into two sets, neither of which is permitted to be empty. We call them ``incoming'' and ``outgoing.'' The valence of each
switch is at least 3, except possibly for one bivalent switch in a closed
curve component. Finally, we require that the components $R$ of $\Sigma_{g,m}\setminus \tau$ have negative generalized Euler characteristic, defined as the usual Euler characteristic $\chi(R)$ minus $1/2$ for every outward-pointing cusp, plus $1/2$ for every inward-pointing cusp (cf. \cite[\S1.1]{PennerHarer1992}).

For a train track $\tau$, a \emph{weight} on $\tau$ is a function $\mu$ on the set of branches of $\tau$ such that, at every switch $v$ of $\tau$, it satisfies the following \emph{switch condition}
\[
\mu(b_1)+\cdots+\mu(b_r)=\mu(b_{r+1})+\cdots+\mu(b_{r+t}),
\]
where $b_1,\ldots,b_r$ are the incoming branch ends incident at $v$ and $b_{r+1},\ldots,b_{r+t}$ are the outgoing branch ends.
We denote by $W(\tau)$ the vector space of weights on $\tau$ and by $E(\tau)\subset W(\tau)$ the cone of non-negative weights on $\tau$.
An element of $E(\tau)$ is also called a \emph{transverse measure} on $\tau$. A pair $(\tau,\mu)$ consisting of a train track $\tau$ and a transverse measure $\mu$ on $\tau$ is called a \emph{measured train track}.
A measured train track $(\tau,\mu)$ is said to be \emph{positive} if $\mu(b)>0$ for every branch $b$ of $\tau$.

A train track $\tau$ is called \emph{recurrent} if it admits a positive transverse measure (cf. \cite[\S1.3]{PennerHarer1992}). Let $b(\tau)$ and $s(\tau)$ denote the numbers of branches and switches of $\tau$, respectively, and let $o(\tau)$ denote the number of orientable connected components of $\tau$. Then $\dim W(\tau)=b(\tau)-s(\tau)+o(\tau)$ (cf. \cite[Lemma 2.1.1]{PennerHarer1992}). In particular, if $\tau$ is connected, then $\dim W(\tau)=b(\tau)-s(\tau)$ when $\tau$ is non-orientable, whereas $\dim W(\tau)=b(\tau)-s(\tau)+1$ when $\tau$ is orientable.

We say that a measured lamination $\lambda$ on $\Sigma_{g,m}$ is \emph{carried} by a train track $\tau$ if there is a $C^1$ map $\phi\colon \Sigma_{g,m}\to \Sigma_{g,m}$, called a \emph{supporting map}, such that $\phi(|\lambda|)\subset \tau$, $\phi$ is homotopic to the identity, and the restriction of the differential $d\phi$ to the tangent line to leaves of $|\lambda|$ at $p$ is nonzero for every $p\in |\lambda|$. 
For a train track $\tau$, we denote by $\ml(\tau)$ the set of measured laminations on $\Sigma_{g,m}$ that are carried by $\tau$.

A measured lamination $\lambda$ carried by $\tau$ determines a transverse measure $\mu=\mu_\lambda\in E(\tau)$ independently of the choice of supporting map (cf. \cite[Proposition 1.7.5]{PennerHarer1992}).
Conversely, every transverse measure $\mu\in E(\tau)$ determines a measured lamination $\lambda_\mu$ on $\Sigma_{g,m}$, with the zero weight corresponding to the zero measured lamination (cf. \cite[Construction 1.7.7 and Corollary 2.7.3]{PennerHarer1992}).

By \cite[Theorem 2.7.4]{PennerHarer1992}, when $\tau$ is recurrent,
the correspondence above gives a natural continuous bijection, which we
call the \emph{realizing map}:
\begin{equation}
\label{eq:realizing_map}
R_\tau\colon E(\tau)\ni \mu\mapsto \lambda_\mu\in \ml(\tau)\subset \ml.
\end{equation}
The collection $\{(E(\tau),R_\tau)\mid \text{$\tau$ is recurrent}\}$ gives integral piecewise-linear coordinates on $\ml$ (cf. \cite[Proposition 9.5.8]{Thurston-LectureNote} and \cite[Theorem 3.1.4]{PennerHarer1992}).


\subsection{Hubbard--Masur differentials and extremal length}
\label{subsec:HMdiffExtremalLength}
Let $x=(M,f)\in\teich_{g,m}$.
We denote by $\mathcal{Q}_x$ the finite-dimensional complex
vector space of integrable holomorphic quadratic
differentials on $M$, equipped with the $L^1$-norm
\[
\|q\|=\int_M |q(z)|\,\frac{|dz\wedge d\overline{z}|}{2},
\]
where $q=q(z)\,dz^2$ is the local expression for
$q\in\mathcal{Q}_x$.
Equivalently, after filling in the punctures of $M$, elements of $\mathcal{Q}_x$ are meromorphic quadratic differentials with at most simple poles at the punctures.
For $q=q(z)\,dz^2\in\mathcal{Q}_x$, we define a measured lamination $v_x(q)$, called the \emph{vertical lamination} (or vertical foliation), by the condition
\[
i\bigl(v_x(q),\alpha\bigr)
=
\inf_{\alpha'\in\alpha}
\int_{f(\alpha')}
\left|\operatorname{Re}\bigl(\sqrt{q(z)}\,dz\bigr)\right|
\]
for every $\alpha\in\mathcal{S}$.

Hubbard and Masur \cite{HubbardMasur1979} showed that the map
\begin{equation}
\label{eq:HM-map}
v_x\colon \mathcal{Q}_x\longrightarrow \ml
\end{equation}
is a homeomorphism. Accordingly, throughout this paper we refer to the map \eqref{eq:HM-map} as the \emph{Hubbard--Masur map}, although the map itself was known before the result of Hubbard and Masur. For each $\lambda\in\ml$, the \emph{Hubbard--Masur differential} $q_{\lambda,x}\in\mathcal{Q}_x$ is defined by $v_x(q_{\lambda,x})=\lambda$.

The \emph{extremal length} of $\lambda\in \ml$ on $x\in \teich_{g,m}$ is defined by
\[
\ext_x(\lambda)=\|q_{\lambda,x}\|.
\]
It satisfies $\ext_x(t\lambda)=t^2\ext_x(\lambda)$ for every $t\ge 0$.
The extremal length function
\[
\teich_{g,m}\times \ml\ni (x,\lambda)\mapsto \ext_x(\lambda)
\]
is continuous and satisfies the \emph{quasiconformal invariance}
\begin{equation}
\label{eq:QC-inv_extremal_length}
e^{-2d_T(x,y)}\le \frac{\ext_x(\lambda)}{\ext_y(\lambda)}\le e^{2d_T(x,y)}
\end{equation}
for all $x,y\in \teich_{g,m}$ and all $\lambda\in\ml\setminus\{0\}$.

\subsection{Teichm\"uller rays}
Let $x=(M,f)\in\teich_{g,m}$ and $[\lambda]\in\pml$. We define the \emph{Teichm\"uller (geodesic) ray} $\ray^\lambda_x\colon[0,\infty)\to\teich_{g,m}$ of direction $[\lambda]$ emanating from $x$ by
\[
\ray^{\lambda}_x(t)=(M_t,f_t\circ f),
\]
where $\ray^{\lambda}_x(0)=x$ and $f_t\colon M\to M_t$ is the quasiconformal mapping with Beltrami coefficient $\tanh(t)\overline{q_{\lambda,x}}/|q_{\lambda,x}|$ for $t\ge 0$, interpreted as $0$ at the zeros of $q_{\lambda,x}$. This definition depends only on the projective class $[\lambda]$.
It is known that $\ray^\lambda_x$ is a geodesic ray with respect to the Teichm\"uller distance:
\[
d_T(\ray^\lambda_x(t),\ray^\lambda_x(s))=|t-s|\quad (t,s\ge 0).
\]
Furthermore, for $x\in \teich_{g,m}$, the map
\[
\pml\times (0,\infty)\ni ([\lambda],t)\mapsto \ray^\lambda_x(t)\in \teich_{g,m}\setminus \{x\}
\]
is a homeomorphism.

\subsection{Thurston compactification}
For $x=(M,f)\in \teich_{g,m}$ and $\alpha\in \mathcal{S}$, we denote by $\hypell_x(\alpha)$ the length of the closed hyperbolic geodesic on $M$ homotopic to $f(\alpha)$.
The hyperbolic length function extends continuously and homogeneously to $\ml$; we denote this extension by $\hypell_x(\lambda)$.
The map 
\[
\teich_{g,m}\ni x\mapsto [\mathcal{S}\ni \alpha\mapsto \hypell_x(\alpha)]\in (\mathbb{R}^{\mathcal{S}}_{\ge 0}-\{0\})/\mathbb{R}_{>0}
\]
is an embedding, and the closure of its image is compact.
We call this compactification of $\teich_{g,m}$ the \emph{Thurston compactification} of $\teich_{g,m}$.
Thurston showed that the compactification is homeomorphic to the closed ball of dimension $2\convgenus$, and the complement of the image of $\teich_{g,m}$ is naturally identified with $\pml$ via \eqref{eq:ml-emb}. 
The boundary $\pml$ of the Thurston compactification is called the \emph{Thurston boundary}.

% \subsection{Action of the mapping class group on $\teich_{g,m}$ and $\pml$}
% The mapping class group $\tmod(g,m)$ is the group of isotopy classes of orientation-preserving self-homeomorphisms of $\Sigma_{g,m}$ that preserve the set of punctures. It acts on $\teich_{g,m}$ by
% \[
% [\omega](M,f)=(M,f\circ \omega^{-1}),
% \]
% where $[\omega]\in \tmod(g,m)$ is represented by a self-homeomorphism $\omega$ of $\Sigma_{g,m}$.

% Under the identification \eqref{eq:ml-emb}, the mapping class group acts on $\ml$ piecewise-linearly. More precisely,
% \begin{equation}
% \label{eq:action_of_mcg-ml}
% i([\omega](\lambda),\alpha)=i(\lambda,[\omega]^{-1}(\alpha))
% \end{equation}
% for $[\omega]\in \tmod(g,m)$, $\lambda\in\ml$, and $\alpha\in\mathcal S$. Thus, $[\omega](\lambda)$ is the measured lamination whose support is the geodesic straightening of $\omega(|\lambda|)$ and whose transverse measure is the pushforward of the transverse measure of $\lambda$ by $\omega$. Since this action commutes with the $\mathbb{R}_{>0}$-action on $\ml$, it induces an action of $\tmod(g,m)$ on $\pml$.
% The action of $\tmod(g,m)$ on $\teich_{g,m}$ extends continuously to the Thurston compactification, and the action on the Thurston boundary coincides with the action on $\pml$.

% \subsubsection*{Teichm\"uller modular group}
% The actions of $\pml$ on $\teich_{g,m}$ and $\pml$ are not effective in few \emph{exceptional cases}. The quotient group $\tmod(g,m)$ of the kernel of natural homomorphism
% \[
% \tmod(g,m)\to \operatorname{Aut}(\teich_{g,m})
% \]
% is called the \emph{Teichm\"uller modular group} of $\Sigma_{g,m}$ (cf. \cite[\S1.3]{ImayoshiTaniguchi1992}).
% The Teichm\"uller modular group acts on $\teich_{g,m}$ and $\pml$ effectively, and the action on $\teich_{g,m}$ is properly discontinuous.
% We assume $\convgenus\ge 1$.
% When $2g-2+m\ge 3$, the kernel is trivial and $\tmod(g,m)$ is isomorphic to $\mcg(\Sigma_{g,m})$. The exceptional cases are $(g,m)=(0,4)$, $(1,1)$, $(1,2)$, and $(2,0)$. Then,
% \begin{align*}
% \tmod(2,0)&\cong \tmod(0,6)\cong \mcg(\Sigma_{2,0})/\mathbb{Z}_2 \\
% \tmod(1,2)&\cong \tmod(0,5)\cong \mcg(\Sigma_{1,2})/\mathbb{Z}_2 \\
% \tmod(1,1)&\cong \tmod(0,4)\cong \operatorname{PSL}(2,\mathbb{Z}),
% \end{align*}
% where the $\mathbb{Z}_2$-actions on the cases are generated by the hyperelliptic involution, which is the center of the mapping class groups in these cases.
% For the proof, See \cite{EarleKra1974}.

\subsection{Action of the mapping class group on $\teich_{g,m}$ and $\pml$}

The mapping class group $\mcg(\Sigma_{g,m})$ consists of the
isotopy classes of orientation-preserving self-homeomorphisms of
$\Sigma_{g,m}$ that preserve the set of punctures. It acts on
$\teich_{g,m}$ by
\begin{equation}
\label{eq:action_of_mcg_Teichmuller}
[\omega](M,f)=(M,f\circ\omega^{-1}),
\end{equation}
where $\omega$ is a representative of
$[\omega]\in\mcg(\Sigma_{g,m})$.

Under the identification \eqref{eq:ml-emb}, the mapping class group
acts on $\ml$ by piecewise-linear homeomorphisms. More precisely,
\begin{equation}
\label{eq:action_of_mcg-ml}
i([\omega](\lambda),\alpha)
=
i(\lambda,[\omega]^{-1}(\alpha))
\end{equation}
for $[\omega]\in\mcg(\Sigma_{g,m})$, $\lambda\in\ml$, and
$\alpha\in\mathcal S$. Equivalently, $[\omega](\lambda)$ is the
measured lamination obtained by geodesically straightening
$\omega(|\lambda|)$ and transporting the transverse measure of
$\lambda$ by $\omega$. Since this action commutes with scalar
multiplication by $\mathbb R_{>0}$, it induces an action of
$\mcg(\Sigma_{g,m})$ on $\pml$.

The action of $\mcg(\Sigma_{g,m})$ on $\teich_{g,m}$ extends
continuously to the Thurston compactification. The induced action on
the Thurston boundary agrees with the action on $\pml$ described
above.

\subsubsection*{Teichm\"uller modular group}
For a few low-complexity surfaces, the actions of the mapping class group
$\mcg(\Sigma_{g,m})$ on $\teich_{g,m}$ and $\pml$ fail to be
faithful. Let
\begin{equation}
\label{eq:mcg_to_tmod}
\mcg(\Sigma_{g,m})
\longrightarrow
\operatorname{Aut}(\teich_{g,m})
\end{equation}
be the natural homomorphism defined by \eqref{eq:action_of_mcg_Teichmuller}.
Let $K_{g,m}$ be the kernel of \eqref{eq:mcg_to_tmod}.
Under the standing assumption $\convgenus\ge1$, this subgroup is
also the kernel of the action on $\pml$. The quotient
\[
\tmod(g,m)
:=
\mcg(\Sigma_{g,m})/K_{g,m}
\]
is called the \emph{Teichm\"uller modular group} of type $(g,m)$, or simply the \emph{modular group} (cf.\ \cite[\S 2]{Bers1981} and \cite[\S 1]{BersBoundary1981}).
When $m=0$, we write $\tmod(g)$ for $\tmod(g,0)$.

The Teichm\"uller modular group $\tmod(g,m)$ acts faithfully on both $\teich_{g,m}$ and $\pml$, and its action on $\teich_{g,m}$ is properly discontinuous.
The term ``Teichm\"uller modular group'' is also used for the
mapping class group itself (cf. \cite{ImayoshiTaniguchi1992}, \cite{Ivanov1992} and \cite{Mumford1967}).

If $2g-2+m\ge3$, then $K_{g,m}$ is trivial, and hence
\[
\tmod(g,m)\cong\mcg(\Sigma_{g,m}).
\]
The exceptional cases are
\[
(g,m)=(0,4),\ (1,1),\ (1,2),\ \text{and }(2,0).
\]
More precisely,
\begin{align*}
K_{2,0}=\langle\iota\rangle,\ 
K_{1,2}=\langle\iota\rangle, \
K_{1,1}=\langle-I\rangle, \
K_{0,4}\cong
(\mathbb Z/2\mathbb Z)^2,
\end{align*}
% \begin{align*}
% K_{2,0}&=\langle\iota\rangle,
% &
% K_{1,2}&=\langle\iota\rangle,\\
% K_{1,1}&=\langle-I\rangle,
% &
% K_{0,4}&\cong
% (\mathbb Z/2\mathbb Z)^2,
% \end{align*}
where $\iota$ denotes the hyperelliptic involution. Consequently,
\begin{align*}
\tmod(2,0)
&\cong
\mcg(\Sigma_{2,0})/\langle\iota\rangle
\cong
\mcg(\Sigma_{0,6})
=
\tmod(0,6),\\
\tmod(1,2)
&\cong
\mcg(\Sigma_{1,2})/\langle\iota\rangle
\cong
\operatorname{Stab}_{\mcg(\Sigma_{0,5})}(p),\\
\tmod(1,1)
&\cong
\mcg(\Sigma_{1,1})/\langle-I\rangle
\cong
\operatorname{PSL}(2,\mathbb Z),\\
\tmod(0,4)
&\cong
\mcg(\Sigma_{0,4})/
(\mathbb Z/2\mathbb Z)^2
\cong
\operatorname{PSL}(2,\mathbb Z),
\end{align*}
where $p$ is the distinguished puncture of $\Sigma_{0,5}$ arising
from the quotient construction. The stabilizer of $p$ has index
five in $\mcg(\Sigma_{0,5})$. For the description of the kernels,
see \cite{BehrstockMargalit2006} and \cite{EarleKra1974}.






\subsection{Thurston measure}
A measured lamination $\lambda\in \ml$ is called \emph{filling} if $i(\lambda,\alpha)>0$ for all $\alpha\in \mathcal{S}$. A nonzero measured lamination $\lambda$ is called \emph{uniquely ergodic} if $i(\lambda,\mu)=0$ for $\mu\in \ml$ implies $\mu=t\lambda$ for some $t\ge 0$.
Let $\ml^{ue}$ be the set of uniquely ergodic measured laminations. We define $\pml^{ue}$ analogously.

It is known that, up to multiplication by a positive constant, there is a unique nonzero locally finite ergodic $\mcg(\Sigma_{g,m})$-invariant Borel measure $\ThursM$ on $\ml$ that gives full measure to the set of filling measured laminations (cf.\ \cite{Masur1985} and \cite{LindenstraussMirzakhani2008}). We call $\ThursM$ the \emph{Thurston measure} on $\ml$.
The set $\ml^{ue}$ has full Thurston measure (cf.\ \cite{Masur1982}).
In train-track coordinates, the Thurston measure belongs to the Lebesgue measure class determined by the canonical Euclidean structure on the space of transverse measures (cf.\ \cite{MoninTelpukhovskiy2019}). The Thurston measure $\ThursM$ satisfies the homogeneity property
\begin{equation}
\label{eq:homogeneous_Thurston_measure}
\ThursM(tE)=t^{2\convgenus}\ThursM(E)
\end{equation}
for any measurable set $E\subset \ml$ and $t\ge 0$, where $tE=\{t\lambda\mid \lambda\in E\}$.

For $x\in\teich_{g,m}$, define
\[
\mathcal{BML}_x=\{\lambda\in\ml\mid\ext_x(\lambda)\le1\},
\]
and set $\mathcal{SML}_x=\partial\mathcal{BML}_x$.
The quotient map $\ml-\{0\}\ni \lambda\mapsto [\lambda]\in \pml$ induces a homeomorphism $\mathcal{SML}_x\to \pml$.
Under this identification, 
we define the \emph{Thurston measure} on $\pml$ with respect to $x\in \teich_{g,m}$ by
\begin{equation}
\label{eq:def_Th_measure}
\PThursM^x(A)=
\dfrac{\ThursM(\{\lambda\in\ml\mid [\lambda]\in A, \ext_x(\lambda)\le 1\})}{\ThursM(\{\lambda\in\ml\mid \ext_x(\lambda)\le 1\})}
\end{equation}
for a Borel set $A\subset \pml$. By definition, each $\PThursM^x$ is a probability measure on $\pml$. For $x,y\in \teich_{g,m}$, $\PThursM^x$ and $\PThursM^y$ are mutually absolutely continuous. In fact, the Radon--Nikodym derivative satisfies
\begin{equation}
\label{eq:Radon-Nikodym}
\dfrac{d\PThursM^y}{d\PThursM^x}([\lambda])=
\left(\frac{\ext_x(\lambda)}{\ext_y(\lambda)}\right)^{\convgenus}
\end{equation}
for $[\lambda]\in \pml$ (cf. \cite[\S2.3.1]{Mirzakhani_etal2012}). Here we use the fact that the normalizing factor $\ThursM(\{\lambda\in\ml\mid\ext_x(\lambda)\le1\})$ is independent of $x$; this is the constancy of the Hubbard--Masur function. For analytically finite surfaces, including punctured surfaces, this follows from \cite[Corollary 15.2]{Miyachi2023}. For closed surfaces, see also \cite[Theorem 5.10]{Dumas2015}.

The following fact is well known; we include a proof for completeness.

\begin{lemma}
\label{lem:quasi-invariance}
Each $\PThursM^x$ is quasi-invariant under the action of the mapping class group.
\end{lemma}

\begin{proof}
Let $[\omega]\in\mcg(\Sigma_{g,m})$ and
$x\in\teich_{g,m}$.
The equivariance of extremal length gives
\[
\ext_{[\omega](x)}([\omega](\lambda))
=\ext_x(\lambda)
\]
for every $\lambda\in\ml$.
Together with the invariance of Thurston measure,
this implies, by \eqref{eq:def_Th_measure}, that
\[
\PThursM^{[\omega](x)}(A)
=\PThursM^x([\omega]^{-1}(A))
=([\omega]_*\PThursM^x)(A)
\]
for every Borel set $A\subset\pml$.
Since $\PThursM^{[\omega](x)}$ and $\PThursM^x$
are mutually absolutely continuous, the claim follows.
\end{proof}

% \begin{proof}
% Let $[\omega]\in \mcg(\Sigma_{g,m})$ and $x\in \teich_{g,m}$.
% From the definition \eqref{eq:def_Th_measure},
% \begin{align*}
% \PThursM^{[\omega](x)}(A)
% &=\dfrac{\ThursM(\{\lambda\in\ml\mid [\lambda]\in A, \ext_{[\omega](x)}(\lambda)\le 1\})}{\ThursM(\{\lambda\in\ml\mid \ext_{[\omega](x)}(\lambda)\le 1\})}
% \\
% &=\dfrac{\ThursM(\{\lambda\in\ml\mid [\lambda]\in A, \ext_{x}([\omega]^{-1}(\lambda))\le 1\})}{\ThursM(\{\lambda\in\ml\mid \ext_{x}([\omega]^{-1}(\lambda))\le 1\})}
% \\
% &=\dfrac{\ThursM(\{\lambda \in\ml\mid [\lambda]\in [\omega]^{-1}(A), \ext_{x}(\lambda)\le 1\})}{\ThursM([\omega](\{\lambda\in\ml\mid \ext_{x}(\lambda)\le 1\}))}
% \\
% &=\dfrac{\ThursM(\{\lambda \in\ml\mid [\lambda]\in [\omega]^{-1}(A), \ext_{x}(\lambda)\le 1\})}{\ThursM(\{\lambda\in\ml\mid \ext_{x}(\lambda)\le 1\})}
% \\
% &=\PThursM^{x}([\omega]^{-1}(A))=[\omega]_*\PThursM^x(A)
% \end{align*}
% for every Borel set $A\subset \pml$,
% since the Thurston measure $\ThursM$ is a $\mcg(\Sigma_{g,m})$-invariant measure on $\ml$. Hence, from \eqref{eq:Radon-Nikodym}, the pushforward measure $[\omega]_*\PThursM^x$ is absolutely continuous with respect to $\PThursM^x$.
% Therefore, the Thurston measure $\PThursM^x$ is quasi-invariant with respect to the action of the mapping class group.
% \end{proof}

The Radon--Nikodym derivative of the pushforward measure $[\omega]_*\PThursM^x$ with respect to $\PThursM^x$ is called the \emph{Jacobian} of $[\omega]$ with respect to $\PThursM^x$ (cf.\ \cite[Appendix D11]{HubbardVol3_2022}).
The preceding proof shows that the Jacobian of $[\omega]$ with respect to $\PThursM^x$ is
\begin{equation}
\label{eq:Radon-Nikodym2}
\dfrac{d([\omega]_*\PThursM^x)}{d\PThursM^x}([\lambda])=
\dfrac{d\PThursM^{[\omega](x)}}{d\PThursM^x}([\lambda])=
\left(\frac{\ext_x(\lambda)}{\ext_{[\omega](x)}(\lambda)}\right)^{\convgenus}
\end{equation}
for $[\lambda]\in \pml$ and $[\omega]\in \mcg(\Sigma_{g,m})$.

The ergodicity of the mapping class group action on $\pml$ with respect to the Lebesgue measure class is due to Masur \cite{Masur1982}; see also Rees \cite{Rees1981} for an alternative proof that applies to surfaces with boundary.
The Thurston measure $\PThursM^x$ is nonatomic. Thus, for each $x\in\teich_{g,m}$, the completion of $(\pml,\PThursM^x)$ is a Lebesgue space in the sense of Rohlin, isomorphic modulo null sets to $[0,1]$ with Lebesgue measure (cf.\ \cite{Rohlin1952}; see also \cite[Theorem 17.41]{Kechris1995}).

\subsection{Busemann functions and horoballs}
The \emph{Busemann function} $\busemann_{\lambda}^{x_0}$ with direction $[\lambda]\in \pml$ and basepoint $x_0\in \teich_{g,m}$ is defined by the limit
\[
\busemann^{x_0}_{\lambda}(x)=\lim_{t\to \infty}(d_T(x,\ray^{\lambda}_{x_0}(t))-t)
\]
(cf. \cite[\S22]{Busemann1955} and \cite[\S3]{EberleinONeill1973}).
By definition, $\busemann^{x_0}_{\lambda}(x_0)=0$, and
\begin{equation}
\label{eq:busemann-basepoint}
\busemann^{[\omega](x_0)}_{\lambda}([\omega](x))
=\busemann^{x_0}_{[\omega]^{-1}(\lambda)}(x)
\end{equation}
Ivanov \cite{Ivanov2001} showed that, for any $x_1,x_2\in \teich_{g,m}$ and $[\lambda]\in\pml$, there exists $D_0>0$ such that $d_T(\ray^\lambda_{x_1}(t),\ray^\lambda_{x_2}(t))\le D_0$ for all $t\ge 0$. Therefore, for another basepoint $x_1\in \teich_{g,m}$,
\begin{equation}
\label{eq:Busemann_function}
|\busemann^{x_0}_{\lambda}(x)-\busemann^{x_1}_{\lambda}(y)|\le D_0+d_T(x,y)
\end{equation}
for all $x,y\in \teich_{g,m}$, where $D_0$ depends on $x_0,x_1$, and $[\lambda]$.
When $[\lambda]$ is uniquely ergodic, the Busemann function is given by
\begin{equation}
\label{eq:busemann_ue}
\busemann^{x_0}_\lambda(x)=\frac{1}{2}(\log \ext_{x}(\lambda)-\log \ext_{x_0}(\lambda))
\end{equation}
(cf. \cite{LiuSu2014} and \cite{Walsh2019}; see also \cite{Miyachi2017}). When $[\lambda]\in \pml^{ue}$, the Busemann function satisfies the following cocycle condition:
\begin{equation}
\label{eq:busemann_cocycle_condition}
\busemann^{x_0}_\lambda(y_0)+\busemann^{y_0}_\lambda(z_0)+\busemann^{z_0}_\lambda(x_0)=0.
\end{equation}
For $r\in \mathbb{R}$, we define the open $r$-\emph{horoball} centered at $[\lambda]\in \pml$ with basepoint $x_0\in \teich_{g,m}$ by
\[
\horos([\lambda],r;x_0)=
\left\{x\in \teich_{g,m}\mid \busemann^{x_0}_\lambda(x)<r\right\}.
\]
Its closure is $\overline{\horos([\lambda],r;x_0)}=\{x\in \teich_{g,m}\mid \busemann^{x_0}_\lambda(x)\le r\}$.
From \eqref{eq:busemann-basepoint},
\begin{equation}
\label{eq:horoball-basepoint}
[\omega](\horos([\lambda],r;x_0))
=\horos([\omega]([\lambda]),r;[\omega](x_0)).
\end{equation}
From the cocycle condition \eqref{eq:busemann_cocycle_condition}, when $[\lambda]\in \pml^{ue}$,
\begin{equation}
\label{eq:horoball-basepoint2}
\horos([\lambda],r;y_0)
=\horos([\lambda],r+\busemann^{x_0}_\lambda(y_0);x_0).
\end{equation}
The boundary $\partial \horos([\lambda],r;x_0)$ is the level set $\{\busemann^{x_0}_\lambda=r\}$ of the Busemann function. We call it the $r$-\emph{horosphere} with basepoint $x_0\in \teich_{g,m}$. When $r=0$, we write $\horos([\lambda];x_0)$ for the $0$-horoball with basepoint $x_0$. The point $x_0$ lies on the boundary of this horoball.

\begin{remark}
In \cite{SuTan2023}, W. Su and D. Tan discuss properties of horoballs and horospheres. In particular, they observe that formula \eqref{eq:busemann_ue} also holds for indecomposable measured laminations (cf. \cite[Proposition 5.7]{SuTan2023}).
\end{remark}


%%%%%%%%%%%%%%%%%%%%%%%%%%%%%%%%%%%%%%%%%%%%%
%%%%%%%%%%%%%%%%%%%%%%%%%%%%%%%%%%%%%%%%%%%%%
%%%%%%%%%%%%%%%%%%%%%%%%%%%%%%%%%%%%%%%%%%%%%
%%%%%%%%%%%%%%%%%%%%%%%%%%%%%%%%%%%%%%%%%%%%%

\section{Dynamics of subgroups}

\subsection{McCarthy--Papadopoulos classification}
McCarthy and Papadopoulos developed a classification of subgroups of mapping class groups modeled on the theory of Kleinian groups in \cite{McCarthyPapadopoulos1989}. We recall their notion of a \emph{sufficiently large subgroup}, which is the mapping class group analogue of a non-elementary Kleinian group.

Although their classification is formulated for subgroups of $\mcg(\Sigma_{g,m})$, we work with subgroups of the Teichm\"uller modular group $\tmod(g,m)$. This amounts to passing to the effective quotient of the actions on $\teich_{g,m}$ and $\pml$, so that both actions are faithful.

\subsubsection{Limit sets}

The definitions of limit points and limit sets are due to McCarthy and Papadopoulos \cite{McCarthyPapadopoulos1989} and Masur \cite{Masur1986}; see also Kent and Leininger \cite[\S3]{KentLeininger2008}.

Let $G$ be a subgroup of $\tmod(g,m)$. A \emph{weak limit point} of the action of $G$ on $\pml$ is a point $[\lambda]\in\pml$ for which there exist $[\lambda']\in\pml$ and an infinite sequence of distinct elements $\{[\omega_n]\}_n\subset G$ such that $[\lambda]=\lim_{n\to\infty}[\omega_n]([\lambda'])$. The \emph{canonical limit set} is the closure of the set of all weak limit points. A \emph{limit set} for $G$ is any closed $G$-invariant subset of the canonical limit set.

We say that $G$ is \emph{dynamically irreducible} if it has a unique nonempty minimal limit set. We call this set the \emph{limit set} of $G$ and denote it by $\Lambda(G)$. We refer to the points of $\Lambda(G)$ as the \emph{limit points} of $G$.

\subsubsection{Sufficiently large subgroups}

A set of pseudo-Anosov mapping classes is said to be \emph{independent} if no two elements of the set have the same fixed-point set in $\pml$. A subgroup $G$ of $\tmod(g,m)$ is called \emph{sufficiently large}\index{sufficiently large} if it contains an independent pair of pseudo-Anosov elements.

Every sufficiently large subgroup is dynamically irreducible. Indeed, let $\Lambda_0(G)\subset\pml$ be the set of fixed points of pseudo-Anosov elements of $G$. If $G$ is sufficiently large, then the closure of $\Lambda_0(G)$ is the unique nonempty minimal closed $G$-invariant subset of $\pml$. Hence, it coincides with the limit set $\Lambda(G)$ (cf.\ \cite[Corollary 4.2]{McCarthyPapadopoulos1989}).

Moreover, suppose that $G$ is a sufficiently large subgroup of $\tmod(g,m)$. If $H<G$ is either an infinite normal subgroup or a subgroup of finite index, then $H$ is also sufficiently large and satisfies $\Lambda(H)=\Lambda(G)$ (cf.\ \cite[Propositions 5.1, 5.4, and 5.5]{McCarthyPapadopoulos1989}).

Let
\[
Z\Lambda(G)=\{[\lambda]\in\pml\mid
i(\lambda,\mu)=0\text{ for some }[\mu]\in\Lambda(G)\}.
\]
McCarthy and Papadopoulos show that if a sequence of distinct points $\{[\omega_n](p)\}_n$ in a $G$-orbit, where $p\in\teich_{g,m}\cup\pml$ and $[\omega_n]\in G$, converges to $[\lambda]\in\pml$, then $[\lambda]\in Z\Lambda(G)$. If, in addition, $[\lambda]\in\pml^{ue}$, then $[\lambda]\in\Lambda(G)$ (cf.\ \cite[Proposition 8.1]{McCarthyPapadopoulos1989}).

\subsubsection*{Examples}
The following examples are discussed in \cite[Example 1 in \S5]{McCarthyPapadopoulos1989}.
\begin{itemize}
\item[(a)]
The Teichm\"uller modular group $\tmod(g,m)$ is sufficiently large. Moreover, its action on $\pml$ is minimal, and hence
\[
\Lambda(\tmod(g,m))=\pml
\]
(cf.\ \cite[Expos\'e 6, VII]{Fathi1979}).
\item[(b)]
For $g\ge 2$, the natural quotient homomorphism \eqref{eq:mcg_to_tmod} restricts to an embedding of the Torelli group $\torelliG_g$.
Indeed, its kernel is trivial when $g\ge 3$. When $g=2$, its kernel is generated by the hyperelliptic involution $\iota$, which acts as $-I$ on $H_1(\Sigma_g,\mathbb Z)$ and therefore does not belong to $\torelliG_g$. We may thus identify $\torelliG_g$ with its image in $\tmod(g)$. Since this image is an infinite normal subgroup of $\tmod(g)$, it follows that $\torelliG_g$ is sufficiently large and
that
$\Lambda(\torelliG_g)=\pml$.
\end{itemize}



\subsection{Conical and horospherical limit points}

Following Kent and Leininger \cite{KentLeininger2008}, we say that a point $[\lambda]\in\pml$ is a \emph{conical limit point} of $G$ if, for every Teichm\"uller geodesic ray $\ray^\lambda_*$ with direction $[\lambda]$, there exist $R>0$ and a $G$-orbit that meets the $R$-neighborhood of $\ray^\lambda_*$ in infinitely many points.
In their definition of conical limit points, Kent and Leininger include the assumption that $[\lambda]\in\Lambda(G)$. It is straightforward to verify that our definition is equivalent to theirs; see \cite{miyachi2026limitsetsmappingclass}. Note that every conical limit point is uniquely ergodic by Masur's criterion \cite{Masur1992}. We denote the set of conical limit points of $G$ by $\Lambda_C(G)$.

% A point $[\lambda]\in\pml$ is called a \emph{big horospherical limit point} of $G$ if there exists a sequence $([\omega_n])_{n=1}^\infty\subset G$ such that
% \begin{equation}
% \label{eq:horo_definition1}
% \frac{\ext_{[\omega_n](x_0)}(\lambda)}{\ext_{x_0}(\lambda)}
% \end{equation}
% remains bounded above and $d_T(x_0,[\omega_n](x_0))\to\infty$ as $n\to\infty$.

We next discuss horospherical limit points. In \cite{miyachi2026limitsetsmappingclass}, Ohshika and the author introduced the following notion. A point $[\lambda]\in\pml$ is called a \emph{small horospherical limit point} of $G$ if there exist a point $x_0\in\teich_{g,m}$ and a sequence $([\omega_n])_{n=1}^\infty\subset G$ such that
\begin{equation}
\label{eq:small-horo_definition}
\busemann^{x_0}_\lambda([\omega_n](x_0))\longrightarrow-\infty
\end{equation}
as $n\to\infty$. 
A point $[\lambda]\in\pml$ is called a \emph{big horospherical limit point} of $G$ if there exist a point $x_0\in\teich_{g,m}$, $b\in \mathbb{R}$ and an infinite sequence $([\omega_n])_{n=1}^\infty\subset G$ of pairwise distinct elements of $G$ such that
\begin{equation}
\label{eq:big_horo_definition}
\busemann^{x_0}_\lambda([\omega_n](x_0))<b.
\end{equation}
for every $n\ge 1$. 

This notion is analogous to its counterpart in the theory of visibility manifolds (cf.\ \cite[\S3]{EberleinONeill1973} and \cite[Section IV]{Sullivan1981}). Geometrically, $[\lambda]\in\pml$ is a small horospherical limit point if and only if every horoball centered at $[\lambda]$ contains infinitely many points of the orbit $G(x_0)$. It follows from \eqref{eq:Busemann_function} that the definition is independent of the choice of basepoint.

We denote the sets of small and big horospherical limit points of $G$ by $\Lambda_h(G)$ and $\Lambda_H(G)$, respectively. We have
\[
\Lambda_C(G)\subset\Lambda_h(G)\subset \Lambda_H(G)
\]
(cf.\ \cite{miyachi2025functiontheorydynamicsergodic}). Suppose that $[\lambda]\in\pml^{ue}$ is a (small) horospherical limit point of $G$. Choose a sequence $([\omega_n])_{n=1}^\infty\subset G$ such that
\begin{equation}
\frac{\ext_{[\omega_n](x_0)}(\lambda)}
{\ext_{x_0}(\lambda)}
\longrightarrow 0.
\end{equation}
Then the sequence $\{[\omega_n](x_0)\}_{n=1}^\infty$ converges to $[\lambda]$ in the Thurston compactification. Indeed, since
\begin{equation*}
e^{-2d_T(x_0,[\omega_n](x_0))}
\ext_{[\omega_n](x_0)}(\lambda)
\longrightarrow 0
\end{equation*}
as $n\to\infty$, and since $\lambda$ is uniquely ergodic, it follows that $[\omega_n](x_0)$ converges to $[\lambda]$ in the Gardiner--Masur compactification (cf.\ \cite[Theorem 1.1]{Miyachi2008} and \cite[Theorem 3]{Miyachi2013}). Consequently, it also converges to $[\lambda]$ in the Thurston compactification (cf.\ \cite[Corollary 1]{Miyachi2013}).

\subsection{Dirichlet points}
We set
\begin{equation}
\label{eq:horo_definition3}
\horo^{x_0}_G([\lambda])
=
\inf_{[\omega]\in G}
\exp\left(2\busemann^{x_0}_\lambda([\omega](x_0))\right).
\end{equation}
By \eqref{eq:busemann_ue}, for $[\lambda]\in\pml^{ue}$, we have
\begin{equation}
\label{eq:horo_definition4}
\horo^{x_0}_G([\lambda])
=
\inf_{[\omega]\in G}
\frac{\ext_{[\omega](x_0)}(\lambda)}
{\ext_{x_0}(\lambda)}.
\end{equation}
A point $[\lambda]\in\pml^{ue}$ is a small horospherical limit point of $G$ if and only if
\begin{equation}
\label{eq:horo_definition2}
\horo^{x_0}_G([\lambda])=0.
\end{equation}
Following the terminology for Kleinian groups, a point $[\lambda]\in\pml$ is called a \emph{Dirichlet point} of $G$ with respect to $x_0$ if $\horo^{x_0}_G([\lambda])>0$ and the infimum defining $\horo^{x_0}_G([\lambda])$ is attained at a point of the orbit $G(x_0)$. 
We denote by $\operatorname{Dir}_G(x_0)$ the set of Dirichlet points of $G$ with respect to $x_0$ and call it the \emph{Dirichlet set} of $G$.
% A point $[\lambda]\in\pml$ is called a \emph{Garnett point} of $G$ if $\horo^{x_0}_G([\lambda])>0$ but this infimum is not attained on $G(x_0)$ (cf.\ \cite[\S2.6]{Nicholls1989}). The Garnett points are contained in $\Lambda_H(G)$.

% \subsection{Horospherical limit points}
% {\color{red}この章は必要か？}

% \begin{proposition}
% \label{prop:horospherical_limit_points}
% Let $G$ be a subgroup of $\mcg(\Sigma_{g,m})$, and let $x_0\in\teich_{g,m}$. For $[\lambda]\in\pml^{ue}$, the following conditions are equivalent:
% \begin{itemize}
% \item[(a)] $[\lambda]\in\Lambda_h(G)$;
% \item[(b)] there exists a sequence $\{[\omega_n]\}_n\subset G$ such that the ratio $\hypell_{[\omega_n](x_0)}(\lambda)/\hypell_{x_0}(\lambda)$ tends to zero as $n\to\infty$, where $\hypell_x$ denotes hyperbolic length on the marked Riemann surface $x\in\teich_{g,m}$; and
% \item[(c)] there exists a sequence $\{[\omega_n]\}_n\subset G$ such that the sequence of intersection-number functions
% \[
% \ml\ni\mu\longmapsto i(\mu,[\omega_n]^{-1}(\lambda))
% \]
% converges uniformly to zero on every compact subset of $\ml$ as $n\to\infty$.
% \end{itemize}
% \end{proposition}

% \begin{proof}
% \noindent
% \textbf{$(a)\Longleftrightarrow(b)$.}
% For every $x\in\teich_{g,m}$, there exists a constant $M_x>0$ such that
% \[
% M_x^{-1}\ext_x(\lambda)
% \le\hypell_x(\lambda)^2
% \le M_x\ext_x(\lambda)
% \]
% for all $\lambda\in\ml$. Indeed, the ratio $\hypell_x(\lambda)^2/\ext_x(\lambda)$ defines a positive continuous function on the compact space $\pml$. By equivariance, the same constant may be used for all points in the orbit of $x_0$. The equivalence of $(a)$ and $(b)$ follows immediately.

% \medskip
% \noindent
% \textbf{$(a)\Longleftrightarrow(c)$.}
% Assume that $(a)$ holds. We may normalize $\lambda$ so that $\ext_{x_0}(\lambda)=1$. Then there exists a sequence $\{[\omega_n]\}_n\subset G$ such that $\ext_{[\omega_n](x_0)}(\lambda)\to 0$ as $n\to\infty$. Let $\epsilon>0$ and let $K\subset\ml$ be compact. Set $M=\max\{1,\max_{\mu\in K}\ext_{x_0}(\mu)\}$. Choose $N>0$ such that $\ext_{[\omega_n](x_0)}(\lambda)<\epsilon^2/M$ whenever $n>N$. Then, for $\mu\in K$ and $n>N$, Minsky's inequality \cite[Lemma 5.1]{Minsky1993} gives
% \begin{align*}
% i(\mu,[\omega_n]^{-1}(\lambda))^2
% &\le \ext_{x_0}(\mu)\ext_{x_0}([\omega_n]^{-1}(\lambda))\\
% &=\ext_{x_0}(\mu)\ext_{[\omega_n](x_0)}(\lambda)\\
% &\le\epsilon^2,
% \end{align*}
% which proves $(c)$.

% Conversely, assume that $(c)$ holds. Let $K=\{\mu\in\ml\mid\ext_{x_0}(\mu)\le 1\}$, which is compact. By assumption, for every $\epsilon>0$, there exists $N>0$ such that $i(\mu,[\omega_n]^{-1}(\lambda))^2<\epsilon$ whenever $n\ge N$ and $\mu\in K$. By the Gardiner--Masur extremal property, extremal length can be expressed as
% \begin{equation}
% \label{eq:tautological_extremal_length}
% \ext_{[\omega_n](x_0)}(\lambda)
% =\ext_{x_0}([\omega_n]^{-1}(\lambda))
% =\max_{[\mu]\in\pml}
% \frac{i(\mu,[\omega_n]^{-1}(\lambda))^2}{\ext_{x_0}(\mu)}
% \le\epsilon
% \end{equation}
% for $n\ge N$ (cf.\ \cite[Theorem 5.1]{Gardiner1991}). Therefore, $\ext_{[\omega_n](x_0)}(\lambda)\to 0$ as $n\to\infty$, and hence $[\lambda]\in\Lambda_h(G)$.
% \end{proof}

% \begin{proposition}
% \label{prop:horospherical_limit_points2}
% Let $G$ be a subgroup of $\mcg(\Sigma_{g,m})$, and let $x_0=(M_0,f_0)\in\teich_{g,m}$. For a uniquely ergodic point $[\lambda]\in\pml$, the following conditions are equivalent:
% \begin{itemize}
% \item[(a)] $[\lambda]\in\Lambda_h(G)$; and
% \item[(b)] there exist sequences $\{t_n\}_n$ and $\{D_n\}_n$ in $[0,\infty)$ satisfying $t_n\nearrow\infty$ and $t_n-D_n\to\infty$, a sequence $\{[\omega_n]\}_n\subset G$, and $e^{2D_n}$-quasiconformal mappings $h_n\colon M_0\to M_{t_n}$ such that $f_{t_n}$ is homotopic to $h_n\circ f_0\circ\omega_n^{-1}$ for every $n$, where $(M_{t_n},f_{t_n})=\ray^\lambda_{x_0}(t_n)$.
% \end{itemize}
% If, in addition, $G$ is torsion-free, condition $(b)$ may be replaced by the following:
% \begin{itemize}
% \item[(b$'$)] there exists an increasing sequence $\{t_n\}_n$ in $[0,\infty)$ with $t_n\nearrow\infty$ such that
% \[
% t_n-d^G_T\bigl(\prG_G(x_0),\prG_G(\ray^\lambda_{x_0}(t_n))\bigr)
% \longrightarrow\infty
% \]
% as $n\to\infty$, where $\prG_G\colon\teich_{g,m}\to\teich_{g,m}/G$ is the quotient projection and $d^G_T$ is the Kobayashi distance on $\teich_{g,m}/G$.
% \end{itemize}
% \end{proposition}

% \begin{proof}
% Suppose first that $(a)$ holds. Choose a sequence $\{[\omega_n]\}_n\subset G$ such that
% \[
% \frac{\ext_{[\omega_n](x_0)}(\lambda)}{\ext_{x_0}(\lambda)}
% \longrightarrow 0
% \]
% as $n\to\infty$. Liu and Su \cite{LiuSu2014} proved that the Gardiner--Masur compactification of $\teich_{g,m}$ is the horofunction compactification with respect to the Teichm\"uller distance. Since $\lambda$ is uniquely ergodic, it is indecomposable and has only one ergodic component. Hence, Walsh's formula for Busemann functions \cite[Theorem 1]{Walsh2019} gives, for each fixed $n\in\mathbb{N}$,
% \[
% \lim_{t\to\infty}
% \left(d_T([\omega_n](x_0),\ray^\lambda_{x_0}(t))-t\right)
% =\frac{1}{2}\log
% \frac{\ext_{[\omega_n](x_0)}(\lambda)}{\ext_{x_0}(\lambda)}
% \]
% (see also \cite{Miyachi2013} and \cite{Miyachi2014}). Therefore, after passing to a subsequence if necessary, we may choose $t_n>0$ with $t_n\nearrow\infty$ such that
% \[
% d_T([\omega_n](x_0),\ray^\lambda_{x_0}(t_n))-t_n
% \longrightarrow-\infty
% \]
% as $n\to\infty$. Set
% \[
% D_n=d_T([\omega_n](x_0),\ray^\lambda_{x_0}(t_n)).
% \]
% Let $h_n\colon M_0\to M_{t_n}$ be the extremal quasiconformal mapping homotopic to $f_{t_n}\circ\omega_n\circ f_0^{-1}$. Then the quasiconformal dilatation of $h_n$ is $e^{2D_n}$, and the preceding estimate gives $t_n-D_n\to\infty$ as $n\to\infty$. Moreover, the homotopy condition implies that $f_{t_n}$ is homotopic to $h_n\circ f_0\circ\omega_n^{-1}$. Thus, $(b)$ holds.

% Conversely, suppose that $(b)$ holds. Since $h_n$ has quasiconformal dilatation $e^{2D_n}$, the extremal length distortion inequality gives
% \[
% \frac{\ext_{[\omega_n](x_0)}(\lambda)}{\ext_{x_0}(\lambda)}\le e^{2D_n}
% \frac{\ext_{\ray^\lambda_{x_0}(t_n)}(\lambda)}{\ext_{x_0}(\lambda)}
% =e^{2D_n-2t_n}
% \longrightarrow 0
% \]
% as $n\to\infty$. This implies that $[\lambda]\in\Lambda_h(G)$.

% Assume, in addition, that $G$ is torsion-free. Then the quotient projection $\prG_G\colon\teich_{g,m}\to\teich_{g,m}/G$ is a holomorphic universal covering map. Since the Teichm\"uller distance is the Kobayashi distance on $\teich_{g,m}$, the Kobayashi distance on $\teich_{g,m}/G$ is the quotient distance induced by the Teichm\"uller distance (cf.\ \cite[Chapter IV, Proposition 1.6]{Kobayashi2005}). Hence,
% \[
% d^G_T\bigl(\prG_G(x_0),\prG_G(\ray^\lambda_{x_0}(t))\bigr)
% =\inf_{[\omega]\in G}
% d_T([\omega](x_0),\ray^\lambda_{x_0}(t)).
% \]
% Therefore, condition $(b)$ is equivalent to condition $(b')$.
% \end{proof}

%%%%%%%%%%%%%%%%%%%%%%%%%%%%%%%%%%%%%%%%%%%%%
%%%%%%%%%%%%%%%%%%%%%%%%%%%%%%%%%%%%%%%%%%%%%
%%%%%%%%%%%%%%%%%%%%%%%%%%%%%%%%%%%%%%%%%%%%%
%%%%%%%%%%%%%%%%%%%%%%%%%%%%%%%%%%%%%%%%%%%%%




% \section{Analytic aspects of the Thurston measure}

% \subsection{Main theorem of this section}
% We call a distance $\distanceml$ on $\ml$ \emph{compatible} if there exists $x_0\in\teich_{g,m}$ such that the Hubbard--Masur map is a bi-Lipschitz homeomorphism from $\mathcal{Q}_{x_0}$, equipped with the distance induced by the $L^1$-norm, onto $(\ml,\distanceml)$. We call such a point $x_0$ a \emph{reference point} for $\distanceml$. For $\lambda_0\in\ml$ and $r>0$, let $B_{ml}(\lambda_0,r)$ denote the $\distanceml$-ball of radius $r$ centered at $\lambda_0$. When necessary, we write $B_{ml}^{\distanceml}(\lambda_0,r)$ to specify the distance. Let $\distanceml_0$ denote the distance induced by the $L^1$-norm on $\mathcal{Q}_{x_0}$ via the Hubbard--Masur map; namely,
% \begin{equation}
% \label{eq:induced_distance}
% \distanceml_0(\lambda_1,\lambda_2)
% =\|q_{\lambda_1,x_0}-q_{\lambda_2,x_0}\|.
% \end{equation}

% Fix a compatible distance $\distanceml$ on $\ml$ with reference point $x_0\in\teich_{g,m}$. Under the identification $\mathcal{SML}_{x_0}\ni\lambda\mapsto[\lambda]\in\pml$, we define a distance on $\pml$ by
% \[
% \hat{\distanceml}([\lambda_1],[\lambda_2])
% =\distanceml(\lambda_1,\lambda_2)
% \quad
% (\lambda_1,\lambda_2\in\mathcal{SML}_{x_0}\subset\ml).
% \]
% We call $\hat{\distanceml}$ a \emph{compatible distance} on $\pml$ with reference point $x_0$. For $[\lambda_0]\in\pml$ and $r>0$, let $B_{pml}([\lambda_0],r)$ denote the $\hat{\distanceml}$-ball of radius $r$ centered at $[\lambda_0]$. We write $B_{pml}^{\distanceml}([\lambda_0],r)$ when the underlying compatible distance on $\ml$ needs to be specified.

% We prove the following theorem.

% \begin{theorem}[The Thurston measure is a Vitali measure]
% \label{thm:Thurston-vitali}
% Let $\distanceml$ be a compatible distance on $\ml$. Then $(\ml,\distanceml,\ThursM)$ is a Vitali metric measure space, and $(\pml,\hat{\distanceml},\PThursM^x)$ is a Vitali metric measure space for every $x\in\teich_{g,m}$.
% \end{theorem}

% The case $\convgenus\ge2$ is proved in \S\ref{subsec:proof_doubling}; the remaining cases are treated in the appendix.


% \medskip
% \paragraph{{\bf Standing assumption.}}
% For the proof in this section, we assume that $\convgenus\ge 2$. Under the standing assumption $\convgenus>0$, the excluded cases are precisely $(g,m)=(0,4)$ and $(1,1)$; these cases are treated separately in the appendix.

% \subsection{Vitali metric measure spaces}

% Let $(X,d,\mu)$ be a metric measure space. A covering $\mathcal{B}$ of a subset $A\subset X$ by closed balls is called \emph{fine} if
% \[
% \inf\{r>0\mid\overline{B}(x,r)\in\mathcal{B}\}=0
% \]
% for every $x\in A$, where $B(x,r)$ denotes the open ball of radius $r$ centered at $x$. We call $(X,d,\mu)$ a \emph{Vitali metric measure space} if, for every subset $A\subset X$ and every fine covering $\mathcal{B}$ of $A$, there exists a pairwise disjoint subcollection $\mathcal{C}\subset\mathcal{B}$ such that $\mu\bigl(A\setminus\bigcup_{B\in\mathcal{C}}B\bigr)=0$. In this case, $\mu$ is called a \emph{Vitali measure} (cf.\ \cite{Tyson2015}).

% Vitali metric measure spaces satisfy the following Lebesgue differentiation theorem (cf.\ \cite[\S3.4]{Tyson2015}).

% \begin{proposition}[Lebesgue differentiation theorem]
% \label{prop:LDT}
% Let $(X,d,\mu)$ be a Vitali metric measure space, let $V$ be a Banach space, and let $f\colon X\to V$ be locally integrable. Then $\mu$-almost every $x\in X$ is a Lebesgue point of $f$; that is,
% \[
% \lim_{r\to 0}
% \frac{1}{\mu(B(x,r))}
% \int_{B(x,r)}\|f(y)-f(x)\|_V\,d\mu(y)=0
% \]
% for $\mu$-almost every $x\in X$. In particular,
% \[
% \lim_{r\to 0}
% \frac{1}{\mu(B(x,r))}
% \int_{B(x,r)}f(y)\,d\mu(y)=f(x)
% \]
% for $\mu$-almost every $x\in X$.
% \end{proposition}

% We use the following criterion for the Vitali property.

% \begin{proposition}[{\cite[Theorem 3.4.3]{Tyson2015}}]
% \label{prop:mms_to_be_vitali}
% Let $(X,d,\mu)$ be a metric measure space such that
% \begin{equation}
% \label{eq:sufficient-condition-Vitali}
% \limsup_{r\to 0}\frac{\mu(B(x,2r))}{\mu(B(x,r))}<\infty
% \end{equation}
% for $\mu$-almost every $x\in X$. Then $(X,d,\mu)$ is a Vitali metric measure space.
% \end{proposition}

% \subsection{Thurston measure versus Euclidean measure}

% Fix $x_0=(M_0,f_0)\in\teich_{g,m}$ and an identification of $\mathcal{Q}_{x_0}$ with $\mathbb{C}^{\convgenus}$. Let $\euclidvolume_{x_0}$ denote the Euclidean (Lebesgue) measure on $\mathcal{Q}_{x_0}$ induced, via this identification, by Lebesgue measure on $\mathbb{C}^{\convgenus}$. The measure $\euclidvolume_{x_0}$ is uniquely determined up to multiplication by a positive constant.

% Let $\QThurs^{x_0}$ denote the pullback of the Thurston measure $\ThursM$ to $\mathcal{Q}_{x_0}$ under the Hubbard--Masur map $\HM\colon\mathcal{Q}_{x_0}\longrightarrow\ml$.

% \begin{proposition}
% \label{prop:mutually_ac}
% The measures $\euclidvolume_{x_0}$ and $\QThurs^{x_0}$ are mutually absolutely continuous.
% \end{proposition}

% \begin{proof}
% Let $\mathcal{Q}_{x_0}^0\subset\mathcal{Q}_{x_0}$ denote the principal stratum consisting of integrable quadratic differentials whose zeros are all simple and which have a simple pole at every puncture, and set $Z_{x_0}=\mathcal{Q}_{x_0}\setminus\mathcal{Q}_{x_0}^0$.

% Let $q_0\in\mathcal{Q}_{x_0}^0$. For closed surfaces, Dumas showed that there exist an open neighborhood $U$ of $q_0$ and a maximal train track $\tau$ such that $\tau$ carries $\HM(q)$ for every $q\in U$ and the restriction of the Hubbard--Masur map to $U$ is a real-analytic diffeomorphism onto its image in the train-track coordinate chart $\ml(\tau)$ (cf.\ \cite[Theorem 5.8]{Dumas2015}).

% Although Dumas states the result for closed surfaces, the same proof applies to punctured surfaces of finite type. Indeed, after filling in the punctures, every integrable quadratic differential extends to a meromorphic quadratic differential with at most simple poles, and each simple pole is a cone singularity of angle $\pi$. The adapted triangulation used in Dumas's proof is constructed from a Delaunay triangulation of the associated singular flat metric (cf.\ \cite[\S5.5]{Dumas2015}). Masur and Smillie construct such Delaunay triangulations for compact flat surfaces with cone-type singularities; in particular, their construction applies to the flat structures arising from finite-norm meromorphic quadratic differentials (cf.\ \cite[\S4]{MasurSmillie1991}). Thus, in the punctured case, the vertices of the triangulation are taken to be all the singularities of the quadratic differential, namely, its zeros and simple poles. The complementary region of the dual train track corresponding to a simple pole is then a punctured monogon. Moreover, on the canonical double cover, the square root of the quadratic differential extends holomorphically across the preimages of the simple poles, so the relative-period and symplectic arguments introduce no additional pole terms. Consequently, the Hubbard--Masur map remains a local real-analytic diffeomorphism on the principal stratum in the punctured case.

% It follows that, for every $q_0\in\mathcal{Q}_{x_0}^0$, there exist a neighborhood $U$ and a train-track coordinate chart $\ml(\tau)$ with the properties above. Since the Thurston measure belongs to the Lebesgue measure class in train-track coordinates on $\ml(\tau)$, the restrictions of $\euclidvolume_{x_0}$ and $\QThurs^{x_0}$ to $\mathcal{Q}_{x_0}^0$ are mutually absolutely continuous. Thus, for every Borel set $E\subset\mathcal{Q}_{x_0}^0$, we have $\euclidvolume_{x_0}(E)=0$ if and only if $\QThurs^{x_0}(E)=0$.

% Let $E\subset\mathcal{Q}_{x_0}$ be a Borel set such that $\QThurs^{x_0}(E)=0$. The set $Z_{x_0}$ consists of the discriminant locus of quadratic differentials with multiple zeros, together with the union, over all punctures, of the hyperplanes of quadratic differentials that are regular at the specified puncture. Hence $Z_{x_0}$ is a proper complex-analytic subset of $\mathcal{Q}_{x_0}$, and therefore $\euclidvolume_{x_0}(Z_{x_0})=0$. Moreover, $\QThurs^{x_0}(E\setminus Z_{x_0})=0$, and hence the mutual absolute continuity on $\mathcal{Q}_{x_0}^0$ implies that $\euclidvolume_{x_0}(E\setminus Z_{x_0})=0$. Therefore,
% \[
% \begin{aligned}
% \euclidvolume_{x_0}(E)
% &=\euclidvolume_{x_0}(E\setminus Z_{x_0})
% +\euclidvolume_{x_0}(E\cap Z_{x_0})
% =0.
% \end{aligned}
% \]
% Thus, $\euclidvolume_{x_0}$ is absolutely continuous with respect to $\QThurs^{x_0}$.

% We now prove absolute continuity in the other direction. Let $\ml^{ec}$ denote the set of essentially complete measured laminations. By definition, a measured lamination is essentially complete if every component of the complement of its support is either an ideal triangle or a punctured monogon (cf.\ \cite[\S9.5]{Thurston-LectureNote}). In every train-track chart $\ml(\tau)$, the set $\ml^{ec}\cap\ml(\tau)$ contains the complement of a countable union of kernels of linear functions with rational coefficients (cf.\ \cite[Proposition 9.5.11]{Thurston-LectureNote}). It follows that $\ml^{ec}$ has full $\ThursM$-measure. Since $q_{\lambda,x_0}\in\mathcal{Q}_{x_0}^0$ for every $\lambda\in\ml^{ec}$, we have
% \[
% \QThurs^{x_0}(Z_{x_0})
% =\ThursM\bigl(\HM(Z_{x_0})\bigr)
% \le\ThursM(\ml\setminus\ml^{ec})
% =0.
% \]

% Now let $E\subset\mathcal{Q}_{x_0}$ be a Borel set such that $\euclidvolume_{x_0}(E)=0$. Then $\euclidvolume_{x_0}(E\setminus Z_{x_0})=0$, so the mutual absolute continuity on $\mathcal{Q}_{x_0}^0$ gives $\QThurs^{x_0}(E\setminus Z_{x_0})=0$. Together with $\QThurs^{x_0}(Z_{x_0})=0$, this yields
% \[
% \QThurs^{x_0}(E)
% =\QThurs^{x_0}(E\setminus Z_{x_0})
% +\QThurs^{x_0}(E\cap Z_{x_0})
% =0.
% \]
% Therefore, $\QThurs^{x_0}$ is absolutely continuous with respect to $\euclidvolume_{x_0}$, completing the proof.
% \end{proof}

% \subsection{Volume asymptotics}

% Let $\distanceml$ be a compatible distance on $\ml$ with reference point $x_0\in\teich_{g,m}$. For $q_0\in\mathcal{Q}_{x_0}$ and $r>0$, let $B_{\mathcal{Q}_{x_0}}(q_0,r)\subset\mathcal{Q}_{x_0}$ denote the ball of radius $r$ centered at $q_0$ with respect to the $L^1$-norm. Since $\distanceml$ is compatible, there exists $K\ge 1$ such that
% \begin{equation}
% \label{eq:volume_comparison1}
% B^{\distanceml_0}_{ml}(\lambda,r/K)
% \subset B^{\distanceml}_{ml}(\lambda,r)
% \subset B^{\distanceml_0}_{ml}(\lambda,Kr)
% \end{equation}
% for every $r>0$ and $\lambda\in\ml$.

% \begin{lemma}
% \label{lem:volume_asymptotic}
% Suppose that a compatible distance $\distanceml$ on $\ml$ is $K$-bi-Lipschitz equivalent to $\distanceml_0$. For every $\epsilon>0$ and every $\lambda\in\ml$ with $q_{\lambda,x_0}\in\mathcal{Q}_{x_0}^0\setminus\{0\}$, there exist constants $A>0$ and $r_0>0$ such that
% \begin{align*}
% \left(\frac{A}{K^{2\convgenus}}-\epsilon\right)r^{2\convgenus}
% &\le\ThursM\left(B^{\distanceml}_{ml}(\lambda,r)\right)\le\left(AK^{2\convgenus}+\epsilon\right)r^{2\convgenus}
% \end{align*}
% for every $0<r\le r_0$. The constant $A$ depends on $x_0$ and $\lambda$, whereas $r_0$ may depend on $x_0$, $\lambda$, $\epsilon$, and $K$.
% \end{lemma}

% \begin{proof}
% Set $q_0=q_{\lambda,x_0}$. We first consider the case $\distanceml=\distanceml_0$. We claim that there exist $A>0$ and $r_0>0$ such that
% \begin{equation}
% \label{eq:doubling-3}
% \left|
% \ThursM\left(B_{ml}^{\distanceml_0}(\lambda,r)\right)
% -Ar^{2\convgenus}
% \right|
% \le\epsilon r^{2\convgenus}
% \end{equation}
% for every $0<r\le r_0$.

% Since $q_0\in\mathcal{Q}_{x_0}^0$, as discussed in the proof of \Cref{prop:mutually_ac}, there exists a neighborhood $U$ of $q_0$ such that the Hubbard--Masur map $\HM\colon U\to\ml(\tau)\subset\ml$ is a real-analytic diffeomorphism onto its image and $\HM(q_0)=\lambda$. Choose $r_0>0$ sufficiently small that $B_{\mathcal{Q}_{x_0}}(q_0,r_0)\subset U$. Since the distance induced by the $L^1$-norm and the Euclidean measure $\euclidvolume_{x_0}$ are invariant under translation,
% \[
% \euclidvolume_{x_0}\left(
% B_{\mathcal{Q}_{x_0}}(q_{\lambda,x_0},r)
% \right)
% =\euclidvolume_{x_0}\left(
% B_{\mathcal{Q}_{x_0}}(0,1)
% \right)r^{2\convgenus}
% \]
% for every $r>0$. On $U$, write $d\QThurs^{x_0}=J(q)\,d\euclidvolume_{x_0}$. The train-track coordinate description above shows that the density $J$ is positive and continuous near $q_0$. Therefore,
% \begin{align*}
% \ThursM\left(B_{ml}^{\distanceml_0}(\lambda,r)\right)
% &=\QThurs^{x_0}\left(B_{\mathcal{Q}_{x_0}}(q_{\lambda,x_0},r)\right)\\
% &=J(q_0)\euclidvolume_{x_0}\left(B_{\mathcal{Q}_{x_0}}(q_0,r)\right)
% +o(r^{2\convgenus})\\
% &=Ar^{2\convgenus}+o(r^{2\convgenus})
% \end{align*}
% as $r\to 0$, where $A=J(q_0)\euclidvolume_{x_0}(B_{\mathcal{Q}_{x_0}}(0,1))$. This proves \eqref{eq:doubling-3}. The desired comparison follows by applying \eqref{eq:doubling-3} to the two outer balls in \eqref{eq:volume_comparison1}, after decreasing $r_0$ and adjusting the error term if necessary.
% \end{proof}

% \subsection{Proof of \Cref{thm:Thurston-vitali}}
% \label{subsec:proof_doubling}

% Let $\distanceml$ be a compatible distance on $\ml$ with reference point $x_0\in\teich_{g,m}$. We first show that $(\ml,\distanceml,\ThursM)$ is a Vitali metric measure space.

% Suppose that $\distanceml$ is $K$-bi-Lipschitz equivalent to $\distanceml_0$. By \Cref{lem:volume_asymptotic}, for every $\lambda\in\ml$ with $q_{\lambda,x_0}\in\mathcal{Q}_{x_0}^0\setminus\{0\}$ and every $k>0$,
% \begin{equation}
% \label{eq:doubling-1}
% \limsup_{r\to 0}
% \frac{\ThursM(B_{ml}^{\distanceml}(\lambda,kr))}
% {\ThursM(B_{ml}^{\distanceml}(\lambda,r))}
% \le K^{4\convgenus}k^{2\convgenus}.
% \end{equation}
% By \Cref{prop:mutually_ac}, the set $\HM(\mathcal{Q}_{x_0}^0)\subset\ml$ has full $\ThursM$-measure. Therefore, \Cref{prop:mms_to_be_vitali} implies that $(\ml,\distanceml,\ThursM)$ is a Vitali metric measure space.

% We next show that $(\pml,\hat{\distanceml},\PThursM^x)$ is a Vitali metric measure space for every $x\in\teich_{g,m}$. As above, it suffices to verify that, for every $\lambda\in\mathcal{SML}_{x_0}$ with $q_{\lambda,x_0}\in\mathcal{Q}_{x_0}^0$ and every $k>0$,
% \begin{equation}
% \label{eq:doubling-2}
% \limsup_{r\to 0}
% \frac{\PThursM^x(B^{\distanceml}_{pml}([\lambda],kr))}
% {\PThursM^x(B^{\distanceml}_{pml}([\lambda],r))}
% <\infty.
% \end{equation}
% Since the Radon--Nikodym derivative in \eqref{eq:Radon-Nikodym} is positive and continuous, for every $[\lambda_0]\in\pml$ we have
% \begin{equation}
% \label{eq:change2}
% \PThursM^x\left(B_{pml}^{\distanceml}([\lambda_0],r)\right)
% =\left\{
% \left(
% \frac{\ext_{x_0}(\lambda_0)}{\ext_x(\lambda_0)}
% \right)^{\convgenus}
% +o(1)
% \right\}
% \PThursM^{x_0}\left(B_{pml}^{\distanceml}([\lambda_0],r)\right)
% \end{equation}
% as $r\to 0$, where $\lambda_0$ is any representative of $[\lambda_0]$. Thus, it suffices to consider the case $x=x_0$.

% We first assume that $\hat{\distanceml}$ is induced by $\distanceml_0$; see \eqref{eq:induced_distance}. Fix $\lambda_0\in\mathcal{SML}_{x_0}$ such that $q_0=q_{\lambda_0,x_0}\in\mathcal{Q}_{x_0}^0$. For $0<r<1$, define
% \begin{align*}
% B_{\lambda_0}^S(r)
% &=B^{\distanceml_0}_{ml}(\lambda_0,r)
% \cap\mathcal{SML}_{x_0},\\
% E_{\lambda_0}(r)
% &=\{t\lambda\in\ml\mid
% \lambda\in B_{\lambda_0}^S(r),\ 1-r\le t\le 1+r\},\\
% C_{\lambda_0}(r)
% &=\{t\lambda\in\ml\mid
% \lambda\in B_{\lambda_0}^S(r),\ 0\le t\le 1\}.
% \end{align*}
% The set $B^{\distanceml_0}_{pml}([\lambda_0],r)$ is the image of $B_{\lambda_0}^S(r)$ under the projection $\mathcal{SML}_{x_0}\to\pml$. We claim that
% \begin{equation}
% \label{eq:Balls_Sectors}
% B^{\distanceml_0}_{ml}(\lambda_0,r/4)
% \subset E_{\lambda_0}(r)
% \subset B^{\distanceml_0}_{ml}(\lambda_0,4r).
% \end{equation}

% Indeed, let $\lambda\in E_{\lambda_0}(r)$. There exist $\lambda'\in B_{\lambda_0}^S(r)$ and $t\in[1-r,1+r]$ such that $\lambda=t\lambda'$. Since $q_{t\lambda',x_0}=t^2q_{\lambda',x_0}$, we have
% \[
% \begin{aligned}
% \distanceml_0(\lambda_0,\lambda)
% &\le\distanceml_0(\lambda_0,\lambda')
% +\distanceml_0(\lambda',t\lambda')\le r+|t^2-1|\\
% &\le 3r+r^2\le 4r.
% \end{aligned}
% \]

% Conversely, let $\lambda\in B^{\distanceml_0}_{ml}(\lambda_0,r/4)$. Write $\lambda=t\lambda'$ with $\lambda'\in\mathcal{SML}_{x_0}$ and $t>0$. Since $\|q_{\lambda_0,x_0}\|=1$ and $q_{\lambda,x_0}=t^2q_{\lambda',x_0}$, the reverse triangle inequality gives $|t^2-1|<r/4$. Thus,
% \[
% \begin{aligned}
% \distanceml_0(t\lambda_0,\lambda)
% &\le\distanceml_0(t\lambda_0,\lambda_0)
% +\distanceml_0(\lambda_0,\lambda)\\
% &=|t^2-1|+\distanceml_0(\lambda_0,\lambda)
% <r/2.
% \end{aligned}
% \]
% If $t\ge 1$, then
% \begin{align*}
% \distanceml_0(\lambda_0,\lambda')
% &=\|q_{\lambda_0,x_0}-q_{\lambda',x_0}\|
% \le t^2\|q_{\lambda_0,x_0}-q_{\lambda',x_0}\|
% =\|q_{t\lambda_0,x_0}-q_{t\lambda',x_0}\|\\
% &=\distanceml_0(t\lambda_0,\lambda)
% <r/2<r.
% \end{align*}
% If $t<1$, then $t^2>1-r/4$ and
% \[
% \distanceml_0(\lambda_0,\lambda')
% =\frac{1}{t^2}\distanceml_0(t\lambda_0,t\lambda')
% <\frac{r/2}{1-r/4}
% <r.
% \]
% Thus, $\lambda'\in B_{\lambda_0}^S(r)$. Moreover, $|t-1|<r/4$, and hence $1-r<t<1+r$. Therefore, $\lambda\in E_{\lambda_0}(r)$, proving \eqref{eq:Balls_Sectors}.

% We also have
% \begin{align*}
% \ThursM(C_{\lambda_0}(r))
% &=\PThursM^{x_0}\left(
% B^{\distanceml_0}_{pml}([\lambda_0],r)
% \right)\ThursM(\mathcal{BML}_{x_0}),
% \end{align*}
% and, up to a set of Thurston measure zero,
% \[
% E_{\lambda_0}(r)
% =(1+r)C_{\lambda_0}(r)
% \setminus
% (1-r)\operatorname{Int}(C_{\lambda_0}(r)).
% \]
% Moreover, up to a set of Thurston measure zero,
% \[
% E_{\lambda_0}(r)
% =
% (1+r)C_{\lambda_0}(r)
% \setminus
% (1-r)C_{\lambda_0}(r).
% \]
% Indeed, the symmetric difference between these two sets is precisely
% \[
% (1-r)B_{\lambda_0}^S(r).
% \]
% This set has zero Thurston measure. To see this, recall that
% $q_{t\lambda,x_0}=t^2q_{\lambda,x_0}$ and that
% $\|q_{\lambda,x_0}\|=1$ for every
% $\lambda\in\mathcal{SML}_{x_0}$. Hence the inverse image of
% $(1-r)B_{\lambda_0}^S(r)$ under the Hubbard--Masur map is contained
% in the norm sphere
% \[
% \left\{q\in\mathcal{Q}_{x_0}\mid \|q\|=(1-r)^2\right\}.
% \]
% This sphere has zero Euclidean measure. Therefore,
% \Cref{prop:mutually_ac} implies that
% $(1-r)B_{\lambda_0}^S(r)$ has zero Thurston measure.
% Applying the homogeneity of the Thurston measure in
% \eqref{eq:homogeneous_Thurston_measure}, we obtain
% \begin{align*}
% \ThursM(E_{\lambda_0}(r))
% &=\left(
% (1+r)^{2\convgenus}-(1-r)^{2\convgenus}
% \right)\ThursM(C_{\lambda_0}(r))\\
% &=\left(4\convgenus r+O(r^3)\right)
% \PThursM^{x_0}\left(
% B_{pml}^{\distanceml_0}([\lambda_0],r)
% \right)
% \ThursM(\mathcal{BML}_{x_0})
% \end{align*}
% as $r\to 0$, where $2\convgenus=\dim_{\mathbb{R}}\ml$.

% It follows from \eqref{eq:Balls_Sectors} that
% \[
% \begin{aligned}
% \frac{\ThursM\left(
% B^{\distanceml_0}_{ml}(\lambda_0,r/4)
% \right)}{\ThursM(\mathcal{BML}_{x_0})}
% &\le\left(4\convgenus r+O(r^3)\right)
% \PThursM^{x_0}\left(
% B_{pml}^{\distanceml_0}([\lambda_0],r)
% \right)\\
% &\le\frac{\ThursM\left(
% B^{\distanceml_0}_{ml}(\lambda_0,4r)
% \right)}{\ThursM(\mathcal{BML}_{x_0})}.
% \end{aligned}
% \]
% Consequently, for all sufficiently small $r>0$,
% \begin{equation}
% \label{eq:volume-asymptotic2}
% \frac{4k\convgenus r+O(r^3)}{4\convgenus r+O(r^3)}
% \frac{\PThursM^{x_0}\left(
% B^{\distanceml_0}_{pml}([\lambda_0],kr)
% \right)}
% {\PThursM^{x_0}\left(
% B^{\distanceml_0}_{pml}([\lambda_0],r)
% \right)}
% \le
% \frac{\ThursM\left(
% B^{\distanceml_0}_{ml}(\lambda_0,4kr)
% \right)}
% {\ThursM\left(
% B^{\distanceml_0}_{ml}(\lambda_0,r/4)
% \right)}.
% \end{equation}
% This proves \eqref{eq:doubling-2} when $\distanceml=\distanceml_0$.

% Now suppose that $\hat{\distanceml}$ is induced by an arbitrary compatible distance $\distanceml$ on $\ml$. The bi-Lipschitz comparison gives
% \[
% B^{\distanceml_0}_{pml}([\lambda_0],r/K)
% \subset B^{\distanceml}_{pml}([\lambda_0],r)
% \subset B^{\distanceml_0}_{pml}([\lambda_0],Kr).
% \]
% Thus, \eqref{eq:doubling-2} follows from the case of $\distanceml_0$ and \eqref{eq:volume-asymptotic2}. Since the set of points represented by $\lambda\in\mathcal{SML}_{x_0}$ with $q_{\lambda,x_0}\in\mathcal{Q}_{x_0}^0$ has full $\PThursM^x$-measure, \Cref{prop:mms_to_be_vitali} implies that $(\pml,\hat{\distanceml},\PThursM^x)$ is a Vitali metric measure space. This completes the proof of \Cref{thm:Thurston-vitali} when $\convgenus\ge 2$.

% \begin{remark}
% Equation \eqref{eq:change2} also implies that, if a function $f$ on $\pml$ is integrable with respect to one of the measures $\PThursM^x$, and hence with respect to all of them, then the set of Lebesgue points of $f$ is independent of the choice of basepoint $x\in\teich_{g,m}$.
% \end{remark}

% \begin{remark}
% The proofs of \eqref{eq:doubling-1} and \eqref{eq:doubling-2} also show that, when $\ml$ and $\pml$ are equipped with compatible distances having reference point $x_0\in\teich_{g,m}$,
% \begin{align*}
% \liminf_{r\to 0}
% \frac{\ThursM(B_{ml}(\lambda,kr))}
% {\ThursM(B_{ml}(\lambda,r))}
% &>0,\\
% \liminf_{r\to 0}
% \frac{\PThursM^x(B_{pml}([\lambda],kr))}
% {\PThursM^x(B_{pml}([\lambda],r))}
% &>0
% \end{align*}
% for every $k>0$ and every $\lambda\in\ml$ such that $q_{\lambda,x_0}\in\mathcal{Q}_{x_0}^0\setminus\{0\}$.
% \end{remark}

%%%%%%%%%%%%%%%%%%%%%%%%%%%%%%%%%%%%%%%%%%%%%
%%%%%%%%%%%%%%%%%%%%%%%%%%%%%%%%%%%%%%%%%%%%%
%%%%%%%%%%%%%%%%%%%%%%%%%%%%%%%%%%%%%%%%%%%%%
%%%%%%%%%%%%%%%%%%%%%%%%%%%%%%%%%%%%%%%%%%%%%
\section{Level loci of extremal lengths}

\subsection{Level loci}

For $x_1,x_2\in\teich_{g,m}$ and $c\in \mathbb{R}$, we define the \emph{level loci of extremal lengths} by
\begin{align*}
\ml_c(x_1,x_2)
&=\{\lambda\in\ml\mid \ext_{x_1}(\lambda)=e^c\ext_{x_2}(\lambda)\},\\
\pml_c(x_1,x_2)
&=\{[\lambda]\in\pml\mid \ext_{x_1}(\lambda)=e^c\ext_{x_2}(\lambda)\}.
\end{align*}
The second definition is independent of the choice of representative $\lambda$ because extremal length is homogeneous of degree two.
When $c=0$, we simply write $\ml(x_1,x_2)=\ml_0(x_1,x_2)$ and $\pml(x_1,x_2)=\pml_0(x_1,x_2)$.
 
We prove the following theorem.

\begin{theorem}
\label{thm:level_set}
If $x_1\ne x_2$, then $\ml_c(x_1,x_2)$ is null with respect to $\ThursM$, and $\pml_c(x_1,x_2)$ is null with respect to $\PThursM^x$ for every $x\in\teich_{g,m}$.
\end{theorem}

\begin{remark}
For the proof in this section, we assume that $\convgenus\ge2$. Under the standing assumption $\convgenus>0$, the excluded cases are precisely $(g,m)=(0,4)$ and $(1,1)$; these cases are treated separately in the appendix.
\end{remark}


\subsection{Lemmas}
\label{subsec:lemmas}

The \emph{rational depth} of $\lambda\in\ml$ is the dimension of the space of rational linear functions that vanish at $\lambda$ in a train-track coordinate chart containing $\lambda$ (cf.\ \cite[Definition 9.5.10]{Thurston-LectureNote}). If a measured lamination has rational depth zero, then it is essentially complete (cf.\ \cite[Proposition 9.5.11]{Thurston-LectureNote}). Since the Thurston measure belongs to the Lebesgue measure class in each train-track chart, and since there is a countable train-track atlas, the set of measured laminations of rational depth zero has full $\ThursM$-measure in $\ml$.

\begin{lemma}[Local real analyticity of the
Hubbard--Masur map]
\label{lem:HM_local_analytic}
Let $x\in\teich_{g,m}$, and let $q_0\in\mathcal Q_x\setminus \{0\}$
have only simple zeros and a simple pole at every
puncture.
Then there exist a neighborhood $U$ of $q_0$ and a
train-track coordinate chart such that the
Hubbard--Masur map restricted to $U$ is a
real-analytic diffeomorphism onto an open subset
of that chart.
\end{lemma}

\begin{proof}
For closed surfaces, this follows from Dumas's
adapted train-track construction
\cite[Lemma 5.7 and Theorem 5.8]{Dumas2015}.

For punctured surfaces, we use the same construction
with simple poles included among the singularities
of the flat metric.
After filling in the punctures, an integrable
quadratic differential extends meromorphically
with at most simple poles.
A simple pole is a cone singularity of angle $\pi$.
The Delaunay triangulation construction of
Masur and Smillie applies to these finite-area
flat metrics
\cite[\S4]{MasurSmillie1991}.
The complementary region of the adapted train track
corresponding to a simple pole is a punctured monogon.

On the canonical double cover, the square root of
the quadratic differential extends holomorphically
across the preimages of the simple poles.
Thus the relative-period and symplectic arguments
in Dumas's proof apply without additional pole terms,
giving the required local real-analytic
diffeomorphism.
\end{proof}

\begin{lemma}[Common train track]
\label{lem:common_splitting}
Let $x_1,x_2\in\teich_{g,m}$, and let $\lambda_0\in\ml$ be a measured lamination of rational depth zero. Then there exist a complete train track $\tau$
(in the sense of \cite{PennerHarer1992}), a strictly positive transverse measure $\mu_0\in E(\tau)$ representing $\lambda_0$, and a neighborhood $U$ of $\mu_0$ in $W(\tau)$ such that $U\subset E(\tau)$ and for $i=1,2$, the map
\[
U\ni\mu\longmapsto q_{R_\tau(\mu),x_i}\in\mathcal Q_{x_i}
\]
is a real-analytic diffeomorphism onto an open neighborhood of $q_{\lambda_0,x_i}$, where $R_\tau\colon E(\tau)\to\ml$ is the realizing map defined in \eqref{eq:realizing_map}.
\end{lemma}

\begin{proof}
Set $q_i=q_{\lambda_0,x_i}$ for $i=1,2$.
Since $\lambda_0$ has rational depth zero, it is
essentially complete, and each $q_i$ belongs to
the principal stratum.
By \Cref{lem:HM_local_analytic}, choose positively
measured generic train tracks $(\tau_i,\mu_i)$
representing $\lambda_0$ and neighborhoods $U_i$
of $\mu_i$ such that
\[
U_i\ni\mu\longmapsto q_{R_{\tau_i}(\mu),x_i}
\]
is a real-analytic diffeomorphism onto an open neighborhood of $q_i$.

By \cite[Theorem 2.8.5]{PennerHarer1992}, the measured train tracks $(\tau_1,\mu_1)$ and $(\tau_2,\mu_2)$ are equivalent. The common-subtrack theorem therefore yields a positively measured generic train track $(\tau_3,\mu_3)$ obtained from each $(\tau_i,\mu_i)$ by a finite sequence of measured splits, after performing shifts if necessary (cf.\ \cite[Theorem 2.3.1]{PennerHarer1992}).

No collision can occur in either splitting sequence. Indeed, at the first collision, the equality of the two relevant branch weights would give a nontrivial rational linear relation, independent of the switch equations, for the weights representing $\lambda_0$. Pulling this relation back through the preceding shifts and non-collision splits would imply that $\lambda_0$ has positive rational depth, contrary to our assumption. Consequently, the dimensions of $W(\tau_3)$ and $W(\tau_i)$ agree.

Let $C_i\colon W(\tau_3)\to W(\tau_i)$ be the linear carrying map induced by the splitting sequence. Since carrying does not change the represented measured lamination, we have $R_{\tau_i}\circ C_i=R_{\tau_3}$. The injectivity of $R_{\tau_3}$ implies that $C_i$ is injective. Since the source and target have the same dimension, $C_i$ is an isomorphism.

Choose a sufficiently small connected neighborhood $U$ of $\mu_3$ contained in $C_1^{-1}(U_1)\cap C_2^{-1}(U_2)$ and in the positive weight cone. For $i=1,2$, the map
\[
U\ni\mu\longmapsto q_{R_{\tau_3}(\mu),x_i}
=q_{R_{\tau_i}(C_i(\mu)),x_i}
\]
is then a real-analytic diffeomorphism onto an open neighborhood of $q_i$. In particular, $R_{\tau_3}(U)$ is an open neighborhood of $\lambda_0$ in $\ml$. By the standard characterization of complete train tracks, $\tau_3$ is therefore complete (cf.\ \cite[Propositions 1.4.9 and 1.7.6]{PennerHarer1992}). Taking $\tau=\tau_3$ and $\mu_0=\mu_3$ proves the lemma.
\end{proof}

\begin{lemma}
\label{lem:coincidence}
Let $x_1,x_2\in\teich_{g,m}$ and $c\in\mathbb R$. If there is a nonempty open set $V\subset\ml$ such that $\ext_{x_1}(\lambda)=e^c\ext_{x_2}(\lambda)$ for every $\lambda\in V$, then $c=0$ and $x_1=x_2$.
\end{lemma}

\begin{proof}
Choose a measured lamination $\lambda_0\in V$ of rational depth zero. By \Cref{lem:common_splitting}, there exist a complete train track $\tau$, a strictly positive measure $\mu_0\in E(\tau)$ representing $\lambda_0$, and a neighborhood $U$ of $\mu_0$ such that $R_\tau(U)\subset V$ and the conclusions of that lemma hold.

Suppose first that $c=0$.
Since $\lambda_0$ is essentially complete, it is maximal. For each $i=1,2$, let $\Sigma_{\lambda_0,\tau}(x_i)\in W(\tau)$ be the element defined in \cite[\S6.2]{MiyachiOhshika2017}. Let $\omega_{Th}$ denote the Thurston symplectic form on $W(\tau)$.
The differential formula for extremal length gives
\[
d_{\mu_0}(\ext_{x_i}\circ R_\tau)(\xi)
=2\omega_{Th}\bigl(\xi,\Sigma_{\lambda_0,\tau}(x_i)\bigr)
\]
for every $\xi\in W(\tau)$ (cf.\ \cite[Theorem 1.1]{Miyachi2013b} and \cite[(1.1) and (6.5)]{MiyachiOhshika2017}). Since the two extremal length functions agree on $R_\tau(U)$, their differentials at $\mu_0$ agree. The nondegeneracy of the Thurston symplectic form therefore implies that
\[
\Sigma_{\lambda_0,\tau}(x_1)
=\Sigma_{\lambda_0,\tau}(x_2).
\]
The map $x\mapsto\Sigma_{\lambda_0,\tau}(x)$ is an embedding by \cite[Th\'eor\`eme 6.1]{MiyachiOhshika2017}. Hence $x_1=x_2$.

Now suppose that $c\ne0$. The same differential argument gives
\[
\Sigma_{\lambda_0,\tau}(x_1)
=e^c\Sigma_{\lambda_0,\tau}(x_2).
\]
Let $F_i$ be the horizontal foliation of $q_{\lambda_0,x_i}$.
From the proof of \cite[Th\'eor\`eme 6.1]{MiyachiOhshika2017}, we have
\begin{align*}
i(\gamma,F_1)=\omega_{Th}(\gamma,\Sigma_{\lambda_0,\tau}(x_1))=e^c\omega_{Th}(\gamma,\Sigma_{\lambda_0,\tau}(x_2))
=e^ci(\gamma,F_2)=i(\gamma,e^cF_2)
\end{align*}
for sufficiently many simple closed curves $\gamma\in E(\tau)$, and hence $F_1=e^cF_2$.
Thus $x_1$ and $x_2$ lie on a Teichm\"uller geodesic
with direction $[\lambda_0]$, and
$d_T(x_1,x_2)=|c|/2>0$.

The same argument applies to every measured lamination
of rational depth zero in $V$.
By the uniqueness of the Teichm\"uller geodesic joining
the distinct points $x_1$ and $x_2$, the projective
classes of all such laminations must belong to the
two-element set determined by the vertical and horizontal
foliations of this geodesic.
Consequently, the set of measured laminations of rational
depth zero in $V$ is contained in the union of at most
two rays in $\ml$.

Each of these rays has zero Thurston measure.
On the other hand, measured laminations of rational
depth zero have full Thurston measure in $V$, and the
nonempty open set $V$ has positive Thurston measure.
This is a contradiction.
% Thus $x_1$ and $x_2$ lie on the same Teichm\"uller geodesic with direction $\lambda_0$ and $d_T(x_1,x_2)=|c|/2$. Moreover,
% \[
% \frac{\ext_{x_1}(F_1)}{\ext_{x_2}(F_2)}=\frac{i(\lambda_0,F_1)}{i(\lambda_0,F_2)}
% =e^c=e^{2\epsilon_c\,d_T(x_1,x_2)},
% \]
% where $\epsilon_c$ is the sign of $c$.
% Therefore, the level locus $\ml_c(x_1,x_2)$ coincides with the ray $\{tF_1\mid t\ge 0\}\subset \ml$, which cannot contain a nonempty open set. This is a contradiction.
\end{proof}

\subsection{Proof of \Cref{thm:level_set}}
Fix distinct points $x_1,x_2\in\teich_{g,m}$ and
$c\in\mathbb R$.
Let $\mathcal R\subset\ml$ be the full-measure set
of measured laminations of rational depth zero.

For each $\lambda_0\in\mathcal R$,
\Cref{lem:common_splitting} gives a complete train
track $\tau$ and a connected open neighborhood $U$
of the weight representing $\lambda_0$ such that
both maps
\[
\mu\longmapsto q_{R_\tau(\mu),x_i},
\qquad i=1,2,
\]
are real analytic.
After shrinking $U$, we may assume that both
$q_{R_\tau(\mu),x_1}$ and $q_{R_\tau(\mu),x_2}$
belong to the principal stratum for every $\mu\in U$.
The function
\[
F(\mu)
=
\ext_{x_1}(R_\tau(\mu))
-e^c\ext_{x_2}(R_\tau(\mu))
\]
is therefore real analytic on $U$, since the
$L^1$-norm is real analytic on the principal stratum
(cf.\ \cite[\S\S5.3--5.4]{Dumas2015}).

By \Cref{lem:coincidence}, $F$ cannot vanish
identically on $U$.
Its zero set consequently has Lebesgue measure
zero \cite{Mityagin2020}.
Since the Thurston measure belongs to the Lebesgue
measure class in train-track coordinates,
\[
\ThursM\bigl(
\ml_c(x_1,x_2)\cap R_\tau(U)
\bigr)=0.
\]

The open sets $R_\tau(U)$ cover $\mathcal R$.
Since $\ml$ is second countable, a countable
subcollection covers $\mathcal R$.
Thus $\ml_c(x_1,x_2)\cap\mathcal R$ is null.
Since $\ml\setminus\mathcal R$ is also null,
$\ml_c(x_1,x_2)$ has zero Thurston measure.

Finally, extremal length is homogeneous of degree
two, so the cone over $\pml_c(x_1,x_2)$ inside
$\mathcal{BML}_x$ is contained in $\ml_c(x_1,x_2)$.
By the definition of the projectivized Thurston
measure in \eqref{eq:def_Th_measure}, we obtain
\[
\PThursM^x(\pml_c(x_1,x_2))=0.
\]
This completes the proof when $\convgenus\ge2$.

%%%%%%%%%%%%%%%%%%%%%%%%%%%%%%%%%%%%%%%%%%%%%
%%%%%%%%%%%%%%%%%%%%%%%%%%%%%%%%%%%%%%%%%%%%%
%%%%%%%%%%%%%%%%%%%%%%%%%%%%%%%%%%%%%%%%%%%%%
%%%%%%%%%%%%%%%%%%%%%%%%%%%%%%%%%%%%%%%%%%%%%

% \section{A dichotomy in the dynamics of subgroups}

% In this section, we characterize the conservative and dissipative sets for actions of torsion-free subgroups of the mapping class group. We also characterize the dissipative part under the weaker assumption that the stabilizer of a chosen basepoint is trivial.

% \subsection{Conservativity and dissipativity}

% Let $(X,P)$ be a probability space equipped with an action of a countable group $\Gamma$ by nonsingular transformations; that is, the measure $P$ is quasi-invariant under the action. The action of $\Gamma$ on a $\Gamma$-invariant measurable subset $Y\subset X$ is said to be \emph{conservative} if, for every measurable subset $B\subset Y$ with $P(B)>0$, there are infinitely many $\gamma\in\Gamma$ such that
% $P(B\cap\gamma(B))>0$.
% The action of $\Gamma$ on a $\Gamma$-invariant measurable subset $Z\subset X$ is said to be \emph{dissipative} if there exists a measurable subset $B\subset Z$ such that
% $P(B\cap\gamma(B))=0$ for every $\gamma\in\Gamma\setminus\{\mathrm{id}\}$,
% and
% \begin{align*}
% \bigcup_{\gamma\in\Gamma}\gamma(B)=Z
% \end{align*}
% (cf.\ \cite[Definition D11.3]{HubbardVol3_2022}).

% For a subgroup $G$ of the mapping class group, we define
% \begin{align*}
% \conservativepart(G)
% &=\conservativepart_x(G)
% =\left\{[\lambda]\in\pml\mathrel{}\middle|\mathrel{}
% \sum_{[\omega]\in G}
% \left(
% \frac{\ext_x(\lambda)}
% {\ext_{[\omega](x)}(\lambda)}
% \right)^{\convgenus}
% =\infty
% \right\},\\
% \dissipativepart(G)
% &=\dissipativepart_x(G)
% =\left\{[\lambda]\in\pml\mathrel{}\middle|\mathrel{}
% \sum_{[\omega]\in G}
% \left(
% \frac{\ext_x(\lambda)}
% {\ext_{[\omega](x)}(\lambda)}
% \right)^{\convgenus}
% <\infty
% \right\}.
% \end{align*}
% It follows from \eqref{eq:QC-inv_extremal_length} that the convergence or divergence of these series is independent of the choice of $x\in\teich_{g,m}$. Thus, the sets $\conservativepart_x(G)$ and $\dissipativepart_x(G)$ are independent of $x$. Nevertheless, we retain the subscript $x$ when it is useful to indicate the choice of basepoint explicitly.

% \begin{theorem}[Dichotomy]
% \label{thm:Hopf_decomposition}
% Let $G\le\mcg(\Sigma_{g,m})$ be a torsion-free subgroup. Then the following dichotomy holds.
% \begin{enumerate}
% \item The action of $G$ on $\pml$ is conservative on $\conservativepart(G)$.
% \item The action of $G$ on $\pml$ is dissipative on $\dissipativepart(G)$.
% \end{enumerate}
% \end{theorem}

% \begin{proof}
% Fix $x\in\teich_{g,m}$. Since $G$ is torsion-free and the stabilizer of every point of Teichm\"uller space in the mapping class group is finite, $\operatorname{Stab}_G(x)$ is trivial. Hence, by \Cref{lem:Kleinian-like}, the action of $G$ on $(\pml,\PThursM^x)$ is Kleinian-like. The result follows from \cite[Proposition D11.5]{HubbardVol3_2022} and \eqref{eq:Radon-Nikodym2}; the torsion-free assumption is used in the identification of the conservative part.
% \end{proof}

% The following is a direct analogue of Sullivan's characterization of the dissipative part of the measurable action of a Kleinian group (cf.\ \cite[Theorem III]{Sullivan1981}).

% \begin{theorem}[Dirichlet points and Dissipative set]
% \label{thm:Dissipative_part}
% Let $G$ be a nontrivial subgroup of the mapping class group, and let $x_0\in\teich_{g,m}$. Suppose that $[\omega](x_0)\ne x_0$ for every $[\omega]\in G\setminus\{\mathrm{id}\}$. Then the symmetric difference
% \[
% \dissipativepart_{x_0}(G)\bigtriangleup \bigcup_{[\omega]\in G}[\omega](\Dirichlet^\infty_G(x_0))
% \]
% is null with respect to $\PThursM^{x_0}$.
% \end{theorem}
% The set $\Dirichlet^\infty_G(x_0)$ in the statement of \Cref{thm:Dissipative_part} is the ideal boundary of the Dirichlet polyhedron $\Dirichlet_G(x_0)$ for $G$ centered at $x_0$, which will be discussed in \S\ref{subsec:DirichletDomain_Ideal}.
% The proof of \Cref{thm:Dissipative_part} is given in \S\ref{subsec:Proof-Dissipative_part}.

% \subsection{Ideal boundary of the Dirichlet polyhedron}
% \label{subsec:DirichletDomain_Ideal}
% In \cite{McCarthyPapadopoulos1996}, McCarthy and Papadopoulos discuss fundamental domains for subgroups of the mapping class group that are analogous to Dirichlet polyhedra for hyperbolic manifolds (see, e.g., \cite{Beardon1995}).

% Let $G\le\mcg(\Sigma_{g,m})$, and let $x_0\in\teich_{g,m}$ satisfy $\operatorname{Stab}_G(x_0)=\{\mathrm{id}\}$. The \emph{Dirichlet polyhedron} $\Dirichlet_G(x_0)$ for $G$, centered at $x_0$, is defined by
% \[
% \Dirichlet_G(x_0)
% =\left\{x\in\teich_{g,m}\mathrel{}\middle|\mathrel{}
% d_T(x,x_0)<d_T(x,[\omega](x_0))
% \text{ for every }[\omega]\in G\setminus\{\mathrm{id}\}\right\}.
% \]
% McCarthy and Papadopoulos prove the following:
% \begin{itemize}
% \item The closure of the Dirichlet polyhedron satisfies
% \[
% \overline{\Dirichlet_G(x_0)}
% =\left\{x\in\teich_{g,m}\mathrel{}\middle|\mathrel{}
% d_T(x,x_0)\le d_T(x,[\omega](x_0))
% \text{ for every }[\omega]\in G\right\},
% \]
% and $x_0$ is an interior point of this closure (\cite[Propositions 5.1 and 5.2]{McCarthyPapadopoulos1996}).
% \item The Dirichlet polyhedron $\Dirichlet_G(x_0)$ and its closure are starlike with respect to $x_0$ (\cite[Proposition 5.3]{McCarthyPapadopoulos1996}).
% \item As in the case of hyperbolic spaces (see \cite[Chapter 9]{Beardon1995}), $\Dirichlet_G(x_0)$ is a locally finite fundamental domain in $\teich_{g,m}$ (\cite[Theorem 5.9]{McCarthyPapadopoulos1996}).
% \item The boundary $\partial \Dirichlet_G(x_0)$ is a properly embedded, locally flat submanifold of codimension one in $\teich_{g,m}$, and every boundary point is contained in at least one hypersurface of the form
% \[
% L_{[\omega]}(x_0)
% =\{x\in \teich_{g,m}\mid d_T(x,x_0)=d_T(x,[\omega](x_0))\}
% \]
% for some $[\omega]\in G\setminus\{\mathrm{id}\}$ (\cite[Propositions 5.2 and 5.5]{McCarthyPapadopoulos1996}).
% \end{itemize}
% A $[\omega]$-\emph{face} $F_{[\omega]}(x_0)$ of $\Dirichlet_G(x_0)$ is the subset of $\partial \Dirichlet_G(x_0)$ consisting of points that lie on $L_{[\omega]}(x_0)$ but on no other such hypersurface. The union of the faces is a dense open subset of $\partial \Dirichlet_G(x_0)$ (\cite[Proposition 5.16]{McCarthyPapadopoulos1996}).
% Let $\Dirichlet^\infty_G(x_0)\subset \pml$ denote the set of accumulation points of $\Dirichlet_G(x_0)$ in the Thurston compactification.

% \begin{lemma}
% \label{lem:Dirichlet_geodesic_horoball}
% Let $[\lambda]\in\pml^{ue}$, and consider the following conditions:
% \begin{itemize}
% \item[(a)] The Teichm\"uller ray $\ray^{\lambda}_{x_0}$ is contained in $\Dirichlet_G(x_0)$.
% \item[(b)] $[\lambda]\in \Dirichlet^\infty_G(x_0)$.
% \item[(c)] The $0$-horoball $\horos([\lambda];x_0)$ is disjoint from the $G$-orbit $G(x_0)$ of $x_0$.
% \end{itemize}
% Condition ${\rm (a)}$ implies ${\rm (c)}$, and conditions ${\rm (a)}$ and ${\rm (b)}$ are equivalent. Moreover, if $\overline{\horos([\lambda];x_0)}\cap G(x_0)=\{x_0\}$, then ${\rm (a)}$ holds.
% \end{lemma}

% \begin{proof}
% \noindent
% \textbf{${\rm (a)}\Rightarrow {\rm (c)}$.}
% Let $[\omega]\in G$. Since $\ray^\lambda_{x_0}(t)\in \Dirichlet_G(x_0)$ for every $t\ge 0$,
% \[
% d_T(\ray^\lambda_{x_0}(t),[\omega](x_0))
% \ge d_T(\ray^\lambda_{x_0}(t),x_0)=t,
% \]
% and hence
% \[
% \busemann^{x_0}_\lambda([\omega](x_0))
% =\lim_{t\to\infty}
% \bigl(d_T(\ray^\lambda_{x_0}(t),[\omega](x_0))-t\bigr)
% \ge 0.
% \]
% Thus, $[\omega](x_0)$ lies outside the $0$-horoball.

% We next prove the additional assertion. Assume that $x_0$ is the only point of the orbit contained in the closure of the $0$-horoball, and let $[\omega]\in G\setminus\{\mathrm{id}\}$. By assumption,
% $\busemann^{x_0}_\lambda([\omega](x_0))>0$. The function
% $t\mapsto d_T([\omega](x_0),\ray^\lambda_{x_0}(t))-t$ is nonincreasing and converges to $\busemann^{x_0}_\lambda([\omega](x_0))$. Consequently, for every $t\ge0$,
% \[
% d_T([\omega](x_0),\ray^\lambda_{x_0}(t))-t
% \ge\busemann^{x_0}_\lambda([\omega](x_0))>0.
% \]
% Since this holds for every $[\omega]\in G\setminus\{\mathrm{id}\}$, we have $\ray^\lambda_{x_0}(t)\in\Dirichlet_G(x_0)$ for every $t\ge0$. Thus, ${\rm (a)}$ holds.

% \medskip
% \noindent
% \textbf{${\rm (b)}\Rightarrow {\rm (a)}$.}
% Since $[\lambda]\in \Dirichlet^\infty_G(x_0)$, there exists a sequence $\{x_n\}_{n=1}^\infty\subset \Dirichlet_G(x_0)$ such that $x_n\to[\lambda]$ in the Thurston compactification. Set $t_n=d_T(x_0,x_n)$. Then $x_n$ also converges to $[\lambda]$ in the Gardiner--Masur compactification, and
% \[
% e^{-2t_n}\ext_{x_n}(\mu)
% \longrightarrow \frac{i(\mu,\lambda)^2}{\ext_{x_0}(\lambda)}
% \]
% uniformly on compact subsets of $\ml$ (cf.\ \cite[Corollaries 1 and 2]{Miyachi2013}).

% For each $n\in\mathbb{N}$, choose $[\lambda_n]\in\pml$ such that
% $\ray^{\lambda_n}_{x_0}(t_n)=x_n$. Choose representatives $\lambda_n$ normalized by $\ext_{x_0}(\lambda_n)=1$. Then $t_n\to\infty$, and, after passing to a subsequence, we may assume that $\lambda_n$ converges to some $\lambda_0\in\ml$. Moreover,
% \[
% e^{-2t_n}\ext_{x_n}(\lambda_n)
% =e^{-4t_n}\ext_{x_0}(\lambda_n)\longrightarrow 0.
% \]
% Hence $i(\lambda_0,\lambda)=0$, and thus $[\lambda_0]=[\lambda]$. By Teichm\"uller's theorem, the geodesic rays $\ray^{\lambda_n}_{x_0}$ converge uniformly on compact subsets of $[0,\infty)$ to $\ray^\lambda_{x_0}$.

% Suppose, for contradiction, that the ray is not contained in $\Dirichlet_G(x_0)$. Since $x_0\in\Dirichlet_G(x_0)$ and the ray is continuous, there exists $t_0>0$ such that $\ray^\lambda_{x_0}(t_0)\in\partial\Dirichlet_G(x_0)$. Then there exists $[\omega]\in G\setminus\{\mathrm{id}\}$ such that
% \[
% d_T(\ray^\lambda_{x_0}(t_0),x_0)
% =d_T(\ray^\lambda_{x_0}(t_0),[\omega](x_0)).
% \]
% By \cite[Lemma 4.4]{McCarthyPapadopoulos1996},
% \[
% d_T(x_0,\ray^\lambda_{x_0}(t_0+1))
% >d_T([\omega](x_0),\ray^\lambda_{x_0}(t_0+1)).
% \]
% Therefore, for all sufficiently large $n$, we have $t_0+1<t_n$ and
% $\ray^{\lambda_n}_{x_0}(t_0+1)\notin\Dirichlet_G(x_0)$. This contradicts the starlikeness of $\Dirichlet_G(x_0)$ with respect to $x_0$, because $x_n\in\Dirichlet_G(x_0)$.

% \medskip
% \noindent
% \textbf{${\rm (a)}\Rightarrow {\rm (b)}$.}
% Since $\lambda$ is uniquely ergodic, $\ray^\lambda_{x_0}(t)\to[\lambda]$ in the Thurston compactification (cf.\ \cite{Masur1980}). By ${\rm (a)}$, every point of this ray lies in $\Dirichlet_G(x_0)$; hence $[\lambda]\in\Dirichlet^\infty_G(x_0)$.
% \end{proof}

% \begin{lemma}
% \label{lem:almost-unique}
% Assume that $[\omega](x_0)\ne x_0$ for every $[\omega]\in G\setminus\{\mathrm{id}\}$.
% For $\PThursM^{x_0}$-almost every $[\lambda]\in\Dirichlet^\infty_G(x_0)$, we have $\overline{\horos([\lambda];x_0)}\cap G(x_0)=\{x_0\}$.
% \end{lemma}

% \begin{proof}
% Let $[\lambda]\in\Dirichlet^\infty_G(x_0)\cap\pml^{ue}$.
% Suppose that $[\omega](x_0)\in\overline{\horos([\lambda];x_0)}$ for some $[\omega]\in G\setminus\{\mathrm{id}\}$. By \Cref{lem:Dirichlet_geodesic_horoball}, conditions ${\rm (a)}$ and ${\rm (b)}$ are equivalent for $[\lambda]$, and condition ${\rm (a)}$ implies ${\rm (c)}$. Hence $[\omega](x_0)$ lies on the boundary of the $0$-horoball. Therefore,
% \begin{align*}
% \frac{\ext_{[\omega](x_0)}(\lambda)}{\ext_{x_0}(\lambda)}
% &=\exp\left(2\busemann^{x_0}_\lambda([\omega](x_0))\right)\\
% &=\exp\left(2\busemann^{x_0}_\lambda(x_0)\right)=1.
% \end{align*}
% Thus, $[\lambda]\in\pml(x_0,[\omega](x_0))$. Since $G$ is countable, \Cref{thm:level_set} shows that the union of these level loci over $[\omega]\in G\setminus\{\mathrm{id}\}$ is null. The complement of $\pml^{ue}$ is also null. The assertion follows.
% \end{proof}

% \subsection{Proof of \Cref{thm:Dissipative_part}}
% \label{subsec:Proof-Dissipative_part}

%%%%%%%%%%%%%%%%%%%%%%%%%%%%%%%%%%%%%%%%%%%%%
%%%%%%%%%%%%%%%%%%%%%%%%%%%%%%%%%%%%%%%%%%%%%
%%%%%%%%%%%%%%%%%%%%%%%%%%%%%%%%%%%%%%%%%%%%%
%%%%%%%%%%%%%%%%%%%%%%%%%%%%%%%%%%%%%%%%%%%%%




\section{The Hopf decomposition for subgroup actions}
\label{sec:Hopf_decomposition}
In this section, we prove \Cref{thm:Hopf_decomposition_main,thm:main_classfication_limit_sets}.
We first recall the measure-theoretic terminology and state a refinement
of the description of the dissipative part in terms of the ideal boundary
of a Dirichlet polyhedron. We then establish the preliminary lemmas and
prove the assertions in turn. All null sets on $\pml$ are understood
with respect to the Thurston measure class.

\subsection{Conservativity and dissipativity}
\label{subsec:conservativity_dissipativity}
Let $(X,P)$ be a standard probability space, completed if necessary.
Suppose that a countable group $\Gamma$ acts on $X$ by nonsingular
transformations; that is, $P$ is quasi-invariant under the action.

A measurable subset $B\subset X$ is called a \emph{wandering set} if $P(B\cap\gamma(B))=0$ for every $\gamma\in\Gamma\setminus\{\mathrm{id}\}$.
The action on a $\Gamma$-invariant measurable subset $Z\subset X$ is called \emph{dissipative} if there exists a measurable wandering set $B\subset Z$ such that
\[
P\left(Z\setminus\bigcup_{\gamma\in\Gamma}\gamma(B)\right)=0.
\]
The action on a $\Gamma$-invariant measurable subset $Y\subset X$ is called \emph{conservative} if, for every measurable subset $B\subset Y$ of positive measure, $P(B\cap\gamma(B))>0$ for infinitely many $\gamma\in\Gamma$.

For a subgroup $G\le\tmod(g,m)$, we call its action
\emph{conservative} or \emph{totally dissipative} if its action
on the whole space $\pml$ is conservative or dissipative,
respectively.
For background, see \cite{Pommerenke1976}, \cite{Sullivan1981} and \cite{Nicholls1989}. See also \cite{Kaimanovich2010} and \cite[Definition D11.3]{HubbardVol3_2022}.

\begin{remark}
Our terminology differs slightly from that of \cite{Kaimanovich2010}. Our notion of conservativity corresponds to his notion of \emph{infinite conservativity}, whereas our notion of dissipativity corresponds to his notion of \emph{complete dissipativity}. In his terminology, an action is conservative if every measurable set of positive measure is recurrent.
\end{remark}

The action of a group $\Gamma$ on a measurable space $(X,\mu)$ is said to be \emph{essentially free} if the fixed point set $\operatorname{Fix}(\gamma)=\{x\in X\mid \gamma(x)=x\}$ is null for all $\gamma\in \Gamma\setminus \{\mathrm{id}\}$.
For an essentially free action, Kaimanovich's conservative and infinitely conservative parts agree modulo null sets. This distinction therefore disappears for the subgroup actions on $\pml$ considered below, by \Cref{lem:Fix_null}.

\subsection{Proof strategy and a refinement of the main results}
\label{subsec:Main_resut_of_SectionDichotomy}
We use the sets $\conservativepart(G)=\conservativepart_x(G)$ and
$\dissipativepart(G)=\dissipativepart_x(G)$ defined in the Introduction
by divergence and convergence, respectively, of the sum of the
Radon--Nikodym derivatives. Their independence of $x$ and their
$G$-invariance are proved in \Cref{lem:Cons-Diss-ElementaryProperty}.

The proof of \Cref{thm:Hopf_decomposition_main} is given in
\S\ref{subsec:ProofOfDecomposition}.
The identification of the conservative part with $\Lambda_H(G)$,
which is the first assertion of \Cref{thm:main_classfication_limit_sets},
is proved in \S\ref{subsec:Proof_of_ConservativePart_BigHorospherical}.
For the dissipative part, we prove the following stronger statement,
which also describes it in terms of the ideal boundary of a Dirichlet
polyhedron. It is analogous to Sullivan's description for Kleinian
groups (cf.\ \cite[Corollary in Section IV]{Sullivan1981}).

\begin{theorem}[The dissipative part and Dirichlet polyhedra]
\label{thm:Dissipative_part}
Let $G$ be a nontrivial subgroup of the Teichm\"uller modular group, and let $x_0\in\teich_{g,m}$. Suppose that $\operatorname{Stab}_G(x_0)=\{\mathrm{id}\}$. Then the symmetric differences
\[
\dissipativepart(G)\bigtriangleup \bigcup_{[\omega]\in G}[\omega](\Dirichlet^\infty_G(x_0)),\quad
\dissipativepart(G)\bigtriangleup \operatorname{Dir}_G(x_0)
\]
are null with respect to $\PThursM^{x}$ for all $x\in \teich_{g,m}$.
\end{theorem}

The set $\Dirichlet^\infty_G(x_0)$ appearing in \Cref{thm:Dissipative_part} is the ideal boundary of the Dirichlet polyhedron $\Dirichlet_G(x_0)$ centered at $x_0$; it is discussed in \S\ref{subsec:DirichletDomain_Ideal}. We show that $\Dirichlet^\infty_G(x_0)$ is a wandering set and that almost every point of $\Dirichlet^\infty_G(x_0)$ is a Dirichlet point (cf. \Cref{lem:almost-unique}).
The proof of \Cref{thm:Dissipative_part} is given in \S\ref{subsec:Proof-Dissipative_part}.

\begin{remark}
The proof of \cite[Proposition D11.5]{HubbardVol3_2022} shows that, if $\operatorname{Stab}_G(x)=\{\mathrm{id}\}$, then the action of $G$ on $\dissipativepart(G)$ is dissipative with respect to $\PThursM^x$. Although that proposition is stated for torsion-free groups, its proof of dissipativity only requires the stabilizer at the basepoint to be trivial in our case. Moreover, \Cref{thm:Hopf_decomposition_main} shows that $\dissipativepart(G)$ is the maximal dissipative invariant subset modulo null sets.
\end{remark}

\subsubsection{A characterization of conservative actions}
As a consequence of \Cref{thm:Hopf_decomposition_main,thm:main_classfication_limit_sets,thm:Dissipative_part}, we obtain the following characterization of conservative subgroup actions, analogous to \cite[Corollary in Section IV]{Sullivan1981}.

\begin{corollary}[A characterization of conservative actions]
\label{coro:conservativity_Dirichletset}
Let $G$ be a nontrivial subgroup of the Teichm\"uller modular group, and let $x_0\in\teich_{g,m}$ satisfy $\operatorname{Stab}_G(x_0)=\{\mathrm{id}\}$. 
For every $x\in\teich_{g,m}$, the following five conditions are equivalent:
\begin{itemize}
\item[(a)] the action of $G$ on $(\pml,\PThursM^{x})$ is conservative;
\item[(b)] $\Lambda_H(G)$ has full $\PThursM^{x}$-measure;
\item[(c)] $\dissipativepart(G)$ has zero $\PThursM^{x}$-measure;
\item[(d)] $\operatorname{Dir}_G(x_0)$ has zero $\PThursM^{x}$-measure;
\item[(e)] $\Dirichlet^\infty_G(x_0)$ has zero $\PThursM^{x}$-measure.
\end{itemize}
\end{corollary}

\begin{proof}
By \Cref{thm:Hopf_decomposition_main}, conditions (a) and (c) are
equivalent. The first assertion of
\Cref{thm:main_classfication_limit_sets} gives the equivalence with
(b), and \Cref{thm:Dissipative_part} gives the equivalence with (d).
Finally, a countable union of translates of
$\Dirichlet^\infty_G(x_0)$ is null if and only if
$\Dirichlet^\infty_G(x_0)$ is null, by nonsingularity of the action.
Another application of \Cref{thm:Dissipative_part} proves the
equivalence with (e).
\end{proof}

\subsection{Preliminary lemmas}
We first establish two lemmas needed for the proof of \Cref{thm:Hopf_decomposition_main}.
\begin{lemma}
\label{lem:Cons-Diss-ElementaryProperty}
Let $G$ be a subgroup of the Teichm\"uller modular group.
\begin{itemize}
\item[(a)]
The sets $\conservativepart(G)$ and $\dissipativepart(G)$ are independent of the choice of basepoint $x\in \teich_{g,m}$.
\item[(b)]
$\conservativepart(G)$ and $\dissipativepart(G)$ are invariant under
the action of $G$.
\end{itemize}
\end{lemma}

\begin{proof}
\textbf{(a).}
By \eqref{eq:QC-inv_extremal_length} and the invariance of $d_T$ under the action of $G$, we have
\[
e^{-4\convgenus d_T(x,y)}
\left(
\frac{\ext_x(\lambda)}
{\ext_{[\omega](x)}(\lambda)}
\right)^{\convgenus}
\le
\left(
\frac{\ext_y(\lambda)}
{\ext_{[\omega](y)}(\lambda)}
\right)^{\convgenus}
\le
e^{4\convgenus d_T(x,y)}
\left(
\frac{\ext_x(\lambda)}
{\ext_{[\omega](x)}(\lambda)}
\right)^{\convgenus}
\]
for all $x,y\in\teich_{g,m}$ and $[\omega]\in G$.
Summing over $[\omega]\in G$ shows that convergence or divergence
of the defining series is independent of the choice of basepoint.

\medskip
\noindent
\textbf{(b).}
For $[\eta]\in G$, the equivariance of extremal length gives
\begin{align*}
&\sum_{[\omega]\in G}
\left(
\frac{\ext_x([\eta](\lambda))}
     {\ext_{[\omega](x)}([\eta](\lambda))}
\right)^{\convgenus} =
\sum_{[\omega]\in G}
\left(
\frac{\ext_{[\eta]^{-1}(x)}(\lambda)}
     {\ext_{[\eta]^{-1}\circ[\omega](x)}(\lambda)}
\right)^{\convgenus}=
\sum_{[\gamma]\in G}
\left(
\frac{\ext_{[\eta]^{-1}(x)}(\lambda)}
     {\ext_{[\gamma]\circ[\eta]^{-1}(x)}(\lambda)}
\right)^{\convgenus},
\end{align*}
where $[\gamma]=[\eta]^{-1}\circ[\omega]\circ[\eta]$.
Since the convergence or divergence of the defining series is independent of the choice of basepoint, this proves~(b).
\end{proof}

\begin{lemma}[Essential freeness]
\label{lem:Fix_null}
The action of $\tmod(g,m)$ on $\pml$ is essentially free.
Namely, for $[\omega]\in \tmod(g,m)\setminus \{\mathrm{id}\}$, 
\[
\operatorname{Fix}([\omega])
=
\{[\lambda]\in\pml\mid
[\omega]([\lambda])=[\lambda]\}
\]
is $\PThursM^x$-null for every $x\in\teich_{g,m}$.
\end{lemma}

\begin{proof}
Suppose first that $[\omega]$ has finite order. Since the action of $\tmod(g,m)$ on $\teich_{g,m}$ is faithful and $[\omega]\ne\mathrm{id}$, we may choose $x_0\in\teich_{g,m}$ such that $[\omega](x_0)\ne x_0$.
We also choose $k\ge1$ such that $[\omega]^k=\mathrm{id}$. If $[\lambda]\in\operatorname{Fix}([\omega])$, then
\[
[\omega](\lambda)=c\lambda
\]
for some $c>0$. Since $[\omega]^k(\lambda)=c^k\lambda=\lambda$, we have
$c=1$.

By equivariance of extremal length,
\begin{align*}
\operatorname{Fix}([\omega])
&\subset
\left\{
[\lambda]\in\pml\mathrel{}\middle|\mathrel{}
\ext_{x_0}(\lambda)
=
\ext_{x_0}([\omega](\lambda))
\right\} \\
&=
\left\{
[\lambda]\in\pml\mathrel{}\middle|\mathrel{}
\ext_{x_0}(\lambda)
=
\ext_{[\omega]^{-1}(x_0)}(\lambda)
\right\} \\
&=
\pml\bigl(x_0,[\omega]^{-1}(x_0)\bigr).
\end{align*}
This set is $\PThursM^{x_0}$-null by \Cref{thm:level_set}. Since the
measures $\PThursM^x$ are mutually absolutely continuous, it is
$\PThursM^x$-null for every $x\in\teich_{g,m}$.

Next, suppose that $[\omega]$ has infinite order. If $[\omega]$ is
pseudo-Anosov, then it has exactly two fixed points in $\pml$, namely
the projective classes of its stable and unstable measured laminations.
Since $\PThursM^x$ has no atoms, its fixed-point set is null.

Suppose that $[\omega]$ is reducible. Choose $n\ge1$ such that
$[\varphi]=[\omega]^n$ fixes an essential reduction class $\alpha$ (cf. \cite[\S2]{BirmanLubotzkyMcCarthy1983}).
Since $[\omega]$ has infinite order, $[\varphi]$ still acts nontrivially
on $\teich_{g,m}$. Moreover,
\[
\operatorname{Fix}([\omega])
\subset \operatorname{Fix}([\varphi]).
\]

Let $[\lambda]\in\operatorname{Fix}([\varphi])$, and write
\[
[\varphi](\lambda)=c\lambda
\]
for some $c>0$. Since $[\varphi](\alpha)=\alpha$, invariance of the
intersection number gives
\[
i(\alpha,\lambda)
=
i([\varphi](\alpha),[\varphi](\lambda))
=
c\,i(\alpha,\lambda).
\]
Consequently, either $i(\alpha,\lambda)=0$ or $c=1$. It follows that
\[
\operatorname{Fix}([\omega])
\subset
\{[\lambda]\in\pml\mid i(\alpha,\lambda)=0\}
\cup
\pml\bigl(x_0,[\varphi]^{-1}(x_0)\bigr),
\]
where $x_0$ is chosen so that $[\varphi](x_0)\ne x_0$. Compare \cite[Corollary~2.17]{Ivanov1992}.

The first set is null because the set of filling measured laminations
has full Thurston measure. The second set is null by
\Cref{thm:level_set}. Hence $\operatorname{Fix}([\omega])$ is
$\PThursM^x$-null for every $x\in\teich_{g,m}$.
\end{proof}

\subsection{Proof of \Cref{thm:Hopf_decomposition_main}}
\label{subsec:ProofOfDecomposition}
We begin by recalling the relevant terminology from \cite{Kaimanovich2010}.
Let $\mathscr{C}$ and $\mathscr{D}$ denote the unions of the
ergodic components whose conditional probability measures are
purely nonatomic and purely atomic, respectively. Following Kaimanovich \cite{Kaimanovich2010}, we call $\mathscr{C}$ and $\mathscr{D}$ the \emph{continual} and \emph{discontinual} parts of the action of $G$.
Set
$\mathscr{D}_{\infty}=\{[\lambda]\in\mathscr{D}:\#\operatorname{Stab}_G([\lambda])=\infty\}$,
\[
\mathtt{cons}_{\infty}
=
\mathscr{C}\cup\mathscr{D}_{\infty}.
\]
Also, let
$A_0=
\{[\lambda]\in\pml:
\operatorname{Stab}_G([\lambda])\ne\{\mathrm{id}\}\}$.
By \cite[Theorem~14 and Corollary~15]{Kaimanovich2010},
$\mathtt{cons}_{\infty}$ is the infinitely conservative part,
and $\mathtt{cons}_{\infty}\cup A_0$ is the conservative part
in Kaimanovich's terminology, modulo null sets.

Fix $x\in\teich_{g,m}$, and suppose first that $G$ is infinite.
By \eqref{eq:Radon-Nikodym2} and
\cite[Theorem~29]{Kaimanovich2010},
$\conservativepart(G)$ agrees with
$\mathtt{cons}_{\infty}$ modulo null sets.
By \Cref{lem:Fix_null} and the countability of $G$, the set
$A_0$ is null. Hence the action on $\conservativepart(G)$
is conservative in our terminology, and the action on its
complement $\dissipativepart(G)$ is dissipative.
The assertion now follows from the Hopf decomposition for general nonsingular actions
\cite[Theorem~14 and Corollary~15]{Kaimanovich2010}.

If $G$ is finite, the defining series converges everywhere.
Thus $\conservativepart(G)=\varnothing$ and $\dissipativepart(G)=\pml$.
By essential freeness, we may restrict the action to a $G$-invariant conull Borel subset on which it is free. A free action of a finite group on a standard Borel space admits a Borel fundamental domain. Hence the action is totally dissipative.
\qed

\begin{remark}
The proof of \Cref{thm:Hopf_decomposition_main} gives the following description in terms of ergodic components of the action of a nontrivial subgroup $G$ of $\tmod(g,m)$:
\begin{align*}
(\text{the conservative part of the action of $G$})&=\mathscr{C} \\
(\text{the dissipative part of the action of $G$})&= \mathscr{D}_1=\mathscr{D}\setminus A_0
\end{align*}
modulo $\PThursM^x$-null sets. Indeed, since the action of $\tmod(g,m)$ is essentially free (cf. \Cref{lem:Fix_null}), $\mathscr{D}_\infty$ and $A_0$ are null. Compare \cite[Corollary 16]{Kaimanovich2010}.
\end{remark}

\subsection{The conservative part}
\label{subsec:Proof_of_ConservativePart_BigHorospherical}
We prove the first assertion of \Cref{thm:main_classfication_limit_sets}.

If $G$ is finite, then both $\conservativepart(G)$ and
$\Lambda_H(G)$ are empty. We may therefore assume that $G$ is infinite.

Fix $x\in\teich_{g,m}$. 
By \cite[Theorem~29]{Kaimanovich2010}, the infinitely conservative part $\mathtt{cons}_\infty$ agrees modulo null sets with
\[
\left\{
[\lambda]\in\pml
\mathrel{}\middle|\mathrel{}
\exists a>0\text{ s.t. }
\#\left\{[\omega]\in G:\frac{d([\omega]_*\PThursM^x)}
     {d\PThursM^x}([\lambda])>a\right\}=\infty
\right\}.
\]
For $[\lambda]\in\pml^{ue}$, equations \eqref{eq:Radon-Nikodym2} and \eqref{eq:busemann_ue} show that this condition is equivalent to the existence of $b\in\mathbb{R}$ such that
\[
\#\left\{
[\omega]\in G
\mathrel{}\middle|\mathrel{}
\busemann^x_\lambda([\omega](x))<b
\right\}
=\infty.
\]
Since $\operatorname{Stab}_G(x)$ is finite, this is equivalent to the
condition that some horoball centered at $[\lambda]$ contains infinitely
many points of the orbit $G(x)$. Thus it is precisely the condition
$[\lambda]\in\Lambda_H(G)$.

The set $\pml^{ue}$ has full $\PThursM^x$-measure. Moreover, by
\Cref{lem:Fix_null}, the action of $G$ on $\pml$ is essentially free.
Consequently, its conservative and infinitely conservative parts agree
modulo null sets. It follows that
\[
\conservativepart(G)=\Lambda_H(G)
\]
modulo $\PThursM^x$-null sets.
\qed

\subsection{Ideal boundary of the Dirichlet polyhedron}
\label{subsec:DirichletDomain_Ideal}
In \cite{McCarthyPapadopoulos1996}, McCarthy and Papadopoulos construct fundamental domains for subgroups of the mapping class group that are analogous to Dirichlet polyhedra for hyperbolic manifolds (see, e.g., \cite{Beardon1995}).

Let $G\le \tmod(g,m)$, and let $x_0\in\teich_{g,m}$ satisfy $\operatorname{Stab}_G(x_0)=\{\mathrm{id}\}$. The \emph{Dirichlet polyhedron} $\Dirichlet_G(x_0)$ for $G$, centered at $x_0$, is defined by
\[
\Dirichlet_G(x_0)
=\left\{x\in\teich_{g,m}\mathrel{}\middle|\mathrel{}
d_T(x,x_0)<d_T(x,[\omega](x_0))
\text{ for every }[\omega]\in G\setminus\{\mathrm{id}\}\right\}.
\]
McCarthy and Papadopoulos prove the following:
\begin{itemize}
\item The closure in $\teich_{g,m}$ of the Dirichlet polyhedron satisfies
\[
\overline{\Dirichlet_G(x_0)}
=\left\{x\in\teich_{g,m}\mathrel{}\middle|\mathrel{}
d_T(x,x_0)\le d_T(x,[\omega](x_0))
\text{ for every }[\omega]\in G\right\},
\]
and $x_0$ is an interior point of this closure (\cite[Propositions 5.1 and 5.2]{McCarthyPapadopoulos1996}).
\item The Dirichlet polyhedron $\Dirichlet_G(x_0)$ and its closure are starlike with respect to $x_0$ (\cite[Proposition 5.3]{McCarthyPapadopoulos1996}).
\item As in the case of hyperbolic spaces (see \cite[Chapter 9]{Beardon1995}), $\Dirichlet_G(x_0)$ is a locally finite fundamental domain in $\teich_{g,m}$ (\cite[Theorem 5.9]{McCarthyPapadopoulos1996}).
\item The boundary $\partial \Dirichlet_G(x_0)$ is a properly embedded, locally flat submanifold of codimension one in $\teich_{g,m}$, and every boundary point is contained in at least one hypersurface of the form
\[
L_{[\omega]}(x_0)
=\{x\in \teich_{g,m}\mid d_T(x,x_0)=d_T(x,[\omega](x_0))\}
\]
for some $[\omega]\in G\setminus\{\mathrm{id}\}$ (\cite[Propositions 5.2 and 5.5]{McCarthyPapadopoulos1996}).
\end{itemize}
Let $\Dirichlet^\infty_G(x_0)$ denote the intersection of the closure of $\Dirichlet_G(x_0)$ in the Thurston compactification with $\pml$.

\begin{lemma}
\label{lem:Dirichlet_geodesic_horoball}
Let $[\lambda]\in\pml^{ue}$, and consider the following conditions:
\begin{itemize}
\item[(a)] The Teichm\"uller ray $\ray^{\lambda}_{x_0}$ is contained in $\Dirichlet_G(x_0)$.
\item[(b)] $[\lambda]\in \Dirichlet^\infty_G(x_0)$.
\item[(c)] The $0$-horoball $\horos([\lambda];x_0)$ is disjoint from the $G$-orbit $G(x_0)$ of $x_0$.
\end{itemize}
Condition ${\rm (a)}$ implies ${\rm (c)}$, and conditions ${\rm (a)}$ and ${\rm (b)}$ are equivalent. Moreover, if $\overline{\horos([\lambda];x_0)}\cap G(x_0)=\{x_0\}$, then ${\rm (a)}$ holds.
\end{lemma}

\begin{proof}
\noindent
\textbf{${\rm (a)}\Rightarrow {\rm (c)}$.}
Let $[\omega]\in G$. Since $\ray^\lambda_{x_0}(t)\in \Dirichlet_G(x_0)$ for every $t\ge 0$,
\[
d_T(\ray^\lambda_{x_0}(t),[\omega](x_0))
\ge d_T(\ray^\lambda_{x_0}(t),x_0)=t,
\]
and hence
\[
\busemann^{x_0}_\lambda([\omega](x_0))
=\lim_{t\to\infty}
\bigl(d_T(\ray^\lambda_{x_0}(t),[\omega](x_0))-t\bigr)
\ge 0.
\]
Thus, $[\omega](x_0)$ lies outside the $0$-horoball.

We next prove the additional assertion. Assume that $x_0$ is the only point of the orbit contained in the closure of the $0$-horoball, and let $[\omega]\in G\setminus\{\mathrm{id}\}$. By assumption,
$\busemann^{x_0}_\lambda([\omega](x_0))>0$. The function
$t\mapsto d_T([\omega](x_0),\ray^\lambda_{x_0}(t))-t$ is nonincreasing and converges to $\busemann^{x_0}_\lambda([\omega](x_0))$. Consequently, for every $t\ge0$,
\[
d_T([\omega](x_0),\ray^\lambda_{x_0}(t))-t
\ge\busemann^{x_0}_\lambda([\omega](x_0))>0.
\]
Since this holds for every $[\omega]\in G\setminus\{\mathrm{id}\}$, we have $\ray^\lambda_{x_0}(t)\in\Dirichlet_G(x_0)$ for every $t\ge0$. Thus, ${\rm (a)}$ holds.

\medskip
\noindent
\textbf{${\rm (b)}\Rightarrow {\rm (a)}$.}
Since $[\lambda]\in \Dirichlet^\infty_G(x_0)$, there exists a sequence $\{x_n\}_{n=1}^\infty\subset \Dirichlet_G(x_0)$ such that $x_n\to[\lambda]$ in the Thurston compactification. Set $t_n=d_T(x_0,x_n)$. Then $x_n$ also converges to $[\lambda]$ in the Gardiner--Masur compactification, and
\[
e^{-2t_n}\ext_{x_n}(\cdot)
\longrightarrow \frac{i(\cdot,\lambda)^2}{\ext_{x_0}(\lambda)}
\]
uniformly on compact subsets of $\ml$ (cf.\ \cite[Corollaries 1 and 2]{Miyachi2013}).

For each $n\in\mathbb{N}$, choose $[\lambda_n]\in\pml$ such that
$\ray^{\lambda_n}_{x_0}(t_n)=x_n$. Choose representatives $\lambda_n$ normalized by $\ext_{x_0}(\lambda_n)=1$. Then $t_n\to\infty$, and, after passing to a subsequence, we may assume that $\lambda_n$ converges to some $\lambda_0\in\ml$. Moreover,
\[
e^{-2t_n}\ext_{x_n}(\lambda_n)
=e^{-4t_n}\ext_{x_0}(\lambda_n)\longrightarrow 0.
\]
Hence $i(\lambda_0,\lambda)=0$, and thus $[\lambda_0]=[\lambda]$. By Teichm\"uller's theorem, the geodesic rays $\ray^{\lambda_n}_{x_0}$ converge uniformly on compact subsets of $[0,\infty)$ to $\ray^\lambda_{x_0}$.

Suppose, for contradiction, that the ray is not contained in $\Dirichlet_G(x_0)$. Since $x_0\in\Dirichlet_G(x_0)$ and the ray is continuous, there exists $t_0>0$ such that $\ray^\lambda_{x_0}(t_0)\in\partial\Dirichlet_G(x_0)$. Then there exists $[\omega]\in G\setminus\{\mathrm{id}\}$ such that
\[
d_T(\ray^\lambda_{x_0}(t_0),x_0)
=d_T(\ray^\lambda_{x_0}(t_0),[\omega](x_0)).
\]
By \cite[Lemma 4.5]{McCarthyPapadopoulos1996},
\[
d_T(x_0,\ray^\lambda_{x_0}(t_0+1))
>d_T([\omega](x_0),\ray^\lambda_{x_0}(t_0+1)).
\]
Therefore, for all sufficiently large $n$, we have $t_0+1<t_n$ and
$\ray^{\lambda_n}_{x_0}(t_0+1)\notin\Dirichlet_G(x_0)$. This contradicts the starlikeness of $\Dirichlet_G(x_0)$ with respect to $x_0$, because $x_n\in\Dirichlet_G(x_0)$.

\medskip
\noindent
\textbf{${\rm (a)}\Rightarrow {\rm (b)}$.}
Since $\lambda$ is uniquely ergodic, $\ray^\lambda_{x_0}(t)\to[\lambda]$ in the Thurston compactification (cf.\ \cite{Masur1980}). By ${\rm (a)}$, every point of this ray lies in $\Dirichlet_G(x_0)$; hence $[\lambda]\in\Dirichlet^\infty_G(x_0)$.
\end{proof}

\begin{lemma}[Dirichlet points]
\label{lem:almost-unique}
Assume that $[\omega](x_0)\ne x_0$ for every $[\omega]\in G\setminus\{\mathrm{id}\}$.
Then $\Dirichlet^\infty_G(x_0)$ is a wandering set for the action of $G$. Moreover, for $\PThursM^{x_0}$-almost every $[\lambda]\in\Dirichlet^\infty_G(x_0)$, we have $\overline{\horos([\lambda];x_0)}\cap G(x_0)=\{x_0\}$.
\end{lemma}

\begin{proof}
Fix $[\omega]\in G\setminus\{\mathrm{id}\}$ and let $[\lambda]\in \Dirichlet^\infty_G(x_0)\cap [\omega](\Dirichlet^\infty_G(x_0))\cap \pml^{ue}$.
Since $[\omega](\Dirichlet^\infty_G(x_0))=
\Dirichlet^\infty_G([\omega](x_0))$, \Cref{lem:Dirichlet_geodesic_horoball} gives
\[
\busemann^{x_0}_\lambda([\omega](x_0))\ge0,
\qquad
\busemann^{[\omega](x_0)}_\lambda(x_0)\ge0.
\]
From the cocycle condition \eqref{eq:busemann_cocycle_condition},
\[
0=\busemann^{x_0}_\lambda(x_0)=\busemann^{x_0}_\lambda([\omega](x_0))+\busemann^{[\omega](x_0)}_\lambda(x_0),
\]
we have $\busemann^{x_0}_\lambda(x_0)=\busemann^{x_0}_\lambda([\omega](x_0))=0$.
Therefore,
\begin{align*}
\frac{\ext_{[\omega](x_0)}(\lambda)}{\ext_{x_0}(\lambda)}
&=\exp\left(2\busemann^{x_0}_\lambda([\omega](x_0))\right)=\exp\left(2\busemann^{x_0}_\lambda(x_0)\right)=1.
\end{align*}
Thus,
\[
\Dirichlet^\infty_G(x_0)\cap
[\omega](\Dirichlet^\infty_G(x_0))\cap\pml^{ue}
\subset\pml(x_0,[\omega](x_0)).
\]
The right-hand side is null by \Cref{thm:level_set}, and
$\pml\setminus\pml^{ue}$ is null. Hence
$\Dirichlet^\infty_G(x_0)$ is wandering.

Let $[\lambda]\in\Dirichlet^\infty_G(x_0)\cap\pml^{ue}$.
Suppose that $[\omega](x_0)\in\overline{\horos([\lambda];x_0)}$ for some $[\omega]\in G\setminus\{\mathrm{id}\}$. 
By \Cref{lem:Dirichlet_geodesic_horoball},
$\busemann^{x_0}_\lambda([\omega](x_0))\ge0$, whereas membership in the
closed horoball gives the opposite inequality. Hence this value is zero,
and \eqref{eq:busemann_ue} yields
$[\lambda]\in\pml(x_0,[\omega](x_0))$. 
Since $G$ is countable, 
we have $\overline{\horos([\lambda];x_0)}\cap G(x_0)=\{x_0\}$ for $\PThursM^{x_0}$-almost every $[\lambda]\in \Dirichlet^{\infty}_G(x_0)$.
\end{proof}

\begin{lemma}[Dirichlet sets and Dirichlet polyhedra]
\label{lem:DirichletSets_Polyhedrons}
Let $G$ be a nontrivial subgroup of the Teichm\"uller modular group, and let
$x_0\in \teich_{g,m}$ satisfy $\operatorname{Stab}_G(x_0)=\{\mathrm{id}\}$.
Then the symmetric difference
\begin{equation}
\label{eq:symmetric_differenice_Dirichletset_Dirichlet_polyhedron}
\operatorname{Dir}_G(x_0)\bigtriangleup \bigcup_{[\omega]\in G}[\omega](\Dirichlet^\infty_G(x_0))
\end{equation}
is $\PThursM^{x_0}$-null.
\end{lemma}

\begin{proof}
Let $E\subset\Dirichlet^\infty_G(x_0)$ consist of the uniquely ergodic points $[\lambda]$ satisfying
\[
\overline{\horos([\lambda];x_0)}\cap G(x_0)=\{x_0\}.
\]
From the proof of \Cref{lem:almost-unique}, $[\omega](E)\cap E=\emptyset$ for every $[\omega]\in G\setminus\{\mathrm{id}\}$ and $\Dirichlet^\infty_G(x_0)\setminus E$ is $\PThursM^{x_0}$-null.
By definition, $E\subset \operatorname{Dir}_G(x_0)$. Moreover,
\[
\operatorname{Dir}_G(x_0)
\setminus\bigcup_{[\omega]\in G}[\omega](E)
\subset
(\pml\setminus\pml^{ue})
\cup
\bigcup_{\substack{[\omega_1],[\omega_2]\in G\\
[\omega_1]\ne[\omega_2]}}
\pml([\omega_1](x_0),[\omega_2](x_0)).
\]
Indeed, outside the exceptional set on the right, a Dirichlet
point is uniquely ergodic and has a unique minimizing orbit
point. Applying \Cref{lem:Dirichlet_geodesic_horoball} at
that orbit point and translating back to $x_0$, we see that
it belongs to $[\omega](E)$ for some $[\omega]\in G$.
The exceptional set is null by \Cref{thm:level_set} and countability
of $G$; the triviality of $\operatorname{Stab}_G(x_0)$ ensures that
$[\omega_1](x_0)\ne[\omega_2](x_0)$ whenever
$[\omega_1]\ne[\omega_2]$.
Conversely, every point of $[\omega](E)$ is uniquely ergodic and has
$[\omega](x_0)$ as its unique minimizing orbit point, by equivariance
of extremal length. Thus
\[
\bigcup_{[\omega]\in G}[\omega](E)\subset\operatorname{Dir}_G(x_0).
\]
Finally, replacing $E$ by $\Dirichlet^\infty_G(x_0)$ changes this
union only by a null set, because the action is nonsingular and $G$
is countable. This proves the assertion.
\end{proof}

\subsection{Proof of \Cref{thm:Dissipative_part}}
\label{subsec:Proof-Dissipative_part}
Let $\dissipativepart^0_{x_0}(G)$ be the subset of $\dissipativepart_{x_0}(G)\cap\pml^{ue}$ consisting of the points for which the infimum defining $\horo^{x_0}_G([\lambda])$ is attained at a unique point of the orbit $G(x_0)$.

We first show that $\dissipativepart^0_{x_0}(G)$ is conull in $\dissipativepart_{x_0}(G)$. Indeed, if $[\lambda]\in\dissipativepart_{x_0}(G)\cap\pml^{ue}$, then the defining series converges. A summable family of positive terms attains its supremum; otherwise, infinitely many terms would be bounded below by half of that supremum. It follows that $\horo^{x_0}_G([\lambda])>0$ and that the infimum defining it is attained. If the minimum is attained at two distinct orbit points $y,z\in G(x_0)$, then $[\lambda]\in\pml(y,z)$. Since the orbit is countable, \Cref{thm:level_set}, together with the fact that $\pml^{ue}$ has full measure, shows that the set of exceptions is null.

Let $[\lambda]\in\dissipativepart^0_{x_0}(G)$. We choose $[\omega]\in G$ such that the infimum defining $\horo^{x_0}_G([\lambda])$ is attained at $[\omega](x_0)$.
Set $r_0=\busemann^{x_0}_\lambda([\omega](x_0))=\frac{1}{2}\log\horo^{x_0}_G([\lambda])$. By the uniqueness of the minimizing orbit point and \eqref{eq:horoball-basepoint2},
\[
G(x_0)\cap\overline{\horos([\lambda];[\omega](x_0))}
=G(x_0)\cap\overline{\horos([\lambda],r_0;x_0)}
=\{[\omega](x_0)\}.
\]
Applying \Cref{lem:Dirichlet_geodesic_horoball} with basepoint $[\omega](x_0)$, we obtain $[\lambda]\in\Dirichlet^\infty_G([\omega](x_0))$. Since
$\Dirichlet^\infty_G([\omega](x_0))=[\omega](\Dirichlet^\infty_G(x_0))$, it follows that
\[
\dissipativepart_{x_0}(G)
\setminus\bigcup_{[\omega]\in G}
[\omega](\Dirichlet^\infty_G(x_0))
\]
is $\PThursM^{x_0}$-null.

For the converse inclusion, let $E\subset\Dirichlet^\infty_G(x_0)$ consist of the uniquely ergodic points $[\lambda]$ satisfying
$\overline{\horos([\lambda];x_0)}\cap G(x_0)=\{x_0\}$. 
From the proof of \Cref{lem:almost-unique}, $[\omega](E)\cap E=\emptyset$ for every $[\omega]\in G\setminus\{\mathrm{id}\}$ and $\Dirichlet^\infty_G(x_0)\setminus E$ is $\PThursM^{x_0}$-null.

The set $E$ is measurable: on $\pml^{ue}$ its defining condition is
\[
\ext_{[\omega](x_0)}(\lambda)>\ext_{x_0}(\lambda)
\quad\text{for every }[\omega]\in G\setminus\{\mathrm{id}\},
\]
a countable family of measurable conditions.
We now show directly that $E\setminus\dissipativepart_{x_0}(G)$ is null.
By the construction of $E$, the sets $[\omega]^{-1}(E)$, $[\omega]\in G$, are pairwise disjoint.
Hence, by Tonelli's theorem and \eqref{eq:Radon-Nikodym2},
\begin{align*}
\int_E
\sum_{[\omega]\in G}
\left(
\frac{\ext_{x_0}(\lambda)}
{\ext_{[\omega](x_0)}(\lambda)}
\right)^{\convgenus}
d\PThursM^{x_0}([\lambda])
&=
\sum_{[\omega]\in G}
\int_E
\frac{d([\omega]_*\PThursM^{x_0})}{d\PThursM^{x_0}}([\lambda])
d\PThursM^{x_0}([\lambda])\\
&=
\sum_{[\omega]\in G}
\PThursM^{x_0}([\omega]^{-1}(E)) \\
&\le \PThursM^{x_0}(\pml)=1<\infty.
\end{align*}
Therefore, the defining series for $\dissipativepart_{x_0}(G)$ is finite for $\PThursM^{x_0}$-almost every point of $E$. 
Thus, $E\setminus\dissipativepart_{x_0}(G)$ is null. Together with \Cref{lem:almost-unique}, this shows that $\Dirichlet^\infty_G(x_0)\setminus\dissipativepart_{x_0}(G)$ is null. By \Cref{lem:Cons-Diss-ElementaryProperty},
$\dissipativepart_{x_0}(G)$ is $G$-invariant.
Since the action is nonsingular and $G$ is countable, it follows that
\[
\left(
\bigcup_{[\omega]\in G}
[\omega](\Dirichlet^\infty_G(x_0))
\right)
\setminus\dissipativepart_{x_0}(G)
\]
is also $\PThursM^{x_0}$-null. The two inclusions prove the assertion.
The nullity of the symmetric difference $\dissipativepart(G)\bigtriangleup \operatorname{Dir}_G(x_0)$ follows from \Cref{lem:DirichletSets_Polyhedrons}.

The result for $\PThursM^x$, with arbitrary
$x\in\teich_{g,m}$, follows from the mutual absolute continuity of the
measures $\{\PThursM^x\}_{x\in\teich_{g,m}}$.
This also proves the second assertion of
\Cref{thm:main_classfication_limit_sets}.
\qed

\section{Dynamics of the Torelli group}
\label{sec:dynamics_of_Torelli_group}

\subsection{Siegel upper half space}
We recall the Siegel upper half space, its realization as the Siegel disk, and the relation between their boundaries under the Cayley transform, following Friedland and Freitas \cite{FriedlandFreitas2004}.
This point of view is classical in the theory of Hermitian symmetric spaces; see, for example, Helgason \cite{Helgason2001} and Satake \cite{Satake1980}.

Let $\mat_{m,n}(\mathbb{C})$ denote the set of $m\times n$ complex matrices, and let
$\symm(g,\mathbb{C})\subset \mat_{g,g}(\mathbb{C})$ denote the set of symmetric matrices.
The \emph{Siegel upper half space} $\mathfrak{H}_g$ of degree $g$ is the open subset of
$\symm(g,\mathbb{C})$ consisting of symmetric matrices whose imaginary parts are positive definite.
The symplectic group
\[
\spl(2g,\mathbb{R})=\left\{
M\in \gl(2g,\mathbb{R})\mid M^TJ_gM=J_g\right\}
\quad
\left(
J_g=\begin{pmatrix} 0 & I_g \\ -I_g & 0\end{pmatrix}
\right)
\]
acts on $\mathfrak{H}_g$ by
\[
M(Z)=(AZ-B)(-CZ+D)^{-1}
\quad
\left(M=\begin{pmatrix} A & B \\ C & D\end{pmatrix}\right).
\]
This convention is compatible with the action of the symplectic group on the first homology group of $\Sigma_g$
(cf. \S\ref{subsec:symplectic_representation} and \eqref{eq:symplectic_rep}; see also \cite[(6.8)]{Earle1978}).

% The boundary $\partial\mathfrak{H}_g$ of the Siegel upper half spaceconsists of symmetric matrices $Z$ with
% ${\rm Im}(Z)\ge 0$ and $\rank({\rm Im}(Z))<g$.
% For $k=1,\ldots,g$, let
% \[
% \partial^{(k)}\mathfrak{H}_g
% =
% \{Z\in \partial\mathfrak{H}_g\mid \rank({\rm Im}(Z))=g-k\}.
% \]
% Then the rank decomposition $\{\partial^{(k)}\mathfrak{H}_g\}_{k=1}^{g}$ gives a stratification of
% the boundary $\partial\mathfrak{H}_g$.
% Moreover, this stratification is invariant under the action of $\spl(2g,\mathbb{R})$
% in the following sense: if
% $Z\in \partial^{(k)}\mathfrak{H}_g$ and
% $M=\begin{pmatrix} A & B \\ C & D\end{pmatrix}\in \spl(2g,\mathbb{R})$ satisfy
% $\det(-CZ+D)\ne 0$, then $M(Z)\in \partial^{(k)}\mathfrak{H}_g$.

Let
\[
\mathfrak{D}_g=\{W\in \symm(g,\mathbb{C})\mid I_g-W\overline{W}>0\}.
\]
This domain is called the \emph{Siegel disk}, or the \emph{generalized unit disk}, of degree $g$
(cf. \cite{Siegel1943}).
It is known that the \emph{Cayley transform}
\begin{equation}
\label{eq:cayray}
% \Cayley\colon\mathfrak{D}_g\ni W\mapsto \sqrt{-1}(I_g+W)(I_g-W)^{-1}\in \mathfrak{H}_g
\Cayley\colon\mathfrak{H}_g\ni Z\mapsto (Z-\sqrt{-1}I_g)(Z+\sqrt{-1}I_g)^{-1}\in \mathfrak{D}_g
\end{equation}
is biholomorphic. Via the Cayley transform, $\spl(2g,\mathbb{R})$ acts naturally on
$\mathfrak{D}_g$, and this action extends to the boundary $\partial \mathfrak{D}_g$ in
$\symm(g,\mathbb{C})$.
We use the following characterization (cf. Friedland and Freitas \cite{FriedlandFreitas2004}).

\begin{proposition}
\label{prop:F-F-lemma3.1}
A sequence $\{Z_n\}_{n=1}^\infty$ in $\mathfrak{H}_g$ converges to a point in
$\partial \mathfrak{H}_g\subset \symm(g,\mathbb{C})$ if and only if the sequence
$\{\Cayley(Z_n)\}_{n=1}^\infty$ converges to a point
$W_\infty\in \partial \mathfrak{D}_g$ with $\det(I_g-W_\infty)\ne 0$.
\end{proposition}

\subsection{Symplectic representation, Torelli group and Torelli space}
\label{subsec:symplectic_representation}
The action of the mapping class group on $H_1(\Sigma_g,\mathbb{Z})$ gives a homomorphism
\[
\hat{\symrep}\colon \mcg(\Sigma_g)\to
{\rm Aut}(H_1(\Sigma_g,\mathbb Z)),
\]
called the \emph{symplectic representation} of the mapping class group (cf. \cite[Chapter Six]{FarbMargalit2012}). The kernel $\torelliG_g<\mcg(\Sigma_g)$ of the symplectic representation is called the \emph{Torelli group}. We now fix a canonical symplectic homology basis $\{\mathbf{a},\mathbf{b}\}=\{\mathbf{a}_1,\cdots,\mathbf{a}_g,\mathbf{b}_1,\cdots, \mathbf{b}_g\}$ on $\Sigma_g$. With respect to this basis, the symplectic representation takes the form
\begin{equation}
\label{eq:symp_homo1}
\hat{\symrep}\colon \mcg(\Sigma_{g})\to \syp(2g,\mathbb{Z})
\end{equation}
with values in the integral symplectic group. The induced action on $H_1(\Sigma_g,\mathbb{Z})=\bigoplus_{j=1}^g(\mathbb{Z}\mathbf{a}_j\oplus \mathbb{Z}\mathbf{b}_j)$ is given by
 \[
 (\omega(\mathbf{a}),\omega(\mathbf{b}))=(\mathbf{a},\mathbf{b}) \hat{\symrep}([\omega]_*)
 \]
 for $[\omega]_*\in \mcg(\Sigma_{g})$. Composing \eqref{eq:symp_homo1} with the natural projection $\syp(2g,\mathbb Z)\to\psyp(2g,\mathbb Z)$ gives a homomorphism
\begin{equation}
\label{eq:symp_homo2}
\symrep\colon \mcg(\Sigma_{g})\to \psyp(2g,\mathbb{Z})\subset \operatorname{Aut}(\mathfrak{H}_g).
\end{equation}

The quotient space
\[
\torelli_{g}=\teich_{g}/\torelliG_{g}
\]
is called the \emph{Torelli space} of genus $g$. It is known that the Torelli group $\torelliG_{g}$ is torsion-free and acts freely on $\teich_{g}$. Hence $\torelli_{g}$ is a complex manifold of dimension $3g-3$ and the natural projection $\teich_{g} \to \torelli_{g}$ is a holomorphic universal covering map (cf. \cite[Theorem 6.12]{FarbMargalit2012}).
The Torelli space $\torelli_{g}$ can also be viewed as the deformation space of framed Riemann surfaces of genus $g$, where a \emph{framed Riemann surface} of genus $g$ is a closed Riemann surface $M$ of genus $g$ together with a symplectic basis of $H_1(M,\mathbb{Z})$ (e.g. \cite[\S2]{Hain2006}).

\subsection{Period map}
\label{subsec:period_map}
We use the symplectic basis fixed in \S\ref{subsec:symplectic_representation}. 
 For $x=(M,f)\in \teich_{g}$,
 we denote by $\omega_x^j$ the holomorphic $1$-form on $M$ satisfying
\[
\int_{f_*(\mathbf{a}_k)}\omega_x^j=\delta_{jk},
\]
where $\delta_{jk}$ denotes the Kronecker delta. We define the \emph{period matrix} $\Pi(x)=\begin{pmatrix} \Pi_{jk}(x)\end{pmatrix}$ of $x$ with respect to this basis by
\[
\Pi_{jk}(x)=\int_{f_*(\mathbf{b}_k)}\omega_x^j
\]
(cf. \cite{FarkasKra1992}).
It is known that the \emph{period map} $\Pi \colon \teich_{g}\to \mathfrak{H}_g$ is holomorphic (\cite{Ahlfors1960}).
The period map satisfies the equivariance relation
\begin{equation}
\label{eq:symplectic_rep}
\symrep([\omega]_*)(\Pi(x))=\Pi\circ [\omega]_*(x)
\end{equation}
for $x\in \teich_{g}$ and $[\omega]\in \tmod(g)$.

The period map descends to a holomorphic map
$\Pi'\colon \torelli_g\to\mathfrak H_g$.
For $g=2$, this map is an embedding.
For $g\ge3$, it is generically two-to-one onto its image,
with branching along the hyperelliptic locus.
Set
\[
\mathcal{J}_g=\Pi(\teich_g),
\qquad
\mathcal{J}_g^c
=\overline{\mathcal{J}_g}^{\,\mathfrak H_g}.
\]
Let $\torelli_g^c$ denote the Torelli space of stable
curves of compact type, namely stable curves whose
dual graphs are trees, equipped with a symplectic
homology basis.
The period map extends to a proper holomorphic map
\[
\Pi^c\colon\torelli_g^c\longrightarrow \mathcal{J}_g^c
\]
whose restriction to $\torelli_g$ is $\Pi'$
(cf.\ \cite[\S\S2 and 4]{Hain2006}).

The Jacobian of a curve of compact type is the
product, with its product polarization, of the
Jacobians of its irreducible components.
Thus $\mathcal{J}_g^c\setminus \mathcal{J}_g$ consists of the period
matrices of singular curves of compact type;
these principally polarized abelian varieties
are decomposable.
Here and below, decomposability refers to the
principal polarization.
We obtain the following consequence.

\begin{proposition}
\label{prop:compacttype}
Let $\{y_n\}$ be a sequence in $\torelli_g$ that
leaves every compact subset.
If $\Pi'(y_n)$ converges to $Z\in\mathfrak H_g$,
then the principally polarized abelian variety
represented by $Z$ is decomposable.

Conversely, let $\{x_n\}$ be a sequence in $\teich_g$
such that $\Pi(x_n)$ converges to the period matrix
of a decomposable principally polarized abelian
variety.
Then the images of $x_n$ in the moduli space
$\mathcal M_g$ leave every compact subset.
\end{proposition}

\begin{proof}
For the first assertion, set
\[
K=\{Z\}\cup\{\Pi'(y_n):n\ge1\}.
\]
This is a compact subset of $\mathcal{J}_g^c$.
By properness of $\Pi^c$, after passing to a
subsequence, $y_n$ converges to a point
$y_\infty\in\torelli_g^c$.
Since $y_n$ leaves every compact subset of
$\torelli_g$, the limit lies in
$\torelli_g^c\setminus\torelli_g$.
Its Jacobian is therefore decomposable, and
$\Pi^c(y_\infty)=Z$.

For the converse, let
\[
\mathcal A_g
=
\operatorname{Sp}(2g,\mathbb Z)\backslash\mathfrak H_g,
\]
and let $j\colon\mathcal M_g\to\mathcal A_g$
be the map sending a smooth curve to its
principally polarized Jacobian.
Suppose that the images of $x_n$ in $\mathcal M_g$
do not leave every compact subset.
After passing to a subsequence, these images
converge to a smooth curve $C\in\mathcal M_g$.
By continuity of $j$, their Jacobians converge
to $j(C)$.
They also converge to the class of the assumed
decomposable limit.
Since $\mathcal A_g$ is Hausdorff, these limits
coincide.
This contradicts the indecomposability of the
principally polarized Jacobian of a smooth curve.
\end{proof}

\subsection{Radial limit of the period map}
We prove the following statement.
\begin{lemma}[Radial limit of the period map]
    \label{lem:radial_limit}
There exists a measurable set $\mathcal{E}_1\subset \pml$ of full measure with the following properties:
\begin{itemize}
    \item[(a)] $\mathcal{E}_1$ consists of uniquely ergodic measured laminations and is invariant under the action of the mapping class group
    $\tmod(g)$;
    \item[(b)] there exists a measurable map, called the radial limit of the period map,
    \[
    \Pi^*\colon \mathcal{E}_1\to \overline{\mathfrak{H}_g}
    \subset \symm(g,\mathbb{C})
    \]
    such that, for any $[\lambda]\in \mathcal{E}_1$ and any $x\in \teich_{g}$,
    \[
    \lim_{t\to \infty}\Pi(\ray^{\lambda}_x(t))=\Pi^*([\lambda]).
    \]
    In particular, this limit is independent of the choice of $x$;
    \item[(c)] for every $[\lambda]\in \mathcal{E}_1$ and every
    $\begin{bmatrix} A & B \\ C & D\end{bmatrix}\in \syp(2g,\mathbb{Z})$, the radial limit satisfies
    \[
    \det(-C\Pi^*([\lambda])+D)\ne 0;
    \]
    \item[(d)] for every $1\le i,j\le g$ and $[\lambda]\in \mathcal{E}_1$, the $(i,j)$-entry $\Pi^*_{ij}([\lambda])$ of the limit $\Pi^*([\lambda])$ is non-zero; and
    \item[(e)] for every $[\lambda]\in \mathcal{E}_1$ and every $x\in \teich_{g}$, the Teichm\"uller geodesic ray $\ray^{\lambda}_x$ is recurrent in the sense that its projection to the moduli space $\teich_{g}/\tmod(g)$ returns infinitely often to a compact subset of the moduli space.
\end{itemize}
\end{lemma}

\begin{remark}
By \Cref{lem:null_siegel} below, the radial limit $\Pi^*([\lambda])$ lies in the boundary $\partial\mathfrak{H}_g$ for almost every $[\lambda]\in \mathcal{E}_1$.
\end{remark}

\begin{proof}[Proof of \Cref{lem:radial_limit}]
By \cite[Theorem 3.1]{Miyachi2024Bounded},
there exist a $\tmod(g)$-invariant full-measure set
$\mathcal{E}_0\subset \mathcal{PML}^{ue}$ and a measurable map
\[
(\Cayley\circ\Pi)^*\colon
\mathcal E_0\to\overline{\mathfrak D_g}
\]
such that the following hold:
\begin{enumerate}
\item[(1)] for any $x\in \teich_{g}$ and $[\lambda]\in \mathcal{E}_0$,
\[
\lim_{t\to \infty}(\Cayley\circ\Pi)(\ray^{\lambda}_x(t))
=
(\Cayley\circ\Pi)^*([\lambda]);
\]
\item[(2)] for any $[\omega]_*\in \tmod(g)$,
\[
\symrep([\omega]_*)\bigl((\Cayley\circ\Pi)^*([\lambda])\bigr)
=
(\Cayley\circ\Pi)^*([\omega]_*[\lambda])
\]
for all $[\lambda]\in \mathcal{E}_0$.
\end{enumerate}
The set $\mathcal E_0$ therefore satisfies condition (a).

For
$M=\begin{pmatrix} A & B \\ C & D\end{pmatrix}\in \syp(2g,\mathbb{Z})$, define
\[
f_M(x)
=
\det\bigl(-\sqrt{-1}C+D-(\sqrt{-1}C+D)\Cayley(\Pi(x))\bigr).
\]
Notice that
\[
f_{I_{2g}}(x)=\det(I_g-\Cayley(\Pi(x))).
\]
The function $f_M$ is a bounded holomorphic function on $\teich_{g}$, since
$\Cayley(\Pi(x))\in \mathfrak{D}_g$.
Moreover, writing $W(x)=\Cayley(\Pi(x))$, we have
\[
f_M(x)
=
\det(-C\Pi(x)+D)\det(I_g-W(x))
\ne0
\]
for every $x\in\teich_g$.
Indeed, the first factor is nonzero because the
symplectic action is defined on $\mathfrak H_g$,
and the second is nonzero because
$W(x)\in\mathfrak D_g$.

By (1), its radial limit is given by
\[
f_M^*([\lambda])
=
\det\bigl(-\sqrt{-1}C+D-(\sqrt{-1}C+D)(\Cayley\circ \Pi)^*([\lambda])\bigr)
\]
for $[\lambda]\in \mathcal{E}_0$.
By the identity theorem for radial limit functions \cite[Theorem 1.2]{Miyachi2024Bounded}, the set
\[
\mathcal{E}_1^M
=
\{[\lambda]\in \mathcal{E}_0\mid f_M^*([\lambda])\ne 0\}
\]
has full measure in $\pml$.

Since $\syp(2g,\mathbb{Z})$ is countable, the intersection
\[
\mathcal{E}'_1
=
\bigcap_{M\in \syp(2g,\mathbb{Z})}\mathcal{E}_1^M
\]
also has full measure in $\pml$ and is invariant under the action of $\tmod(g)$.
In particular, since $f_{I_{2g}}^*([\lambda])\ne 0$ on $\mathcal{E}'_1$, we have
\[
\det\bigl(I_g-(\Cayley\circ\Pi)^*([\lambda])\bigr)\ne 0
\]
for all $[\lambda]\in \mathcal{E}'_1$.
For $[\lambda]\in\mathcal{E}_1'$, set
\[
W_\lambda=(\Cayley\circ\Pi)^*([\lambda]).
\]
Since $\det(I_g-W_\lambda)\ne0$, the map
\[
W\longmapsto
\sqrt{-1}(I_g+W)(I_g-W)^{-1}
\]
is well-defined and continuous in a neighborhood of $W_\lambda$.
By the inverse Cayley transform formula, for every
$x\in\teich_g$ we therefore have
\[
\lim_{t\to\infty}\Pi(\ray^\lambda_x(t))
=
\sqrt{-1}(I_g+W_\lambda)(I_g-W_\lambda)^{-1}.
\]
Define
\[
\Pi^*([\lambda])
=
\sqrt{-1}(I_g+W_\lambda)(I_g-W_\lambda)^{-1}.
\]
This matrix belongs to $\overline{\mathfrak H_g}$, since it is
a limit of matrices in $\mathfrak H_g$.
Moreover, $\Pi^*$ is measurable, and the limit is independent
of the basepoint $x$, because $W_\lambda$ has these properties.
Thus, $\mathcal E_1'$ satisfies condition (b).
Since the non-vanishing of $f_M^*([\lambda])$ is equivalent to
\[
\det(-C\Pi^*([\lambda])+D)\ne 0,
\]
$\mathcal{E}_1'$ satisfies condition (c).

% Before proving claim (d), we recall that, by \cite[\S5, Lemma]{MR766636}, for any holomorphic map
% $F\colon \mathbb{D}\to \teich_{g}$, the composition $\Pi\circ F$ admits a finite non-tangential limit almost everywhere on $\partial\mathbb{D}$.
We next arrange condition (d) by removing another null set.
Fix indices $1\le i,j\le g$ and a point $x\in \teich_g$.
Suppose, to the contrary, that there exists a measurable set
$A\subset \mathcal{E}'_1$ of positive $\PThursM^x$-measure such that
$\Pi^*_{ij}([\lambda])=0$ for all $[\lambda]\in A$.
By the Rauch variational formula, we have
$\partial \Pi_{ij}|_x=\omega^x_i\omega^x_j$ in $\mathcal{Q}_x=T^*_x\teich_g$
(cf. \cite{Ahlfors1960} and \cite{Rauch1965}).
In particular, $\Pi_{ij}$ is non-constant on $\teich_g$.
Since $\PThursM^x$ and $\PThursM^y$ are mutually absolutely continuous for any
$x,y\in \teich_g$, we may change the basepoint $x$, if necessary, and assume that
$\Pi_{ij}(x)\ne 0$.

By the disintegration method discussed in \cite[\S3.1, \S4]{Miyachi2024Bounded}, there exists a quadratic differential $q\in \mathcal{Q}_x$ such that
\[
A_q=\{e^{i\theta}\in \partial \mathbb{D}\mid [v(e^{-i\theta}q)]\in A\}
\]
has positive angular Lebesgue measure.
Let $\Phi_q\colon \mathbb{D}\to \teich_{g}$ be the Teichm\"uller disk determined by $q$, normalized by $\Phi_q(0)=x$.
Since $\Phi_q$ is holomorphic, by \cite[\S5, Lemma]{Shiga1984}, the composition $\Pi_{ij}\circ \Phi_q$ admits a finite non-tangential limit almost everywhere on $\partial\mathbb{D}$.
Moreover, by the definition of the radial limit $\Pi^*$, the non-tangential boundary value of $\Pi_{ij}\circ \Phi_q$ agrees with
the radial limit
\[
\lim_{r\to 1}\Pi_{ij}(\Phi_q(re^{i\theta}))
=\lim_{t\to\infty}\Pi_{ij}(\ray^{v(e^{-i\theta}q)}_x(t))
=\Pi_{ij}^*([v(e^{-i\theta}q)])=0
\]
for almost every $e^{i\theta}\in A_q$.
Namely, the non-tangential boundary value of $\Pi_{ij}\circ \Phi_q$ vanishes on a subset of $\partial\mathbb{D}$ of positive measure.
Therefore, by the Lusin--Privalov uniqueness theorem, we obtain
$\Pi_{ij}\circ \Phi_q\equiv 0$ on $\mathbb{D}$
(cf. \cite{LusinPriwaloff1925}; see also \cite{Tsuji1944}).
In particular, $\Pi_{ij}(x)=\Pi_{ij}(\Phi_q(0))=0$, which contradicts the choice of $x$.

Therefore, for each pair $(i,j)$, the set
\[
\{[\lambda]\in \mathcal{E}'_1\mid \Pi^*_{ij}([\lambda])=0\}
\]
has $\PThursM^x$-measure zero.
Since there are only finitely many pairs $(i,j)$, after removing the union of these null sets from $\mathcal{E}'_1$, we may assume that
$\Pi^*_{ij}([\lambda])\ne 0$ for all $1\le i,j\le g$ and all $[\lambda]$ in the resulting full-measure set.

We next prove condition (e).
Let $\mathcal M_g=\teich_g/\tmod(g)$, and equip
$Q^1\mathcal M_g$ with the Masur--Veech measure.
Choose a countable compact exhaustion
$\{K_n\}_{n\ge1}$ of $Q^1\mathcal M_g$.
Since the Teichm\"uller geodesic flow preserves
the finite Masur--Veech measure, the Poincar\'e
recurrence theorem implies that, for each $n$,
almost every point of $K_n$ returns to $K_n$
at arbitrarily large positive integer times.
Consequently, almost every point of
$Q^1\mathcal M_g$ determines a trajectory whose
projection returns infinitely often to a compact
subset of $\mathcal M_g$.

Let $\mu$ denote the lifted Masur--Veech measure
on $Q^1\teich_g$, let
$p\colon Q^1\teich_g\to\teich_g$ be the bundle
projection, and put $m=p_*\mu$.
Disintegrate $\mu$ over $m$, and denote the
conditional probability measure on the unit
cotangent sphere $S(x)$ by $s_x$.
The preceding recurrence statement holds for
$s_x$-almost every point of $S(x)$ for
$m$-almost every $x$.

Under the identification
\[
S(x)\longrightarrow\pml,
\qquad q\longmapsto[v(q)],
\]
the pushforward of $s_x$ is equivalent to the
Thurston measure class
(cf.\ \cite[\S2.3, Proposition 2.3(i), and \S2.3.1]
{Mirzakhani_etal2012}).
Choose a basepoint $x_*$ for which the preceding
almost-everywhere recurrence statement holds.
It follows that there is a measurable full-measure
set $\mathcal E_1''\subset\pml^{ue}$ such that
$\ray^\lambda_{x_*}$ projects to a recurrent ray
in $\mathcal M_g$ for every
$[\lambda]\in\mathcal E_1''$.

For any $x\in\teich_g$ and
$[\lambda]\in\mathcal E_1''$, Masur's
asymptoticity theorem implies that
$\ray^\lambda_x$ and $\ray^\lambda_{x_*}$
are asymptotic after a constant shift of their
time parameters (cf.\ \cite{Masur1982}).
Thus the projection of $\ray^\lambda_x$ also
returns infinitely often to a compact subset
of $\mathcal M_g$.
Indeed, a closed bounded neighborhood of a compact
subset of $\mathcal M_g$ is compact in the
quotient Teichm\"uller metric.
This proves condition (e) on $\mathcal E_1''$
for every basepoint $x$.

Let $F$ be the full-measure measurable subset obtained
after imposing conditions (b)--(e), and define
\[
\mathcal E_1
=
\bigcap_{\gamma\in\tmod(g)}\gamma^{-1}(F).
\]
Since $\tmod(g)$ is countable and its action preserves
the Thurston measure class, $\mathcal E_1$ has full measure.
By construction, $\mathcal E_1$ is $\tmod(g)$-invariant.
Moreover, $\mathcal E_1\subset F$, so conditions (b)--(e)
remain valid on $\mathcal E_1$.
Since $F$ consists of uniquely ergodic measured laminations,
condition (a) also holds.
This completes the proof.
\end{proof}



\subsection{Pushforward measures on $\partial \mathfrak{H}_g$}
For $x\in\teich_g$, let $\PSiegelThursM^x$ be the pushforward of the Thurston measure to $\overline{\mathfrak{H}_g}$ under $\Pi^*$:
\[
\PSiegelThursM^x=(\Pi^*)_*(\PThursM^x).
\]
% We first prove the following.

\begin{lemma}
\label{lem:null_siegel}
For any $x\in \teich_{g}$,
$\PSiegelThursM^x\bigl(\mathfrak{H}_g\bigr)=0$.
\end{lemma}

\begin{proof}
Fix $x_0\in\teich_g$.
We first show that $\PSiegelThursM^{x_0}(\{Z\})=0$ for every
$Z\in\mathfrak H_g$.
Since the period map is nonconstant and the Cayley transform is biholomorphic, some component $h=(\Cayley\circ\Pi)_{ij}$ is a nonconstant bounded holomorphic function on $\teich_g$.
If $\PSiegelThursM^{x_0}(\{Z\})>0$, then the radial limit of $h$ equals
the constant $(\Cayley(Z))_{ij}$ on a set of positive
Thurston measure.
By the identity theorem for radial limits
\cite[Theorem 1.2]{Miyachi2024Bounded}, this would imply
that $h$ is constant, a contradiction.

Suppose now that $\PSiegelThursM^{x_0}(\mathfrak H_g)>0$, and set
\[
E=\{[\lambda]\in\mathcal E_1
       \mid \Pi^*([\lambda])\in\mathfrak H_g\}.
\]
The equivariance of $\Pi^*$ implies that $E$ is invariant
under the mapping class group.
By ergodicity, $\PThursM^{x_0}(E)=1$.

Since $\psyp(2g,\mathbb{Z})$ acts properly on $\mathfrak{H}_g$, its orbit space is locally compact, second countable, and Hausdorff, and hence is a standard Borel space.
Let
\[
q\colon\mathfrak{H}_g\longrightarrow
\psyp(2g,\mathbb{Z})\backslash\mathfrak{H}_g
\]
be the quotient map.
The measurable map $q\circ\Pi^*$ on $E$ is invariant
under the mapping class group.
By ergodicity, the pushforward of $\PThursM^{x_0}|_E$ under $q\circ\Pi^*$ assigns either zero or one to every Borel set. A probability measure with this property on a standard Borel space is a Dirac measure. Thus $q\circ\Pi^*$ is essentially constant, and there exists $Z_0\in\mathfrak{H}_g$ such that
\[
\PThursM^{x_0}\bigl(\{[\lambda]\in E
          \mid \Pi^*([\lambda])\in \mathcal{O}\}\bigr)=1,
\]
where $\mathcal{O}$ is the $\psyp(2g,\mathbb{Z})$-orbit of $Z_0$.
However, the orbit $\mathcal{O}$ is countable, and the first part of the proof shows that
\[
\PSiegelThursM^{x_0}(\mathcal{O})
\le
\sum_{A\in\psyp(2g,\mathbb{Z})}\PSiegelThursM^{x_0}(\{A(Z_0)\})
=0.
\]
This is a contradiction.
Thus $\PSiegelThursM^{x_0}(\mathfrak H_g)=0$.
Since $x_0\in \teich_g$ was arbitrary, the lemma follows.
\end{proof}
% \begin{proof}
% Suppose, to the contrary, that $\PSiegelThursM^{x_0}(\mathfrak{H}_g)>0$.
% Let $F_0\subset \mathfrak{H}_g$ be a measurable fundamental set for the action of $\psyp(2g,\mathbb{Z})$.
% Namely, $F_0$ is a measurable subset of $\mathfrak{H}_g$ such that $F_0\cap A(F_0)=\emptyset$ for every $A\in \psyp(2g,\mathbb{Z})-\{id\}$, and
% \[
% \bigcup_{A\in \syp(2g,\mathbb{Z})} A(F_0)=\mathfrak{H}_g
% \]
%  (cf. \cite[IV]{Siegel1943}). From our assumption, $\PSiegelThursM^{x_0}(F_0)>0$.

% Take a measurable set $F'_0\subset F_0$ such that
% \begin{equation}
% \label{eq:positive-Fprime}
% 0<\PSiegelThursM^{x_0}(F'_0)
% <\frac{1}{2}\PSiegelThursM^{x_0}(F_0),
% \end{equation}
% and set
% \[
% F'=(\Pi^*)^{-1}\left(\bigcup_{A\in \psyp(2g,\mathbb{Z})}A(F'_0)\right).
% \]

% We first claim that $F'$ is a measurable $\mcg_{g}$-invariant set.
% %Indeed, let $[\omega]_*\in \mcg_{g}$
% %, and let $A_\varphi\in \syp(2g,\mathbb{Z})$ be the symplectic representation of $\varphi$. 
% If $[\lambda]\in F'$, then by definition there exists $A\in \psyp(2g,\mathbb{Z})$ such that
% \[
% (\Pi^*)^{-1}([\lambda])\in A(F'_0).
% \]
% Hence
% \[
% (\Pi^*)^{-1}([\omega]_*([\lambda]))
% =\symrep([\omega]_*)\circ (\Pi^*)^{-1}([\lambda])
% \in \symrep([\omega]_*)\circ A(F'_0),
% \]
% which implies that $[\omega]_*([\lambda])\in F'$.
% Therefore, $F'$ is $\mcg_{g}$-invariant.
% %By \eqref{eq:positive-Fprime}, we have $\PThursM^{x_0}(F')>0$.
% %Since the action of $\mcg_{g}$ on $\pml$ is ergodic, it follows that $\PThursM^{x_0}(F')=1$.
% %
% %On the other hand, using \eqref{eq:positive-Fprime} again, we obtain
% By \eqref{eq:positive-Fprime}, we have $\PThursM^{x_0}(F')>0$.
% Since the action of $\mcg_{g}$ on $\pml$ is ergodic, it follows that $\PThursM^{x_0}(F')=1$.

% Since $F'_0\subset F_0$ and $F_0$ is a fundamental set, the sets $A(F'_0)$, $A\in \psyp(2g,\mathbb{Z})$, are pairwise disjoint.
% Hence the corresponding preimages under $\Pi^*$ are also pairwise disjoint.
% Using \eqref{eq:positive-Fprime} again, we obtain
% \begin{align*}
% 1
% &=\PThursM^{x_0}(F')
% =\PSiegelThursM^{x_0}\left(\bigcup_{A\in \psyp(2g,\mathbb{Z})}A(F'_0)\right)=\sum_{A\in \psyp(2g,\mathbb{Z})}\PSiegelThursM^{x_0}(A(F'_0))\\
% &=\PSiegelThursM^{x_0}(F'_0)
% +\sum_{A\in \psyp(2g,\mathbb{Z})-\{id\}}\PSiegelThursM^{x_0}(A(F'_0)) \\
% &\le \frac{1}{2}\PSiegelThursM^{x_0}(F_0)
% +\sum_{A\in \psyp(2g,\mathbb{Z})-\{id\}}\PSiegelThursM^{x_0}(F'_0)(A(F_0))
% =1-\frac{1}{2}\PSiegelThursM^{x_0}(F_0)<1,
% \end{align*}
% which is a contradiction.
% \end{proof}

% \subsection{Value distribution of the radial limit and the support of the measure}
% Let
% {\color{red}この辺り，きちんと書き直すこと}
% \[
% \partialfinitefull\mathfrak{H}_g=
% \{
% \llangle Z_1,Z_2\rrangle\in \partialfinite\mathfrak{H}_g\mid \rank({\rm Im}(Z_1))=g-1
% \}
% \]
% (the superscript ``$f\!\!f$'' stands for ``finite and full rank'').
% The purpose of this section is to prove the following.

% \begin{theorem}[Radial limits are finite and full rank]
% \label{thm:rank}
% For almost all $[\lambda]\in \mathcal{E}_0$, we have
% \[
% (\pro\circ \Pi)^*([\lambda])\in \partialfinitefull\mathfrak{H}_g.
% \]
% Equivalently, for almost all $[\lambda]\in \mathcal{E}_0$, the period map $\Pi\colon\teich_{g}\to \mathfrak{H}_g$ admits a radial limit $\Pi^*([\lambda])$ in the boundary $\partial\mathfrak{H}_g$ of $\mathfrak{H}_g$ in $\symm(g,\mathbb{C})$, and this limit satisfies $\det(C\Pi^*([\lambda])+D)\ne 0$ for all $\begin{bmatrix} A & B \\ C & D\end{bmatrix}\in \syp(2g,\mathbb{R})$.
% \end{theorem}

% \begin{proof}
% By \Cref{lem:null_siegel}, we may assume that $\Pi^*([\lambda])\in \ppartial\mathfrak{H}_g$ for all $[\lambda]\in \mathcal{E}_0$.

% Let
% \[
% E_1=\{[\lambda]\in \mathcal{E}_0\mid (\pro\circ \Pi)^*([\lambda])\in \partial^{\infty}_{\mathbb{P}}\mathfrak{H}_g\}.
% \]
% Consider the holomorphic function $f_1$ on $\teich_{g}$ defined by
% \begin{align*}
% f_1(x)
% &=\det(I_g-(\Pi(x)-\sqrt{-1}I_g)(\Pi(x)+\sqrt{-1}I_g)^{-1}) \\
% &=(2\sqrt{-1})^g/\det(\Pi(x)+\sqrt{-1}I_g).
% \end{align*}
% Since $(\Pi(x)-\sqrt{-1}I_g)(\Pi(x)+\sqrt{-1}I_g)^{-1}=\Cayley^{-1}(\Pi(x))\in \mathfrak{D}_g$ for all $x\in \teich_{g}$, the function $f_1$ is bounded on $\teich_{g}$.
% Since the imaginary part of $\Pi(x)+\sqrt{-1}I_g$ is positive definite for all $x\in \teich_{g}$, the matrix $\Pi(x)+\sqrt{-1}I_g$ is non-singular (cf. \cite{MR1139765}).
% Therefore, $f_1$ does not vanish on $\teich_{g}$.

% On the other hand, by \Cref{prop:F-F-lemma3.1}\,(4), the radial limit $f_1^*$ of $f_1$ vanishes on $E_1$.
% Hence, if $E_1$ has positive measure, then the identity theorem (\cite[Theorem 1.2]{MR4830064}) implies that $f_1$ vanishes identically on $\teich_{g}$.
% This is a contradiction.

% Thus, $E_1$ has measure zero.
% Hence, by \Cref{prop:F-F-lemma3.1}\,(4), we may further assume that the radial limit $\Pi^*$ of the period map $\Pi\colon\teich_{g}\to \mathfrak{H}_g$ exists on $\mathcal{E}_0$ and satisfies
% \[
% \Pi^*([\lambda])\in \partial\mathfrak{H}_g\subset \symm(g,\mathbb{C})
% \]
% for almost all $[\lambda]\in \mathcal{E}_0$.
% Let
% \[
% E_2=\{[\lambda]\in \mathcal{E}_0\mid \rank(\Pi^*([\lambda]))<g\}.
% \]
% Consider another bounded holomorphic function $f_2$ on $\teich_{g}$ defined by
% \begin{align*}
% f_2(x)
% &=\det(I_g+(\Pi(x)-\sqrt{-1}I_g)(\Pi(x)+\sqrt{-1}I_g)^{-1}) \\
% &=2^g\det(\Pi(x))/\det(\Pi(x)+\sqrt{-1}I_g).
% \end{align*}
% Since ${\rm Im}(\Pi(x))>0$ for $x\in \teich_{g}$, the matrix $\Pi(x)$ is non-singular for all $x\in \teich_{g}$.
% Therefore, $f_2$ does not vanish on $\teich_{g}$.

% Since $\Pi^*([\lambda])\in \partial\mathfrak{H}_g$, its imaginary part is non-negative.
% Hence the imaginary part of $\Pi^*([\lambda])+\sqrt{-1}I_g$ is positive definite, and therefore $\Pi^*([\lambda])+\sqrt{-1}I_g$ is non-singular.
% Thus,
% \[
% f_2^*([\lambda])=2^g\det(\Pi^*([\lambda]))/\det(\Pi^*([\lambda])+\sqrt{-1}I_g)=0
% \]
% for all $[\lambda]\in E_2$.
% Applying the identity theorem again, we obtain $f_2\equiv 0$ on $\teich_{g}$.
% Hence, we obtain a contradiction.
% \end{proof}

% %
% \begin{corollary}
% The following hold:
% \begin{itemize}
% \item[(a)]
% For any $x\in \teich_{g}$, the measure $\PSiegelThursM^x$ is supported on the boundary $\partial\mathfrak{H}_g$ of $\mathfrak{H}_g$ in $\symm(g,\mathbb{C})$;
% \item[(b)]
% For any $x,y\in \teich_{g}$, the measures $\PSiegelThursM^x$ and
% $\PSiegelThursM^y$ are mutually absolutely continuous; and
% \item[(c)]
% For each $x\in \teich_{g}$, the measure $\PSiegelThursM^x$ is quasi-invariant and ergodic with respect to the action of $\syp(2g,\mathbb{Z})$.
% \end{itemize}
% \end{corollary}

% \begin{proof}
% {\bf (a)} It follows from \Cref{lem:null_siegel}.

% \medskip
% \noindent
% {\bf (b)}\quad Since $\PThursM^x$ and $\PThursM^y$ are mutually absolutely continuous, their
% pushforward measures $\PSiegelThursM^x$ and $\PSiegelThursM^y$ are also mutually
% absolutely continuous.

% \medskip
% \noindent
% {\bf (c)}\quad
% From \Cref{lem:null_siegel}, the Radial limit map $\Pi^*$ defines a semiconjugacy $\pml\to \partial \mathfrak{H}_g$ between the actions of $\tmod(g)$ on $\pml$ and $\syp(2g,\mathbb{Z})$ on $\partial\mathfrak{H}_g$ via the symplectic representation $\symrep\colon \mcg(\Sigma_g)\to \syp(2g,\mathbb{Z})$, the quasi-invariance and the ergodicity of $\PSiegelThursM^x$ follows from the quasi-invariance and the ergodicity of $\PThursM^x$. 
% % Indeed, the action of $\syp(2g,\mathbb{Z})$ on $\partial\mathfrak{H}_g$ is the pushforward of the action of $\mcg(\Sigma_g)$ on $\pml$ under the radial limit map $\Pi^*$.
% % We next prove ergodicity. Let
% % $E\subset \partial\mathfrak{H}_g$ be a $\syp(2g,\mathbb{Z})$-invariant measurable set.
% % Since $\PSiegelThursM^x$ is the pushforward of the Thurston measure
% % $\PThursM^x$ under the radial limit map $\Pi^*$, we have
% % \[
% % \PSiegelThursM^x(E)
% % =
% % \PThursM^x\bigl((\Pi^*)^{-1}(E)\bigr).
% % \]
% % Set
% % \[
% % E'=(\Pi^*)^{-1}(E)\cap \mathcal{E}_1\subset \pml.
% % \]
% % By the equivariance of the radial limit map and the surjectivity of the
% % symplectic representation
% % \[
% % \symrep\colon \mcg(\Sigma_g)\to \syp(2g,\mathbb{Z}),
% % \]
% % the set $E'$ is $\mcg(\Sigma_g)$-invariant.
% % Since the action of $\mcg(\Sigma_g)$ on $\pml$ is ergodic with respect to
% % the Thurston measure, we have
% % \[
% % \PThursM^x(E')=0
% % \quad\text{or}\quad
% % \PThursM^x(E')=1.
% % \]
% % Since $\mathcal{E}_1$ has full measure, it follows that
% % \[
% % \PSiegelThursM^x(E)=0
% % \quad\text{or}\quad
% % \PSiegelThursM^x(E)=1.
% % \]
% % This proves that $\PSiegelThursM^x$ is ergodic with respect to the action of
% % $\syp(2g,\mathbb{Z})$.
% \end{proof}


% \section{The limit set of the Torelli group}

\subsection{The conical limit set of the Torelli group}
We characterize the conical limit points of the Torelli group $\torelliG_g$ on $\mathcal E_1$ in terms of the radial limit of the period map.

\begin{proposition}
    \label{prop:conical_limit_set_torelli_group}
    Let $[\lambda]\in \mathcal{E}_1$. The following are equivalent.
    \begin{itemize}
    \item[(a)] $[\lambda]$ is a conical limit point of $\torelliG_{g}$; 
    \item[(b)] the limit $\Pi^*([\lambda])$ lies in $\mathfrak{H}_g$; and
    \item[(c)] the limit $\Pi^*([\lambda])$ lies in the image of the period map $\Pi\colon \teich_{g}\to \mathfrak{H}_g$.
    \end{itemize}
\end{proposition}

\begin{proof}
{\bf (a) $\Rightarrow$ (b).}
There exist a point $x_1\in \teich_{g}$, possibly different from the fixed basepoint $x_0$, a sequence of distinct elements $\{\varphi_n\}_{n=1}^\infty$ in $\torelliG_{g}$, a sequence $\{t_n\}_{n=1}^\infty$ in $[0,\infty)$ with $t_n\to \infty$, and a constant $D>0$ such that
\[
d_T(\ray^{\lambda}_{x_1}(t_n),\varphi_n(x_0))\le D
\]
for all $n\ge 1$.
Since the Teichm\"uller distance is the Kobayashi distance and the period map is holomorphic, the distance-decreasing property implies that
\[
d_{\mathfrak{H}_g}(\Pi(\ray^{\lambda}_{x_1}(t_n)),\Pi(x_0))
=
d_{\mathfrak{H}_g}(\Pi(\ray^{\lambda}_{x_1}(t_n)),\Pi(\varphi_n(x_0)))
\le D
\]
for all $n\ge 1$.
Since $\mathfrak{H}_g$ is a proper complete metric space with respect to its invariant Kobayashi distance, the sequence
$\{\Pi(\ray^{\lambda}_{x_1}(t_n))\}_{n=1}^{\infty}$
has a subsequence converging to a point of $\mathfrak{H}_g$.
On the other hand, by the definition of the radial limit, this sequence converges to $\Pi^*([\lambda])$.
Hence $\Pi^*([\lambda])\in \mathfrak{H}_g$.

\medskip
\noindent
{\bf (b) $\Rightarrow$ (c).}
Assume that $\Pi^*([\lambda])\in \mathfrak{H}_g$.
Since
\[
\Pi^*([\lambda])=
\lim_{t\to \infty}\Pi(\ray^{\lambda}_{x_1}(t)),
\]
we have
$\Pi^*([\lambda])\in \overline{\Pi(\teich_g)}\cap \mathfrak{H}_g$,
where the closure is taken in $\mathfrak{H}_g$.

Suppose, to the contrary, that
$\Pi^*([\lambda])\in \overline{\Pi(\teich_g)}\setminus \Pi(\teich_g)$.
Then, 
% by the description of the closure of the Torelli locus recalled in
% \S\ref{subsec:period_map}, this limit is the period matrix of a reducible principally polarized abelian variety, equivalently to a stable curve of compact type.
% Since the closure $\overline{\Pi(\teich_g)}$ of the image is the covering space of the moduli space $\tilde{\mathcal{M}}_g$ of stable curves of compact type, whose covering projection is the quotient map by the properly discontinuous action of $\syp(2g,\mathbb{Z})$, 
by \Cref{prop:compacttype},
the underlying Riemann surfaces along the ray
$\ray^{\lambda}_{x_1}(t)$ leave every compact subset of moduli space as $t\to \infty$.
On the other hand, since
$[\lambda]\in \mathcal{E}_1$, the ray $\ray^{\lambda}_{x_1}$ is recurrent in the moduli space by \Cref{lem:radial_limit}\,(e).
This contradicts the preceding divergence in moduli space.
% Therefore,
% \[
% \Pi^*([\lambda])\in \Pi(\teich_g).
% \]
% {\bf (b) $\Rightarrow$ (c).}
% Assume that $\Pi^*([\lambda])\in \mathfrak{H}_g$.
% Since
% \[
% \Pi^*([\lambda])=
% \lim_{t\to \infty}\Pi(\ray_{\lambda}^{x_1}(t)),
% \]
% we have
% $\Pi^*([\lambda])\in \overline{\Pi(\teich_g)}\cap \mathfrak{H}_g$,
% where the closure is taken in $\mathfrak{H}_g$.

% Suppose, to the contrary, that
% \[
% \Pi^*([\lambda])\in \overline{\Pi(\teich_g)}\setminus \Pi(\teich_g).
% \]
% By the description of the closure of the Torelli locus recalled in
% \S\ref{subsec:period_map}, this limit corresponds to a reducible principally
% polarized abelian variety, equivalently to a stable curve of compact type.

% On the other hand, since $[\lambda]\in \mathcal{E}_1$, the ray $\ray_{\lambda}^{x_1}$ is recurrent in the moduli space by \Cref{lem:radial_limit}\,(d).
% Therefore, there exist a sequence $\{s_n\}_{n=1}^{\infty}$ in $[0,\infty)$ with $s_n\to\infty$, a sequence $\{\omega_n\}_{n=1}^{\infty}$ in $\mcg(\Sigma_g)$, and a point $y_\infty\in \teich_g$ such that, after passing to a subsequence,
% \[
% \omega_n(\ray_{\lambda}^{x_1}(s_n))\to y_\infty
% \]
% in $\teich_g$.
% Set
% \[
% A_n=\symrep(\omega_n)\in \syp(2g,\mathbb{Z}).
% \]
% Since $\syp(2g,\mathbb{Z})$ is countable, we may pass to a further subsequence so that either $\{A_n\}_{n=1}^{\infty}$ is constant or it consists of pairwise distinct elements.

% First suppose that $A_n=A$ for all $n$.
% By the equivariance and continuity of the period map, we have
% \[
% A\bigl(\Pi^*([\lambda])\bigr)
% =
% \lim_{n\to\infty}A\bigl(\Pi(\ray_{\lambda}^{x_1}(s_n))\bigr)
% =
% \lim_{n\to\infty}\Pi(\omega_n(\ray_{\lambda}^{x_1}(s_n)))
% =
% \Pi(y_\infty).
% \]
% By \Cref{lem:radial_limit}\,(c), the inverse action $A^{-1}$ is defined at $\Pi(y_\infty)$.
% Thus
% \[
% \Pi^*([\lambda])=A^{-1}(\Pi(y_\infty))\in \Pi(\teich_g),
% \]
% which contradicts our assumption.

% It remains to rule out the case where the elements $A_n$ are pairwise distinct.
% In this case, the convergence
% \[
% A_n\bigl(\Pi(\ray_{\lambda}^{x_1}(s_n))\bigr)
% =
% \Pi(\omega_n(\ray_{\lambda}^{x_1}(s_n)))
% \to \Pi(y_\infty)
% \]
% would give a sequence of distinct symplectic transforms of points converging to the boundary point $\Pi^*([\lambda])$ whose images remain in a compact subset of $\mathfrak{H}_g$.
% This contradicts the proper discontinuity of the action of $\syp(2g,\mathbb{Z})$ on $\mathfrak{H}_g$, together with \Cref{lem:radial_limit}\,(c), which guarantees that the symplectic action is defined at $\Pi^*([\lambda])$.
% Therefore this case cannot occur.

% Hence the assumption
% $\Pi^*([\lambda])\in \overline{\Pi(\teich_g)}\setminus \Pi(\teich_g)$
% is impossible, and so
% \[
% \Pi^*([\lambda])\in \Pi(\teich_g).
% \]

\medskip
\noindent
{\bf (c) $\Rightarrow$ (a).}
% We recall that, since the quotient map $\pi_{Tor}\colon \teich_{g}\to \torelli_{g}$ is a universal covering map, it follows from \cite[Chapter IV, Proposition 1.6]{MR2194466} that the Kobayashi distance on $\torelli_{g}$ is the quotient metric of the Teichm\"uller distance on $\teich_{g}$.
Suppose that $\Pi^*([\lambda])\in \Pi(\teich_g)$, and let $x_1\in \teich_{g}$.
Since
\[
\Pi(\ray^{\lambda}_{x_1}(t))\to \Pi^*([\lambda])
\]
as $t\to \infty$, we claim that the ray
\[
\{\prG_{Tor}(\ray^{\lambda}_{x_1}(t))\in \torelli_g \mid t\ge 0\}
\]
is contained in a compact subset of $\torelli_g$, where $\prG_{Tor}\colon \teich_{g}\to \torelli_{g}$ is the quotient map.
Indeed, otherwise there would exist an increasing sequence $\{t_n\}_{n=1}^{\infty}$ in $[0,\infty)$ with $t_n\to\infty$ such that $\prG_{Tor}(\ray^{\lambda}_{x_1}(t_n))$ leaves every compact subset of $\torelli_g$.
By \Cref{prop:compacttype}, the limit
$\Pi^*([\lambda])$ would represent a decomposable
principally polarized abelian variety.
This contradicts
$\Pi^*([\lambda])\in\Pi(\teich_g)$, since the
principally polarized Jacobian of a smooth curve
is indecomposable.
Thus $\prG_{Tor}(\ray^{\lambda}_{x_1}(t))$ is contained in a compact subset of $\torelli_g$.

Hence there is $D>0$ such that, for every $t\ge 0$, $\prG_{Tor}(x_0)$ and $\prG_{Tor}(\ray^{\lambda}_{x_1}(t))$ can be joined by a smooth path of length at most $D$ in $\torelli_g$ with respect to the Kobayashi distance (cf. \cite{Royden1971b}).
By lifting the paths, for every $t\geq 0$, we find an element $\varphi_t\in \torelliG_{g}$ satisfying
\[
d_T(\ray^{\lambda}_{x_1}(t),\varphi_t(x_0))\le D.
\]
Since the Teichm\"uller ray diverges in $\teich_g$, the orbit of $x_0$ under $\torelliG_g$ meets the $D$-neighborhood of the ray $\ray^{\lambda}_{x_1}$ infinitely often.
This proves that $[\lambda]$ is a conical limit point of $\torelliG_{g}$.
\end{proof}


From \Cref{lem:null_siegel} and \Cref{prop:conical_limit_set_torelli_group}, we obtain the following.

\begin{corollary}[The conical limit set has measure zero]
\label{cor:conical_limit_sets}
The conical limit set $\Lambda_C(\torelliG_{g})$ of the Torelli group has measure zero with respect to the Thurston measure.
\end{corollary}

In contrast, the conical limit set of the Teichm\"uller modular group $\tmod(g)$ has full measure. This follows from the conservativity of the Teichm\"uller geodesic flow and the fact that almost every Teichm\"uller geodesic ray returns infinitely often to a compact set (cf. \cite{Masur1982}).
Moreover, the fixed points of any pseudo-Anosov element in a subgroup of $\tmod(g)$ belong to the conical limit set of that subgroup.
The contrast above concerns Thurston measure on the boundary. For results on the genericity of pseudo-Anosov elements, see, for example, \cite{ChoiInhyeok2024, VivekaSoutoJing2020, Maher2011, MalesteinSouto2013, Rivin2008, Rivin2009}.



\begin{corollary}
With respect to the measure induced by the Masur--Veech measure,
almost every trajectory of the Teichm\"uller geodesic flow
on the unit cotangent bundle of the Torelli space
$\teich_g/\torelliG_g$ is non-recurrent in forward time.
\end{corollary}

\begin{proof}
If the projection of a forward trajectory to the
Torelli space returns infinitely often to a compact
subset, its vertical projective foliation belongs
to $\Lambda_C(\torelliG_g)$.
Indeed, compactness in the quotient implies that,
along a sequence of times tending to infinity,
a lift of the trajectory remains within a bounded
Teichm\"uller distance of a fixed
$\torelliG_g$-orbit.

By \Cref{cor:conical_limit_sets}, this set of
directions has zero Thurston measure.
Using the equivalence of the conditional measures
on unit cotangent spheres with the Thurston
measure class, as in the proof of
\Cref{lem:radial_limit}(e), disintegration shows
that the corresponding set of initial conditions
has zero Masur--Veech measure.
The same holds on the quotient by $\torelliG_g$,
which proves the assertion.
\end{proof}

\subsection{Complement of the small horospherical limit set}
We close this section with the following observation.

\begin{proposition}[$\pml\setminus \Lambda_h(\torelliG_g)\ne \emptyset$]
\label{prop:complement_small_horospherical_limit_set}
Let $g\ge 2$ and $x_0=(M_0,f_0)\in\teich_g$. If $[\lambda]\in\pml$ is represented by the vertical measured foliation of a nonzero holomorphic $1$-form on $M_0$, then $[\lambda]\in\operatorname{Dir}_{\torelliG_g}(x_0)\setminus\Lambda_h(\torelliG_g)$.
\end{proposition}

\begin{proof}
Choose an abelian differential $\eta$ on $M_0$ such that $v(\eta^2)=\lambda$. By \cite{Kra1981}, $\Pi\circ \ray^\lambda_{x_0}\colon [0,\infty)\to \mathfrak{H}_g$ is an isometric embedding with respect to the usual distance on $[0,\infty)$ and the Kobayashi distance on $\mathfrak{H}_g$.
By the distance-decreasing property of the Kobayashi distance, for any $[\omega]\in \torelliG_g$ and $t\ge 0$,
\begin{align*}
d_T([\omega](x_0),\ray^\lambda_{x_0}(t))
&\ge d_{\mathfrak{H}_g}(\Pi\circ [\omega](x_0),\Pi\circ \ray^\lambda_{x_0}(t)) \\
&\ge d_{\mathfrak{H}_g}(\Pi(x_0),\Pi\circ \ray^\lambda_{x_0}(t)) \\
&= d_{\mathfrak{H}_g}(\Pi\circ \ray^\lambda_{x_0}(0),\Pi\circ \ray^\lambda_{x_0}(t))=t.
\end{align*}
Therefore, $\busemann^{x_0}_\lambda([\omega](x_0))\ge 0$ for all $[\omega]\in\torelliG_g$, so $[\lambda]\notin\Lambda_h(\torelliG_g)$. Since $\busemann^{x_0}_\lambda(x_0)=0$, we also have $[\lambda]\in\operatorname{Dir}_{\torelliG_g}(x_0)$.
\end{proof}
% \section{Conclusion}

% \subsection{Garnett points}
% A point $[\lambda]\in\pml$ is called a \emph{Garnett point} of $G$ if $\horo^{x_0}_G([\lambda])>0$ but this infimum is not attained on $G(x_0)$ (cf.\ \cite[\S2.6]{Nicholls1989}). The Garnett points are contained in $\Lambda_H(G)$.



\section{Appendix}

\subsection{Coordinates for the once-punctured torus}

We recall explicit coordinates for the Teichm\"uller
space and the space of measured laminations of the
once-punctured torus.
These coordinates will be used to prove the nullity
of level loci in this case.

Henceforth, we identify $\Sigma_{1,1}$ with
$(\mathbb{R}^2\setminus\mathbb{Z}^2)/\mathbb{Z}^2$, equipped with the orientation induced by the standard orientation of $\mathbb{R}^2$. Let $p_0\in\Sigma_{1,1}$ be the image of $(1/2,1/2)$ under the quotient projection. Let $A$ and $B$ be the simple closed curves obtained as the images of the horizontal and vertical lines through $(1/2,1/2)$, respectively. We orient $A$ in the direction of increasing first coordinate and $B$ in the direction of increasing second coordinate. Regarded as loops based at $p_0$, the curves $A$ and $B$ freely generate the fundamental group $\pi_1(\Sigma_{1,1},p_0)$.

In \S\S\ref{subsubsec:TeichmullerSpace11}--\ref{subsubsec:Extremal-length11}, we recall basic facts about the Teichm\"uller space of $\Sigma_{1,1}$; see \cite{miyachi2025functiontheorydynamicsergodic} for proofs.

\subsubsection{Teichm\"uller space of $\Sigma_{1,1}$}
\label{subsubsec:TeichmullerSpace11}
For $\tau\in \mathbb{H}=\{\tau\in \mathbb{C}\mid \operatorname{Im}(\tau)>0\}$, we define $\Gamma_\tau$ as the lattice in $\mathbb{C}$ generated by $1$ and $\tau$.
Let $x=(M,f)\in \teich_{1,1}$.
By the uniformization theorem, there is a unique $\tau=\tau(x)\in \mathbb H$ such that $M$ is biholomorphic to $M_\tau=(\mathbb{C}-\Gamma_\tau)/\Gamma_\tau$ and $f_*(A)$ and $f_*(B)$ correspond to the generators $1$ and $\tau$ of $\Gamma_\tau$.
The correspondence
\[
\teich_{1,1}\ni x\mapsto \tau(x)\in \mathbb{H}
\]
is biholomorphic. Henceforth, we identify $\teich_{1,1}$ with $\mathbb{H}$ in this manner.

\subsubsection{Measured laminations on $\Sigma_{1,1}$}
\label{subsubsec:ML11}
Let $p$, $q\in \mathbb{Z}$. In what follows, we assume that $q\ge 0$ and $\operatorname{gcd}(p,q)=1$ when $q>0$, and $p=1$ when $q=0$.
The \emph{$p/q$-curve} on $\Sigma_{1,1}$ is a simple closed curve homotopic to the quotient curve of the line with direction $(-p,q)$ in $\mathbb{R}^2$. Any nontrivial and non-peripheral simple closed curve is a $p/q$-curve for some $p/q\in \hat{\mathbb{Q}}=\mathbb{Q}\cup\{\infty\}$ where $1/0=\infty$.
In fact, the set of homotopy classes of nontrivial and non-peripheral simple closed curves is identified with $\hat{\mathbb{Q}}$.
Let $\gamma_{p/q}$ be the homotopy class of a $p/q$-curve.
The embedding
\[
\mathbb{R}_{>0}\times \hat{\mathbb{Q}}\ni (t,p/q)\mapsto (-tp,tq)\in \mathbb{R}\times \mathbb{R}_{\ge 0}
\]
induces the identification 
\begin{equation}
\label{eq:ml-coordinates11}
\ml=\ml(\Sigma_{1,1})\cong \mathbb{R}^2/\mathbb{Z}_2
\end{equation}
where $\mathbb{Z}_2=\mathbb{Z}/2\mathbb{Z}$ acts on $\mathbb{R}^2$ by rotation through $\pi$.
The real-analytic structure on $\ml(\Sigma_{1,1})$ induced by the identification with $\mathbb{R}^2/\mathbb{Z}_2$ coincides with the real-analytic structure arising from train-track coordinates.
The Thurston measure on $\ml$ is the quotient of the Lebesgue measure on $\mathbb{R}^2$.

\subsubsection{Hubbard--Masur differentials and extremal lengths}
\label{subsubsec:Extremal-length11}
For $x=(M,f)\in \teich_{1,1}\cong \mathbb{H}$ and $\lambda=[a,b]\in \ml(\Sigma_{1,1})=\mathbb{R}^2/\mathbb{Z}_2$, the Hubbard--Masur differential $q_{\lambda,x}$ is
\begin{equation}
\label{eq:Hubbard-Masur-differential11}
q_{\lambda,x}=-\left(\frac{a+b\overline{\tau(x)}}{\operatorname{Im}(\tau(x))}\right)^2dz^2
\end{equation}
on $M=M_{\tau(x)}$, where $z$ denotes the local coordinate induced by the standard coordinate on the abelian cover $\mathbb{C}-\Gamma_
{\tau(x)}$. 

\subsection{Level loci for $(g,m)=(1,1)$}

We prove \Cref{thm:level_set} for $(g,m)=(1,1)$.
By \eqref{eq:Hubbard-Masur-differential11}, the extremal length
of $\lambda=[a,b]\in\ml$ at $x\in\teich_{1,1}$ is
\begin{equation}
\label{eq:extremal_length-formula11}
\ext_x(\lambda)
=\frac{|a+b\tau(x)|^2}{\operatorname{Im}(\tau(x))}.
\end{equation}

Fix distinct points $x_1,x_2\in\teich_{1,1}$ and $c\in\mathbb R$.
Write $\tau(x_j)=\xi_j+i\eta_j$, where $\eta_j>0$, for $j=1,2$.
Under the identification
$\ml\cong\mathbb R^2/\mathbb Z_2$,
the level locus $\ml_c(x_1,x_2)$ is the image of the zero set
of the homogeneous quadratic polynomial
\[
P_c(a,b)
=
\frac{(a+b\xi_1)^2+b^2\eta_1^2}{\eta_1}
-
e^c\frac{(a+b\xi_2)^2+b^2\eta_2^2}{\eta_2}.
\]

We claim that $P_c$ is not identically zero.
Otherwise, comparison of the coefficients of $a^2$ and $ab$
would give
\[
\eta_2=e^c\eta_1,
\qquad
\xi_1=\xi_2.
\]
Comparison of the coefficients of $b^2$ would then give
$\eta_1=e^c\eta_2$.
Since $\eta_1,\eta_2>0$, these identities imply $c=0$
and $\tau(x_1)=\tau(x_2)$, contradicting $x_1\ne x_2$.

The zero set of a nonzero homogeneous quadratic polynomial
on $\mathbb R^2$ is contained in the union of at most two
lines through the origin.
It therefore has Lebesgue measure zero.
Since the Thurston measure on $\ml$ is the quotient
of Lebesgue measure on $\mathbb R^2$,
the level locus $\ml_c(x_1,x_2)$ is $\ThursM$-null.
Its projectivization $\pml_c(x_1,x_2)$ consists of at most
two points and is consequently $\PThursM^x$-null
for every $x\in\teich_{1,1}$.
This proves \Cref{thm:level_set} for $(g,m)=(1,1)$.

\subsection{Level loci for $(g,m)=(0,4)$}

We deduce the four-punctured sphere case from the
once-punctured torus case.

Fix a topological double cover of the sphere
branched over its four marked points, and choose
one ramification point as the marked point of
the covering torus.
Lifting complex structures gives the standard
identification
\[
L\colon\teich_{0,4}\longrightarrow\teich_{1,1}.
\]
Here the other three ramification points are
filled in.

An essential simple closed curve $\gamma$ on
the four-punctured sphere separates the punctures
into two pairs.
Its inverse image consists of two curves
representing the same essential curve class
$\alpha$ on the once-punctured torus.
The correspondence $\gamma\mapsto\alpha$, extended
homogeneously and continuously using slope
coordinates, gives a homeomorphism
\[
T\colon\ml_{0,4}\longrightarrow\ml_{1,1}.
\]
In particular, it induces a homeomorphism
\[
\mathbb P T\colon\pml_{0,4}\longrightarrow\pml_{1,1}.
\]

For $x\in\teich_{0,4}$ and an essential simple
closed curve $\gamma$, the extremal length
formula for the double cover gives
\[
\ext_{L(x)}(2T(\gamma))=2\ext_x(\gamma)
\]
(cf.\ \cite[Lemma 4.1 and \S4]
{FortierBourqueEtAl2024}).
Since extremal length is homogeneous of degree two,
\[
\ext_x(\gamma)=2\ext_{L(x)}(T(\gamma)).
\]
By homogeneity, continuity, and the density of
weighted simple closed curves, this identity
extends to every $\lambda\in\ml_{0,4}$:
\[
\ext_x(\lambda)=2\ext_{L(x)}(T(\lambda)).
\]

Now let $x_1,x_2\in\teich_{0,4}$ be distinct,
and let $c\in\mathbb R$.
The preceding identity gives
\[
\ext_{x_1}(\lambda)=e^c\ext_{x_2}(\lambda)
\quad\Longleftrightarrow\quad
\ext_{L(x_1)}(T(\lambda))
=e^c\ext_{L(x_2)}(T(\lambda)).
\]
Since $L(x_1)\ne L(x_2)$, the once-punctured
torus calculation shows that the projective
classes satisfying the right-hand equality
form a set of at most two points.
Thus $\pml_c(x_1,x_2)$ also consists of at most
two points and has zero projectivized Thurston
measure.

Moreover, $\ml_c(x_1,x_2)$ is contained in the
union of at most two rays, together with the
origin.
Each ray is locally contained in a one-dimensional
linear subspace in train-track coordinates and
therefore has zero Thurston measure.
This proves \Cref{thm:level_set} for $(g,m)=(0,4)$.
\def\cprime{$'$} \def\cprime{$'$}
\begin{thebibliography}{10}

\bibitem{Ahlfors1960}
Lars~V. Ahlfors.
\newblock {T}he complex analytic structure of the space of closed {R}iemann surfaces.
\newblock In {\em Analytic functions}, pages 45--66. Princeton Univ. Press, Princton, N.J., 1960.

\bibitem{Mirzakhani_etal2012}
Jayadev Athreya, Alexander Bufetov, Alex Eskin, and Maryam Mirzakhani.
\newblock {L}attice point asymptotics and volume growth on {T}eichm{\"u}ller space.
\newblock {\em Duke Math. J.}, 161(6):1055--1111, 2012.

\bibitem{Beardon1995}
Alan~F. Beardon.
\newblock {\em The geometry of discrete groups}, volume~91 of {\em Graduate Texts in Mathematics}.
\newblock Springer-Verlag, New York, 1995.
\newblock Corrected reprint of the 1983 original.

\bibitem{BehrstockMargalit2006}
Jason Behrstock and Dan Margalit.
\newblock Curve complexes and finite index subgroups of mapping class groups.
\newblock {\em Geom. Dedicata}, 118:71--85, 2006.

\bibitem{Bers1981}
Lipman Bers.
\newblock {F}inite-dimensional {T}eichm{\"u}ller spaces and generalizations.
\newblock {\em Bull. Amer. Math. Soc. (N.S.)}, 5(2):131--172, 1981.

\bibitem{BersBoundary1981}
Lipman Bers.
\newblock {T}he action of the modular group on the complex boundary.
\newblock In {\em Riemann surfaces and related topics: Proceedings of the 1978 Stony Brook Conference (State Univ. New York, Stony Brook, N.Y., 1978)}, volume~97 of {\em Ann. of Math. Stud.}, pages 33--52. Princeton Univ. Press, Princeton, N.J., 1981.

\bibitem{BirmanLubotzkyMcCarthy1983}
Joan~S. Birman, Alex Lubotzky, and John McCarthy.
\newblock Abelian and solvable subgroups of the mapping class groups.
\newblock {\em Duke Math. J.}, 50(4):1107--1120, 1983.

\bibitem{BirmanSeries1985}
Joan~S. Birman and Caroline Series.
\newblock {G}eodesics with bounded intersection number on surfaces are sparsely distributed.
\newblock {\em Topology}, 24(2):217--225, 1985.

\bibitem{Bonahon2001}
Francis Bonahon.
\newblock {G}eodesic laminations on surfaces.
\newblock In {\em Laminations and foliations in dynamics, geometry and topology (Stony Brook, NY, 1998)}, volume 269 of {\em Contemp. Math.}, pages 1--37. Amer. Math. Soc., Providence, RI, 2001.

\bibitem{Busemann1955}
Herbert Busemann.
\newblock {\em {T}he geometry of geodesics}.
\newblock Academic Press, Inc., New York, 1955.

\bibitem{ChoiInhyeok2024}
Inhyeok Choi.
\newblock {P}seudo-{A}nosovs are exponentially generic in mapping class groups.
\newblock {\em Geom. Topol.}, 28(4):1923--1955, 2024.

\bibitem{choi2024confinedsubgroupsgroupscontracting}
Inhyeok Choi, Ilya Gekhtman, Wenyuan Yang, and Tianyi Zheng.
\newblock Confined subgroups in groups with contracting elements.
\newblock Preprint arXiv:2405.09070, 2024.

\bibitem{ChoiKim2025InvariantMeasures}
Inhyeok Choi and Dongryul~M. Kim.
\newblock Invariant measures on the space of measured laminations for subgroups of mapping class group.
\newblock Preprint arXiv:2510.23256, 2025.

\bibitem{Fathi1979}
Adrian Douady, Albert Fathi, David Fried, Fran{\c{c}}ois Laudenbach, Valentin Po{\'e}naru, and Michael Shub.
\newblock {\em Travaux de {T}hurston sur les surfaces}, volume~66 of {\em Ast\'erisque}.
\newblock Soci\'et\'e Math\'ematique de France, Paris, 1979.
\newblock S\'eminaire Orsay, With an English summary.

\bibitem{Dumas2015}
David Dumas.
\newblock Skinning maps are finite-to-one.
\newblock {\em Acta Math.}, 215(1):55--126, 2015.

\bibitem{Earle1978}
Clifford~J. Earle.
\newblock {F}amilies of {R}iemann surfaces and {J}acobi varieties.
\newblock {\em Ann. of Math. (2)}, 107(2):255--286, 1978.

\bibitem{EarleKra1974}
Clifford~J. Earle and Irwin Kra.
\newblock {O}n isometries between {T}eichm{\"u}ller spaces.
\newblock {\em Duke Math. J.}, 41:583--591, 1974.

\bibitem{EberleinONeill1973}
P.~Eberlein and B.~O\textquotesingle~Neill.
\newblock {V}isibility manifolds.
\newblock {\em Pacific J. Math.}, 46:45--109, 1973.

\bibitem{VivekaSoutoJing2020}
Viveka Erlandsson, Juan Souto, and Jing Tao.
\newblock {G}enericity of pseudo-{A}nosov mapping classes, when seen as mapping classes.
\newblock {\em Enseign. Math.}, 66(3-4):419--439, 2020.

\bibitem{FarbMargalit2012}
Benson Farb and Dan Margalit.
\newblock {\em {A} primer on mapping class groups}, volume~49 of {\em Princeton Mathematical Series}.
\newblock Princeton University Press, Princeton, NJ, 2012.

\bibitem{FarkasKra1992}
Harschel Farkas and Irwin Kra.
\newblock {\em {R}iemann surfaces}, volume~71 of {\em Graduate Texts in Mathematics}.
\newblock Springer-Verlag, New York, second edition, 1992.

\bibitem{FortierBourqueEtAl2024}
Maxime Fortier~Bourque, D\'idac Mart\'inez-Granado, and Franco Vargas~Pallete.
\newblock The extremal length systole of the {B}olza surface.
\newblock {\em Ann. H. Lebesgue}, 7:1409--1455, 2024.

\bibitem{FriedlandFreitas2004}
Shmuel Friedland and Pedro~J. Freitas.
\newblock {R}evisiting the {S}iegel upper half plane. {I}.
\newblock {\em Linear Algebra Appl.}, 376:19--44, 2004.

\bibitem{Hain2006}
Richard Hain.
\newblock {F}initeness and {T}orelli spaces.
\newblock In {\em Problems on mapping class groups and related topics}, volume~74 of {\em Proc. Sympos. Pure Math.}, pages 57--70. Amer. Math. Soc., Providence, RI, 2006.

\bibitem{Helgason2001}
Sigurdur Helgason.
\newblock {\em {D}ifferential geometry, {L}ie groups, and symmetric spaces}, volume~34 of {\em Graduate Studies in Mathematics}.
\newblock American Mathematical Society, Providence, RI, 2001.
\newblock Corrected reprint of the 1978 original.

\bibitem{HubbardMasur1979}
John Hubbard and Howard Masur.
\newblock {Q}uadratic differentials and foliations.
\newblock {\em Acta Math.}, 142(3-4):221--274, 1979.

\bibitem{HubbardVol3_2022}
John~Hamal Hubbard.
\newblock {\em Teichm\"uller theory and applications to geometry, topology, and dynamics. {V}olume 3. {M}anifolds that fiber over the circle}.
\newblock Matrix Editions, Ithaca, NY, [2022] \copyright 2022.
\newblock With a foreword by William Thurston.

\bibitem{ImayoshiTaniguchi1992}
Yoichi Imayoshi and Masahiko Taniguchi.
\newblock {\em {A}n introduction to {T}eichm{\"u}ller spaces}.
\newblock Springer-Verlag, Tokyo, 1992.

\bibitem{Ivanov1992}
Nikolai~V. Ivanov.
\newblock {\em {S}ubgroups of {T}eichm{\"u}ller modular groups}, volume 115 of {\em Translations of Mathematical Monographs}.
\newblock American Mathematical Society, Providence, RI, 1992.
\newblock Translated from the Russian by E. J. F. Primrose and revised by the author.

\bibitem{Ivanov2001}
Nikolai~V. Ivanov.
\newblock Isometries of {T}eichm\"uller spaces from the point of view of {M}ostow rigidity.
\newblock In {\em Topology, ergodic theory, real algebraic geometry}, volume 202 of {\em Amer. Math. Soc. Transl. Ser. 2}, pages 131--149. Amer. Math. Soc., Providence, RI, 2001.

\bibitem{Kaimanovich2010}
Vadim~A. Kaimanovich.
\newblock Hopf decomposition and horospheric limit sets.
\newblock {\em Ann. Acad. Sci. Fenn. Math.}, 35(2):335--350, 2010.

\bibitem{Kechris1995}
Alexander~S. Kechris.
\newblock {\em Classical descriptive set theory}, volume 156 of {\em Graduate Texts in Mathematics}.
\newblock Springer-Verlag, New York, 1995.

\bibitem{KentLeininger2008}
Richard~P. Kent, IV and Christopher~J. Leininger.
\newblock Shadows of mapping class groups: capturing convex cocompactness.
\newblock {\em Geom. Funct. Anal.}, 18(4):1270--1325, 2008.

\bibitem{Kra1981}
Irwin Kra.
\newblock {T}he {C}arath{\'e}odory metric on abelian {T}eichm{\"u}ller disks.
\newblock {\em J. Analyse Math.}, 40:129--143 (1982), 1981.

\bibitem{LindenstraussMirzakhani2008}
Elon Lindenstrauss and Maryam Mirzakhani.
\newblock {E}rgodic theory of the space of measured laminations.
\newblock {\em Int. Math. Res. Not. IMRN}, (4):Art. ID rnm126, 49, 2008.

\bibitem{LiuSu2014}
Lixin Liu and Weixu Su.
\newblock The horofunction compactification of the {T}eichm{\"u}ller metric.
\newblock In {\em Handbook of {T}eichm{\"u}ller theory. {V}ol. {IV}}, volume~19 of {\em IRMA Lect. Math. Theor. Phys.}, pages 355--374. Eur. Math. Soc., Z{\"u}rich, 2014.

\bibitem{LusinPriwaloff1925}
N.~Lusin and J.~Priwaloff.
\newblock {S}ur l'unicit{\'e} et la multiplicit{\'e} des fonctions analytiques.
\newblock {\em {A}nn. {S}ci. {\'E}cole {N}orm. {S}up. (3)}, 42:143--191, 1925.

\bibitem{Maher2011}
Joseph Maher.
\newblock {R}andom walks on the mapping class group.
\newblock {\em Duke Math. J.}, 156(3):429--468, 2011.

\bibitem{MalesteinSouto2013}
Justin Malestein and Juan Souto.
\newblock {O}n genericity of pseudo-{A}nosovs in the {T}orelli group.
\newblock {\em Int. Math. Res. Not. IMRN}, (6):1434--1449, 2013.

\bibitem{Masur1980}
Howard Masur.
\newblock {U}niquely ergodic quadratic differentials.
\newblock {\em Comment. Math. Helv.}, 55(2):255--266, 1980.

\bibitem{Masur1982}
Howard Masur.
\newblock {I}nterval exchange transformations and measured foliations.
\newblock {\em Ann. of Math. (2)}, 115(1):169--200, 1982.

\bibitem{Masur1985}
Howard Masur.
\newblock {E}rgodic actions of the mapping class group.
\newblock {\em Proc. Amer. Math. Soc.}, 94(3):455--459, 1985.

\bibitem{Masur1986}
Howard Masur.
\newblock Measured foliations and handlebodies.
\newblock {\em Ergodic Theory Dynam. Systems}, 6(1):99--116, 1986.

\bibitem{Masur1992}
Howard Masur.
\newblock Hausdorff dimension of the set of nonergodic foliations of a quadratic differential.
\newblock {\em Duke Math. J.}, 66(3):387--442, 1992.

\bibitem{MasurSmillie1991}
Howard Masur and John Smillie.
\newblock {H}ausdorff dimension of sets of nonergodic measured foliations.
\newblock {\em Ann. of Math. (2)}, 134(3):455--543, 1991.

\bibitem{McCarthyPapadopoulos1989}
John McCarthy and Athanase Papadopoulos.
\newblock Dynamics on {T}hurston's sphere of projective measured foliations.
\newblock {\em Comment. Math. Helv.}, 64(1):133--166, 1989.

\bibitem{McCarthyPapadopoulos1996}
John McCarthy and Athanase Papadopoulos.
\newblock {F}undamental domains in {T}eichm{\"u}ller space.
\newblock {\em Ann. Acad. Sci. Fenn. Math.}, 21(1):151--166, 1996.

\bibitem{Mityagin2020}
B.~S. Mityagin.
\newblock The zero set of a real analytic function.
\newblock {\em Mat. Zametki}, 107(3):473--475, 2020.

\bibitem{Miyachi2008}
Hideki Miyachi.
\newblock Teichm{\"u}ller rays and the {G}ardiner-{M}asur boundary of {T}eichm{\"u}ller space.
\newblock {\em Geom. Dedicata}, 137:113--141, 2008.

\bibitem{Miyachi2013b}
Hideki Miyachi.
\newblock {A} differential formula for extremal length.
\newblock 590:137--152, 2013.

\bibitem{Miyachi2013}
Hideki Miyachi.
\newblock Teichm{\"u}ller rays and the {G}ardiner-{M}asur boundary of {T}eichm{\"u}ller space {II}.
\newblock {\em Geom. Dedicata}, 162:283--304, 2013.

\bibitem{Miyachi2017}
Hideki Miyachi.
\newblock {E}xtremal length functions are log-plurisubharmonic.
\newblock In {\em In the Tradition of Ahlfors--Bers, VII}, volume 696 of {\em Contemp. Math.}, pages 225--250. Amer. Math. Soc., Providence, RI, 2017.

\bibitem{Miyachi2019}
Hideki Miyachi.
\newblock {P}luripotential theory on {T}eichm{\"u}ller space {I}: {P}luricomplex {G}reen function.
\newblock {\em Conform. Geom. Dyn.}, 23:221--250, 2019.

\bibitem{Miyachi2023}
Hideki Miyachi.
\newblock {P}luripotential theory on {T}eichm{\"u}ller space {II}---{P}oisson integral formula.
\newblock {\em Adv. Math.}, 432:Paper No. 109265, 64, 2023.

\bibitem{Miyachi2024Bounded}
Hideki Miyachi.
\newblock {B}ounded {P}luriharmonic {F}unctions and {H}olomorphic {F}unctions on {T}eichm{\"u}ller {S}pace.
\newblock {\em Int. Math. Res. Not. IMRN}, (22):13855--13869, 2024.

\bibitem{miyachi2025bddpluriharmonicfunctionII-ToAppear}
Hideki Miyachi.
\newblock {B}ounded {P}luriharmonic {F}unctions and {H}olomorphic {F}unctions on {T}eichm{\"u}ller {S}pace {II} -- {P}oisson integral formula --.
\newblock Contemporary Mathematics, 2025.
\newblock To appear.

\bibitem{miyachi2025functiontheorydynamicsergodic}
Hideki Miyachi.
\newblock Function theory, dynamics and ergodic theory via thurston theory.
\newblock Preprint, ArXiv:2507.20912, 2025.

\bibitem{miyachi2026limitsetsmappingclass}
Hideki Miyachi and Ken'ichi Ohshika.
\newblock Limit sets of mapping class groups.
\newblock Preprint arXiv:2608.13354, 2026.

\bibitem{MiyachiOhshika2017}
Hideki Miyachi and Ken\textquotesingle~ichi Ohshika.
\newblock {U}ne formule diff{\'e}rentielle de la longueur extr{\'e}male et ses applications.
\newblock {\em Ann. Math. Blaise Pascal}, 24(1):115--133, 2017.

\bibitem{MoninTelpukhovskiy2019}
Leonid Monin and Vanya Telpukhovskiy.
\newblock {O}n normalizations of {T}hurston measure on the space of measured laminations.
\newblock {\em Topology Appl.}, 267:106878, 12, 2019.

\bibitem{Mumford1967}
David Mumford.
\newblock {A}belian quotients of the {T}eichm{\"u}ller modular group.
\newblock {\em J. Analyse Math.}, 18:227--244, 1967.

\bibitem{Nicholls1989}
Peter~J. Nicholls.
\newblock {\em {T}he ergodic theory of discrete groups}, volume 143 of {\em London Mathematical Society Lecture Note Series}.
\newblock Cambridge University Press, Cambridge, 1989.

\bibitem{PennerHarer1992}
R.~C. Penner and J.~L. Harer.
\newblock {\em Combinatorics of train tracks}, volume 125 of {\em Annals of Mathematics Studies}.
\newblock Princeton University Press, Princeton, NJ, 1992.

\bibitem{Pommerenke1976}
Ch. Pommerenke.
\newblock {O}n the {G}reen\textquotesingle s function of {F}uchsian groups.
\newblock {\em Ann. Acad. Sci. Fenn. Ser. A I Math.}, 2:409--427, 1976.

\bibitem{Rauch1965}
H.~E. Rauch.
\newblock {A} transcendental view of the space of algebraic {R}iemann surfaces.
\newblock {\em Bull. Amer. Math. Soc.}, 71:1--39, 1965.

\bibitem{Rees1981}
Mary Rees.
\newblock {A}n alternative approach to the ergodic theory of measured foliations on surfaces.
\newblock {\em Ergodic Theory Dynamical Systems}, 1(4):461--488 (1982), 1981.

\bibitem{Rivin2008}
Igor Rivin.
\newblock {W}alks on groups, counting reducible matrices, polynomials, and surface and free group automorphisms.
\newblock {\em Duke Math. J.}, 142(2):353--379, 2008.

\bibitem{Rivin2009}
Igor Rivin.
\newblock {W}alks on graphs and lattices---effective bounds and applications.
\newblock {\em Forum Math.}, 21(4):673--685, 2009.

\bibitem{Rohlin1952}
V.~A. Rohlin.
\newblock {O}n the fundamental ideas of measure theory.
\newblock {\em Amer. Math. Soc. Translation}, 1952(71):55, 1952.

\bibitem{Royden1971b}
H.~L. Royden.
\newblock {R}emarks on the {K}obayashi metric.
\newblock In {\em Several complex variables, II (Proc. Internat. Conf., Univ. Maryland, College Park, Md., 1970)}, volume Vol. 185 of {\em Lecture Notes in Math.}, pages 125--137. Springer, Berlin-New York, 1971.

\bibitem{Satake1980}
Ichir{\^o} Satake.
\newblock {\em {A}lgebraic structures of symmetric domains}, volume~4 of {\em Kan{\^o} Memorial Lectures}.
\newblock Iwanami Shoten, Tokyo; Princeton University Press, Princeton, NJ, 1980.

\bibitem{Shiga1984}
Hiroshige Shiga.
\newblock {O}n analytic and geometric properties of {T}eichm{\"u}ller spaces.
\newblock {\em J. Math. Kyoto Univ.}, 24(3):441--452, 1984.

\bibitem{Siegel1943}
Carl~Ludwig Siegel.
\newblock {S}ymplectic geometry.
\newblock {\em Amer. J. Math.}, 65:1--86, 1943.

\bibitem{SuTan2023}
Weixu Su and Dong Tan.
\newblock {H}orospheres in {T}eichm{\"u}ller space and mapping class group.
\newblock {\em Ann. Inst. Fourier (Grenoble)}, 73(4):1677--1707, 2023.

\bibitem{Sullivan1981}
Dennis Sullivan.
\newblock {O}n the ergodic theory at infinity of an arbitrary discrete group of hyperbolic motions.
\newblock In {\em Riemann surfaces and related topics: Proceedings of the 1978 Stony Brook Conference (State Univ. New York, Stony Brook, N.Y., 1978)}, volume~97 of {\em Ann. of Math. Stud.}, pages 465--496. Princeton Univ. Press, Princeton, N.J., 1981.

\bibitem{Thurston-LectureNote}
William {Thurston}.
\newblock {\em The {G}eometry and {T}opology of {T}hree-{M}anifolds}.
\newblock 1980.
\newblock {L}ecture {N}ote at {P}rinceton {U}niversity, {A}vailable at {\tt http://library.msri.org/nonmsri/gt3m/}.

\bibitem{Tsuji1944}
Masatsugu Tsuji.
\newblock {T}heory of meromorphic functions in a neighbourhood of a closed set of capacity zero.
\newblock {\em Jpn. J. Math.}, 19:139--154, 1944.

\bibitem{Tukia1997}
Pekka Tukia.
\newblock {C}onservative action and the horospheric limit set.
\newblock {\em Ann. Acad. Sci. Fenn. Math.}, 22(2):387--394, 1997.

\bibitem{Walsh2019}
Cormac Walsh.
\newblock The asymptotic geometry of the {T}eichm\"uller metric.
\newblock {\em Geom. Dedicata}, 200:115--152, 2019.

\end{thebibliography}

\end{document}
